% =====================================================================
%  Глава «Вычислительная линейная алгебра».
%  Самодостаточные сюжеты: SVD, Eigenfaces/Eigencats, PageRank,
%  оценки Липшицевости нейросети, числа с плавающей точкой.
%  Источники: nla360, optim, ds-intro.
%
%  Назначение в книге: новая глава Части II «Математические основы».
%  Если будет принято решение о вставке между Ньютоном (текущая гл. 2)
%  и Частью III, главе автоматически присвоится номер 3 (а текущая
%  глава 3 станет 4-й и т. д.). Все ссылки внутри файла — через \ref,
%  поэтому перенумерация происходит автоматически.
%
%  Подключение из main.tex (предлагаемая позиция):
%      \part{Математические основы}
%      \chapter{Численные методы оптимизации}
\label{ch:opt}

\section{Метод касательных Ньютона и его приложения}
\label{sec:newton}

\subsection{От Ньютона до Канторовича}

Задачи отыскания корня уравнения $f(x)=0$ и поиска минимума функции
$g(x)$ относятся к разным «жанрам» школьной математики, однако
численно они решаются почти одним и тем же приёмом — итерационной
схемой, придуманной Исааком Ньютоном.

\begin{historybox}
В работе \emph{De analysi per aequationes numero terminorum infinitas}
(1669) и затем в \emph{Method of Fluxions} (рукопись 1671~г., издание
1736~г.) Ньютон предложил способ приближённого решения уравнений,
который позже Джозеф Рафсон в 1690~г. переписал в современной
рекуррентной форме. Поэтому в англоязычной литературе метод часто
называют \emph{Newton–Raphson}.

Существенный вклад в строгое исследование метода внёс советский
математик \textbf{Леонид Витальевич Канторович} (1912–1986),
лауреат Нобелевской премии по экономике 1975~г. (за теорию
оптимального распределения ресурсов). В 1948~г. он опубликовал
теорему, носящую теперь его имя: она впервые дала универсальные
\emph{достаточные условия сходимости} метода Ньютона, причём не
только для уравнений с одной переменной, но и для уравнений в
бесконечномерных банаховых пространствах.
\end{historybox}

\paragraph{Геометрическая идея.} Пусть требуется найти корень
уравнения~$f(x)=0$, и нам известна точка $x_k$, близкая к корню.
Заменим график $y=f(x)$ его \emph{касательной} в точке $(x_k, f(x_k))$.
Касательная — это прямая, и её пересечение с осью $x$ найти просто.
Точку этого пересечения и объявим следующим приближением~$x_{k+1}$.
Затем повторим. Так получается метод касательных~(рис.~\ref{fig:newton-geom}).

\begin{figure}[ht]
\centering
\includegraphics[width=0.85\linewidth]{newton_geometry.pdf}
\caption{Три итерации метода Ньютона для уравнения
$x^{3}-2x-5=0$. Каждая касательная (пунктир) пересекает ось абсцисс
в точке, которая берётся в качестве следующего приближения
(квадрат). После трёх шагов корень найден с точностью лучше $10^{-6}$.}
\label{fig:newton-geom}
\end{figure}

Уравнение касательной к графику $y=f(x)$ в точке $x_k$ имеет вид
$y=f(x_k)+f'(x_k)\,(x-x_k)$.
Полагая $y=0$ и выражая $x$, получаем
\begin{equation}\label{eq:newton-root}
\boxed{\;x_{k+1}=x_k-\dfrac{f(x_k)}{f'(x_k)}\;}\qquad k=0,1,2,\dots
\end{equation}

Это и есть \emph{итерационная формула Ньютона}. Точка $x_0$ выбирается
исследователем; от её выбора может зависеть, сойдётся ли метод
вообще.

\subsection{Метод Ньютона для задач оптимизации}

Пусть теперь нужно найти минимум функции $g$. В любой точке минимума
$x^{\star}$ выполнено $g'(x^{\star})=0$. Значит, минимизация
дифференцируемой функции — это решение уравнения
\[
g'(x)=0.
\]
Применим к нему формулу~\eqref{eq:newton-root} с $f=g'$ (и тогда $f'=g''$):
\begin{equation}\label{eq:newton-opt}
\boxed{\;x_{k+1}=x_k-\dfrac{g'(x_k)}{g''(x_k)}\;}
\end{equation}

Эту же формулу можно получить иначе — и это ценная для всей теории
оптимизации точка зрения. Разложим $g$ в окрестности~$x_k$ в
\textbf{ряд Тейлора}:
\begin{equation}\label{eq:taylor}
g(x)\;=\;g(x_k)\;+\;g'(x_k)(x-x_k)\;+\;\tfrac12 g''(x_k)(x-x_k)^2
   \;+\;o\!\bigl((x-x_k)^2\bigr).
\end{equation}

Откуда берётся такое разложение? Идея проста: рядом с точкой $x_k$
функция $g$ ведёт себя «почти как многочлен». Линейное приближение
$g(x_k)+g'(x_k)(x-x_k)$ — это уже знакомая нам касательная к
графику. Если мы хотим добиться большей точности, надо учесть и
\emph{кривизну} графика — а она как раз кодируется второй производной
$g''$. Так возникает квадратичный многочлен в~\eqref{eq:taylor}.

Отбросим бесконечно малое слагаемое и обозначим оставшийся квадратичный
многочлен через
\[
q_k(x)\;=\;g(x_k)+g'(x_k)(x-x_k)+\tfrac12 g''(x_k)(x-x_k)^2.
\]
Минимизировать $q_k$ умеет любой школьник: это квадратное уравнение,
решаемое в одну строку. Условие $q_k'(x)=0$ даёт $g'(x_k)+g''(x_k)(x-x_k)=0$,
откуда $x=x_k-g'(x_k)/g''(x_k)$, что в точности совпадает
с~формулой~\eqref{eq:newton-opt}.

\begin{notebox}[title={«Разделяй и властвуй»}]
Метод Ньютона для оптимизации — классический пример принципа
\emph{разделяй и властвуй}: вместо того чтобы решать сложную задачу
минимизации произвольной функции $g$ целиком, мы заменяем её
на каждом шаге простой подзадачей — минимизацией параболы $q_k$.
Парабола подбирается так, чтобы локально (в окрестности $x_k$)
имитировать поведение $g$, и её минимум вычисляется явно.
\end{notebox}

\begin{figure}[ht]
\centering
\includegraphics[width=0.78\linewidth]{taylor_parabola.pdf}
\caption{Метод Ньютона как минимизация локальной квадратичной модели
$q_k$ (пунктир): следующая точка $x_{k+1}$ — это вершина параболы,
касающейся графика $g$ в точке $x_k$.}
\label{fig:taylor}
\end{figure}

\subsection{Пример 1. Извлечение квадратного корня}
\label{ssec:heron}

Пусть $a>0$. Вычислим $\sqrt{a}$, применяя метод Ньютона
к уравнению
\[
f(x)=x^{2}-a=0.
\]
Тогда $f'(x)=2x$, и формула~\eqref{eq:newton-root} принимает вид
\begin{equation}\label{eq:heron}
\boxed{\;x_{k+1}=\dfrac12\!\left(x_k+\dfrac{a}{x_k}\right)\;}
\end{equation}

Это — знаменитый \textbf{метод Герона}, или вавилонский метод
квадратного корня.

\begin{historybox}
Формулу~\eqref{eq:heron} в явном виде описал \emph{Герон Александрийский}
(I век н.~э.) в трактате «Метрика». Однако эту итерацию умели
применять задолго до него: на вавилонской глиняной табличке
\textsf{YBC~7289} (примерно XVIII–XVI вв. до н.~э.) сохранилось
вычисление $\sqrt{2}$ с точностью до шестого десятичного знака —
именно с помощью среднего арифметического и среднего гармонического,
что эквивалентно~\eqref{eq:heron}. Получается, что одна из
итерационных схем «пережила» три с половиной тысячи лет и
по сей день используется в библиотеках вычисления корня.
\end{historybox}

\subsubsection*{Сходимость через сжимающее отображение}

Запишем~\eqref{eq:heron} как $x_{k+1}=T(x_k)$, где
$T(x)=\tfrac12(x+a/x)$. Точка $\sqrt{a}$ — \emph{неподвижная} для $T$:
$T(\sqrt{a})=\sqrt{a}$.

\begin{theorem}{Сходимость метода Герона}{heron}
Для любого $x_0>0$ последовательность~\eqref{eq:heron} монотонно
убывает (начиная с $k=1$) и сходится к $\sqrt{a}$. Более того, для
ошибки $e_k=x_k-\sqrt{a}$ выполнено
\begin{equation}\label{eq:heron-err}
e_{k+1}\;=\;\dfrac{e_k^2}{2x_k}\,, \qquad
\text{и при }k\geq 1\colon\quad
0\le e_{k+1}\le \dfrac{e_k^{2}}{2\sqrt{a}}.
\end{equation}
\end{theorem}

\begin{proof}
\textbf{Шаг 1 (нижняя оценка).} По неравенству о средних
(AM--GM) для $x_0>0$:
$x_1=\tfrac12(x_0+a/x_0)\geq \sqrt{x_0\cdot a/x_0}=\sqrt{a}$.
Значит, начиная с $x_1$, все члены лежат в~$[\sqrt{a},\infty)$.

\textbf{Шаг 2 (рекуррентность для ошибки).}
\[
e_{k+1}=x_{k+1}-\sqrt{a}=\dfrac{x_k+a/x_k}{2}-\sqrt{a}
=\dfrac{x_k^{2}-2\sqrt{a}\,x_k+a}{2x_k}=\dfrac{(x_k-\sqrt{a})^2}{2x_k}=\dfrac{e_k^{2}}{2x_k}.
\]
При $x_k\geq\sqrt{a}$ имеем $\tfrac{1}{2x_k}\leq\tfrac{1}{2\sqrt{a}}$,
откуда вторая часть~\eqref{eq:heron-err}.

\textbf{Шаг 3 (сжатие).} Производная $T'(x)=\tfrac12(1-a/x^{2})$.
На луче $x\geq\sqrt{a}$ выполнено $T'(x)\in[0,\tfrac12)$. Значит,
$T$~— \emph{сжимающее отображение} (с константой~$\tfrac12$) на
$[\sqrt{a},\infty)$. По теореме Банаха о неподвижной точке
последовательность $x_k=T^{k}(x_1)$ сходится к единственной
неподвижной точке~$\sqrt{a}$.
\end{proof}

Из~\eqref{eq:heron-err} видно, что после первой итерации ошибка
\emph{квадрируется}: каждые два десятичных знака удваиваются за
один шаг. Поэтому при разумном $x_0$ хватает 4--5 итераций, чтобы
получить машинную точность.

\begin{example}{Расстояние до горизонта}{horizon}
Стоящий человек ростом $h=1{,}70$~м смотрит на море. До какой точки
далеко он видит? Из теоремы Пифагора (рис.~\ref{fig:horizon}) для прямой,
касательной к поверхности Земли:
\[
d^{2}+R^{2}=(R+h)^{2},\qquad
d=\sqrt{2Rh+h^{2}}\;\approx\;\sqrt{2Rh}.
\]
При $R=6371$~км, $h=1{,}7\cdot 10^{-3}$~км:
\[
2Rh\;\approx\;2\cdot 6371\cdot 0{,}0017\;=\;21{,}66\;\text{(км$^2$)}.
\]
Применим к этому числу метод Герона со стартом $x_0=5$:
\[
\begin{array}{l|l}
k & x_k \\\hline
0 & 5{,}000\,000 \\
1 & \tfrac{1}{2}(5+21{,}66/5)=4{,}666\,000 \\
2 & 4{,}654\,463 \\
3 & 4{,}654\,460
\end{array}
\]
Получаем $d\approx 4{,}65$~км — расстояние, которое запоминается
любому, кто хоть раз стоял у моря.
\end{example}

\begin{figure}[ht]
\centering
\begin{tikzpicture}[scale=0.9, line cap=round]
  % Земля
  \draw[thick, accent] (0,0) circle (2.6);
  \node[accent!80!black] at (0,-0.4) {Земля};
  % центр
  \fill (0,0) circle (1.2pt) node[below right=1pt] {$O$};
  % точка наблюдателя
  \coordinate (A) at (0,2.6);
  \coordinate (P) at (0,3.05);          % человек на высоте h
  \draw[very thick, accent2] (A) -- (P);
  \fill (P) circle (1.2pt) node[left] {$P$};
  \fill (A) circle (1pt) node[left=1pt] {$A$};
  % касательная: из P касается окружности в точке T
  % используем геометрию: |OP|=R+h, расстояние от O до T равно R, угол OTC=90.
  \pgfmathsetmacro{\Rh}{3.05}
  \pgfmathsetmacro{\R}{2.6}
  \pgfmathsetmacro{\dlen}{sqrt(\Rh*\Rh-\R*\R)}    % длина касательной
  \pgfmathsetmacro{\sina}{\dlen/\Rh}
  \pgfmathsetmacro{\cosa}{\R/\Rh}
  \coordinate (T) at ({\R*\sina}, {\R*\cosa});
  \draw[thick, accent2] (P) -- (T);
  \draw[dashed] (0,0) -- (T);
  \fill (T) circle (1pt) node[right=2pt] {$T$};
  % прямой угол в T
  \draw ($(T)+({-0.18*sin(40)},{-0.18*cos(40)})$) -- 
        ++({0.18*cos(40)},{-0.18*sin(40)}) --
        ++({0.18*sin(40)},{0.18*cos(40)});
  % метки
  \node[accent2] at (0.45,2.8) {$h$};
  \node at (0.85,1.45) {$d$};
  \node[rotate=15] at (-0.7,1.4) {$R$};
\end{tikzpicture}
\caption{Расстояние до горизонта $d=\sqrt{(R+h)^2-R^2}\approx\sqrt{2Rh}$.}
\label{fig:horizon}
\end{figure}

\subsection{Пример 2. Вычисление обратного числа}

Что, если на нашем процессоре нет аппаратного деления, и нужно
вычислить $1/a$ при $a>0$ (только сложениями и умножениями)?
Применим метод Ньютона к уравнению
\[
f(x)=\dfrac{1}{x}-a=0.
\]
Тогда $f'(x)=-1/x^{2}$, и подстановка в~\eqref{eq:newton-root} даёт
\[
x_{k+1}=x_k-\dfrac{1/x_k-a}{-1/x_k^{2}}
       =x_k+x_k^{2}\bigl(1/x_k-a\bigr)
       =x_k\bigl(2-a x_k\bigr).
\]
То есть
\begin{equation}\label{eq:ns}
\boxed{\;x_{k+1}=x_k\,(2-a x_k)\;}
\end{equation}

Эта схема замечательна тем, что в ней нет ни одного деления — только
сложение и умножение. Поэтому именно она исторически использовалась
в первых ЭВМ для реализации операции деления на аппаратном уровне.

\subsubsection*{Анализ сходимости}

Введём «относительную невязку» $u_k=1-a x_k$. Тогда
\[
u_{k+1}=1-a x_{k+1}=1-a x_k(2-ax_k)=1-2 a x_k+(a x_k)^{2}=(1-a x_k)^{2}=u_k^{2}.
\]
Таким образом
\begin{equation}\label{eq:ns-error}
u_k=u_0^{\,2^{k}}\quad\Longrightarrow\quad u_k\to 0\;\Leftrightarrow\;|u_0|<1.
\end{equation}

Условие $|1-a x_0|<1$ эквивалентно $0<x_0<2/a$. \emph{Областью
сходимости} метода является интервал $\bigl(0,\,2/a\bigr)$, и
он \textbf{конечен}: вне этого интервала ($x_0>2/a$ или $x_0<0$)
итерации расходятся. Это принципиальное отличие от метода Герона,
который сходится для любого $x_0>0$.

\begin{notebox}[title={Метод Ньютона--Шульца и оптимизатор Muon}]
Если в~\eqref{eq:ns} заменить число $x_k$ на квадратную матрицу $X_k$,
а число $a$ — на матрицу $A$, получим знаменитую матричную итерацию
\textbf{Ньютона--Шульца}: $X_{k+1}=X_k(2I-A X_k)$. Она вычисляет
$A^{-1}$ без обращения матриц на каждом шаге — лишь умножениями.

Аналогичные полиномиальные итерации применяются для приближённой
\emph{ортогонализации} матриц. Именно так устроен оптимизатор
\textbf{Muon}, появившийся в 2024–2025~гг.\ и используемый при
обучении современных больших языковых моделей: в нём матрица
градиента на каждом шаге быстро ортогонализуется несколькими
итерациями Ньютона--Шульца — это намного дешевле, чем точное
сингулярное разложение.
\end{notebox}

\subsection{Реализация на Python и численный эксперимент}

Запишем универсальную функцию, реализующую метод Ньютона
по формуле~\eqref{eq:newton-root}, и применим её к двум разобранным
примерам. Комментарии в коде даны по-английски (как принято в
большинстве библиотек): \texttt{f}, \texttt{df} — сама функция и
её производная; \texttt{x0} — начальное приближение; \texttt{n\_iter} —
максимальное число итераций; \texttt{tol} — точность остановки;
функция возвращает массив \texttt{xs} с историей итераций.

\begin{lstlisting}[label={lst:newton}]
import numpy as np
import matplotlib.pyplot as plt


def newton(f, df, x0, n_iter=20, tol=1e-15):
    """Newton's tangent method.

    Parameters
    ----------
    f, df  : callables f(x) and f'(x).
    x0     : initial guess.
    n_iter : maximum number of iterations.
    tol    : stopping tolerance on |f(x)|.

    Returns
    -------
    Array of iterates, including x0.
    """
    xs = [x0]
    x = x0
    for _ in range(n_iter):
        fx = f(x)
        if abs(fx) < tol:
            break
        x = x - fx / df(x)
        xs.append(x)
    return np.array(xs)


# ---- Example 1. Square root of a ---------------------------------
a = 2.0
xs_sqrt = newton(lambda x: x**2 - a,
                 lambda x: 2 * x,
                 x0=50.0)          # very poor initial guess
err_sqrt = np.abs(xs_sqrt - np.sqrt(a))

# ---- Example 2. Reciprocal 1/a -----------------------------------
a = 3.0
xs_inv  = newton(lambda x: 1.0 / x - a,
                 lambda x: -1.0 / x**2,
                 x0=0.20)          # inside the convergence radius (0, 2/a)
err_inv = np.abs(xs_inv - 1.0 / a)

# ---- Convergence plot in log scale -------------------------------
plt.semilogy(err_sqrt, 'o-', label=r'sqrt: $|x_k-\sqrt{2}|$')
plt.semilogy(err_inv,  's-', label=r'1/a:  $|x_k-1/3|$')
plt.xlabel('iteration $k$'); plt.ylabel('error')
plt.legend(); plt.grid(True); plt.show()
\end{lstlisting}

\noindent\small\textit{Листинг~2.1. Универсальный метод касательных
Ньютона и его применение.}\normalsize

\paragraph{Что показывает эксперимент.}

На рис.~\ref{fig:sqrt-conv} приведены кривые сходимости метода Герона
для $\sqrt{2}$ при разных стартах $x_0\in\{0{,}1;\,1;\,5;\,50\}$.
Все четыре кривые сходятся к нулю, причём поведение типичное для
квадратичной сходимости: после нескольких «согласующих» шагов
число верных знаков удваивается за итерацию. Это иллюстрирует
\textbf{глобальную сходимость} метода Герона: годится \emph{любое}
начальное приближение из $(0,\infty)$.

\begin{figure}[ht]
\centering
\includegraphics[width=0.85\linewidth]{sqrt_convergence.pdf}
\caption{Глобальная сходимость метода Герона для $\sqrt{2}$ из четырёх
очень разных стартовых точек. Шкала ошибки логарифмическая.}
\label{fig:sqrt-conv}
\end{figure}

Совсем другое поведение демонстрирует метод Ньютона--Шульца для
$1/a$ при $a=3$ (рис.~\ref{fig:rec-conv}): область сходимости — это
интервал $(0,\,2/3)\approx(0,\,0{,}667)$.
Стартуя \emph{внутри} этого интервала, метод сходится квадратично
(нижние три кривые); вблизи правой границы (например, $x_0=0{,}66$) — сходится,
но «впритык»; а уже при $x_0=0{,}70$ ошибка взрывообразно растёт:
$|u_k|=|u_0|^{2^{k}}\to\infty$ в полном согласии с
\eqref{eq:ns-error}. Это \textbf{локальная сходимость}.

\begin{figure}[ht]
\centering
\includegraphics[width=0.85\linewidth]{reciprocal_convergence.pdf}
\caption{Локальная сходимость метода Ньютона--Шульца для $1/3$.
Область сходимости — интервал $(0,\,2/3)$; вне его метод расходится.}
\label{fig:rec-conv}
\end{figure}

\subsection{Достаточные условия сверхлинейной сходимости\texorpdfstring{$^\star$}{*}}
\label{ssec:thm-superlinear}

Сравнение двух примеров наводит на вопрос: \emph{когда именно}
метод Ньютона сходится так быстро, как мы наблюдаем? Ответ
даёт следующая теорема, восходящая к Канторовичу.

\begin{definition}{Сверхлинейная и квадратичная сходимости}{rates}
Говорят, что последовательность $\{x_k\}$ сходится к $x^{\star}$
со \textbf{сверхлинейной скоростью}, если
\[
\lim_{k\to\infty}\dfrac{|x_{k+1}-x^{\star}|}{|x_k-x^{\star}|}=0.
\]
Если же выполнено более сильное условие
$|x_{k+1}-x^{\star}|\leq M |x_k-x^{\star}|^{2}$ для некоторой
константы $M>0$, то сходимость называют \textbf{квадратичной}.
\end{definition}

\begin{definition}{Сильно выпуклая функция}{strconvex}
Функция $g\colon\R\to\R$ называется \textbf{сильно выпуклой} с
константой $\mu>0$, если $g\in C^{2}$ и $g''(x)\geq\mu$ для всех $x$.
Сильная выпуклость гарантирует существование и единственность
точки минимума и, что важно для нас, отделённость второй производной
от нуля.
\end{definition}

\begin{theorem}{Сверхлинейная сходимость метода Ньютона}{newton-rate}
Пусть выполнено хотя бы одно из условий:
\begin{enumerate}\itemsep0pt
  \item[(а)] (\emph{уравнение}) $f\in C^{2}$ в окрестности точки $x^{\star}$,
        $f(x^{\star})=0$ и $f'(x^{\star})\neq 0$;
  \item[(б)] (\emph{сильно выпуклая оптимизация}) $g\in C^{3}$,
        $g$~— сильно выпуклая на интервале, содержащем точку
        минимума $x^{\star}$, и применяется итерация~\eqref{eq:newton-opt}.
\end{enumerate}
Тогда существуют $\delta>0$ и $M>0$ такие, что для всякого
$x_0\in(x^{\star}-\delta,\,x^{\star}+\delta)$ итерации Ньютона
определены и
\begin{equation}\label{eq:rate}
|x_{k+1}-x^{\star}|\;\leq\;M\,|x_k-x^{\star}|^{2}.
\end{equation}
В частности, сходимость квадратичная (а значит, сверхлинейная).
\end{theorem}

\begin{proof}
Достаточно доказать случай (а): задача (б) сводится к нему,
если положить $f=g'$, тогда условие $f'(x^{\star})\neq 0$
эквивалентно $g''(x^{\star})\geq\mu>0$.

Разложим $f$ по формуле Тейлора с остаточным членом в форме
Лагранжа в окрестности~$x_k$ и подставим $x=x^{\star}$:
\[
0\;=\;f(x^{\star})\;=\;f(x_k)+f'(x_k)(x^{\star}-x_k)+\tfrac{f''(\xi_k)}{2}(x^{\star}-x_k)^{2},
\]
где $\xi_k$~— некоторая точка между $x_k$ и $x^{\star}$. Разделим
обе части на $f'(x_k)$ (по непрерывности эта величина не равна
нулю в малой окрестности $x^{\star}$):
\[
0\;=\;\dfrac{f(x_k)}{f'(x_k)}+(x^{\star}-x_k)+\dfrac{f''(\xi_k)}{2 f'(x_k)}(x^{\star}-x_k)^{2}.
\]
Из определения $x_{k+1}=x_k-f(x_k)/f'(x_k)$ получаем
$\frac{f(x_k)}{f'(x_k)}=x_k-x_{k+1}$. Подставляя:
\[
0=(x_k-x_{k+1})+(x^{\star}-x_k)+\dfrac{f''(\xi_k)}{2 f'(x_k)}(x^{\star}-x_k)^{2},
\]
откуда после очевидных переобозначений
\begin{equation}\label{eq:newton-recur}
x_{k+1}-x^{\star}\;=\;\dfrac{f''(\xi_k)}{2 f'(x_k)}\,(x_k-x^{\star})^{2}.
\end{equation}
Это и есть искомая \emph{рекуррентность для расстояния}.

Теперь подберём константы. По непрерывности $f'$ и $f''$ существует
$\delta_1>0$ и числа $L$, $m>0$ такие, что
$|f''(\xi)|\leq L$ и $|f'(x)|\geq m$ для всех $\xi,x\in
[x^{\star}-\delta_1,\,x^{\star}+\delta_1]$.
Положим $M=L/(2m)$ и $\delta=\min\!\bigl(\delta_1,\,\tfrac{1}{2M}\bigr)$.
Тогда из~\eqref{eq:newton-recur} при $|x_k-x^{\star}|<\delta$:
\[
|x_{k+1}-x^{\star}|\leq M|x_k-x^{\star}|^{2}\leq M\delta\cdot|x_k-x^{\star}|\leq\tfrac12|x_k-x^{\star}|<\delta.
\]
Значит, последовательность не покидает $\delta$-окрестность и
оценка~\eqref{eq:rate} выполнена для всех $k$.
\end{proof}

\paragraph{Где условия выполняются для наших примеров?}

\textbf{Метод Герона.} Здесь $f(x)=x^{2}-a$, $f'(x)=2x$,
$f''(x)=2$. В точке $x^{\star}=\sqrt{a}>0$ выполнено $f'(x^{\star})=2\sqrt{a}\neq 0$.
\emph{Условия теоремы выполняются на всём луче $(0,\infty)$,
не содержащем ноль}, и константа в~\eqref{eq:rate} принимает вид
$M=\frac{1}{2\sqrt{a}}$ — в полном согласии с
явным выражением~\eqref{eq:heron-err}, полученным элементарно.

\textbf{Метод Ньютона--Шульца.} Здесь $f(x)=1/x-a$, $f'(x)=-1/x^{2}$,
$f''(x)=2/x^{3}$. В точке $x^{\star}=1/a$ имеем $f'(x^{\star})=-a^{2}\neq 0$,
и теорема применима в малой окрестности этой точки. Однако, в
отличие от метода Герона, оценка снизу $|f'(x)|\geq m>0$ \emph{ломается}
при $x\to\infty$: производная стремится к нулю. Именно поэтому
радиус сходимости в этом примере конечен и равен $1/a$ слева и справа
от $x^{\star}$, что в сумме даёт интервал $(0,\,2/a)$, найденный нами
ранее аналитически.

\medskip

Таким образом, два простейших примера полностью укладываются в
доказанную теорему: метод Ньютона сходится сверхлинейно \emph{там,
где} производная регулярна и отделена от нуля. Глобальную же
сходимость, как у Герона, приходится доказывать дополнительно — например,
ссылаясь на структуру функции (выпуклость) и теорему Банаха о
сжимающем отображении, как мы и сделали в разделе~\ref{ssec:heron}.

\begin{tasksbox}
\begin{enumerate}\itemsep1pt
  \item Применяя метод Ньютона к уравнению $f(x)=x^{n}-a=0$, получите
        итерационную формулу для $\sqrt[n]{a}$. Покажите, что для $n=3$
        она имеет вид $x_{k+1}=\frac{2 x_k+a/x_k^{2}}{3}$.
  \item Докажите, что для метода Герона из любого $x_0>0$ погрешность
        $e_1$ всегда положительна, после чего последовательность
        монотонно убывает. (Указание: используйте AM--GM.)
  \item Для метода Ньютона--Шульца при $a=2$ найдите наибольшее $x_0$,
        при котором итерации ещё сходятся. Сравните численно скорость
        сходимости при $x_0=0{,}1$ и $x_0=0{,}9$.
  \item$^{\star}$ Сформулируйте и докажите аналог
        теоремы~\ref{thm:newton-rate} для метода Ньютона минимизации
        многомерной сильно выпуклой функции $g\colon\R^{n}\to\R$.
        Какую роль играет матрица Гессе $\nabla^{2}g$?
  \item$^{\star}$ Прочитайте о теореме Канторовича в учебнике
        \href{https://en.wikipedia.org/wiki/Kantorovich_theorem}{«Kantorovich theorem»}
        и сформулируйте её в одну фразу: какие именно три величины
        должны быть «в нужном соотношении» для гарантии сходимости?
\end{enumerate}
\end{tasksbox}

%      % =====================================================================
%  Глава «Вычислительная линейная алгебра».
%  Самодостаточные сюжеты: SVD, Eigenfaces/Eigencats, PageRank,
%  оценки Липшицевости нейросети, числа с плавающей точкой.
%  Источники: nla360, optim, ds-intro.
%
%  Назначение в книге: новая глава Части II «Математические основы».
%  Если будет принято решение о вставке между Ньютоном (текущая гл. 2)
%  и Частью III, главе автоматически присвоится номер 3 (а текущая
%  глава 3 станет 4-й и т. д.). Все ссылки внутри файла — через \ref,
%  поэтому перенумерация происходит автоматически.
%
%  Подключение из main.tex (предлагаемая позиция):
%      \part{Математические основы}
%      \chapter{Численные методы оптимизации}
\label{ch:opt}

\section{Метод касательных Ньютона и его приложения}
\label{sec:newton}

\subsection{От Ньютона до Канторовича}

Задачи отыскания корня уравнения $f(x)=0$ и поиска минимума функции
$g(x)$ относятся к разным «жанрам» школьной математики, однако
численно они решаются почти одним и тем же приёмом — итерационной
схемой, придуманной Исааком Ньютоном.

\begin{historybox}
В работе \emph{De analysi per aequationes numero terminorum infinitas}
(1669) и затем в \emph{Method of Fluxions} (рукопись 1671~г., издание
1736~г.) Ньютон предложил способ приближённого решения уравнений,
который позже Джозеф Рафсон в 1690~г. переписал в современной
рекуррентной форме. Поэтому в англоязычной литературе метод часто
называют \emph{Newton–Raphson}.

Существенный вклад в строгое исследование метода внёс советский
математик \textbf{Леонид Витальевич Канторович} (1912–1986),
лауреат Нобелевской премии по экономике 1975~г. (за теорию
оптимального распределения ресурсов). В 1948~г. он опубликовал
теорему, носящую теперь его имя: она впервые дала универсальные
\emph{достаточные условия сходимости} метода Ньютона, причём не
только для уравнений с одной переменной, но и для уравнений в
бесконечномерных банаховых пространствах.
\end{historybox}

\paragraph{Геометрическая идея.} Пусть требуется найти корень
уравнения~$f(x)=0$, и нам известна точка $x_k$, близкая к корню.
Заменим график $y=f(x)$ его \emph{касательной} в точке $(x_k, f(x_k))$.
Касательная — это прямая, и её пересечение с осью $x$ найти просто.
Точку этого пересечения и объявим следующим приближением~$x_{k+1}$.
Затем повторим. Так получается метод касательных~(рис.~\ref{fig:newton-geom}).

\begin{figure}[ht]
\centering
\includegraphics[width=0.85\linewidth]{newton_geometry.pdf}
\caption{Три итерации метода Ньютона для уравнения
$x^{3}-2x-5=0$. Каждая касательная (пунктир) пересекает ось абсцисс
в точке, которая берётся в качестве следующего приближения
(квадрат). После трёх шагов корень найден с точностью лучше $10^{-6}$.}
\label{fig:newton-geom}
\end{figure}

Уравнение касательной к графику $y=f(x)$ в точке $x_k$ имеет вид
$y=f(x_k)+f'(x_k)\,(x-x_k)$.
Полагая $y=0$ и выражая $x$, получаем
\begin{equation}\label{eq:newton-root}
\boxed{\;x_{k+1}=x_k-\dfrac{f(x_k)}{f'(x_k)}\;}\qquad k=0,1,2,\dots
\end{equation}

Это и есть \emph{итерационная формула Ньютона}. Точка $x_0$ выбирается
исследователем; от её выбора может зависеть, сойдётся ли метод
вообще.

\subsection{Метод Ньютона для задач оптимизации}

Пусть теперь нужно найти минимум функции $g$. В любой точке минимума
$x^{\star}$ выполнено $g'(x^{\star})=0$. Значит, минимизация
дифференцируемой функции — это решение уравнения
\[
g'(x)=0.
\]
Применим к нему формулу~\eqref{eq:newton-root} с $f=g'$ (и тогда $f'=g''$):
\begin{equation}\label{eq:newton-opt}
\boxed{\;x_{k+1}=x_k-\dfrac{g'(x_k)}{g''(x_k)}\;}
\end{equation}

Эту же формулу можно получить иначе — и это ценная для всей теории
оптимизации точка зрения. Разложим $g$ в окрестности~$x_k$ в
\textbf{ряд Тейлора}:
\begin{equation}\label{eq:taylor}
g(x)\;=\;g(x_k)\;+\;g'(x_k)(x-x_k)\;+\;\tfrac12 g''(x_k)(x-x_k)^2
   \;+\;o\!\bigl((x-x_k)^2\bigr).
\end{equation}

Откуда берётся такое разложение? Идея проста: рядом с точкой $x_k$
функция $g$ ведёт себя «почти как многочлен». Линейное приближение
$g(x_k)+g'(x_k)(x-x_k)$ — это уже знакомая нам касательная к
графику. Если мы хотим добиться большей точности, надо учесть и
\emph{кривизну} графика — а она как раз кодируется второй производной
$g''$. Так возникает квадратичный многочлен в~\eqref{eq:taylor}.

Отбросим бесконечно малое слагаемое и обозначим оставшийся квадратичный
многочлен через
\[
q_k(x)\;=\;g(x_k)+g'(x_k)(x-x_k)+\tfrac12 g''(x_k)(x-x_k)^2.
\]
Минимизировать $q_k$ умеет любой школьник: это квадратное уравнение,
решаемое в одну строку. Условие $q_k'(x)=0$ даёт $g'(x_k)+g''(x_k)(x-x_k)=0$,
откуда $x=x_k-g'(x_k)/g''(x_k)$, что в точности совпадает
с~формулой~\eqref{eq:newton-opt}.

\begin{notebox}[title={«Разделяй и властвуй»}]
Метод Ньютона для оптимизации — классический пример принципа
\emph{разделяй и властвуй}: вместо того чтобы решать сложную задачу
минимизации произвольной функции $g$ целиком, мы заменяем её
на каждом шаге простой подзадачей — минимизацией параболы $q_k$.
Парабола подбирается так, чтобы локально (в окрестности $x_k$)
имитировать поведение $g$, и её минимум вычисляется явно.
\end{notebox}

\begin{figure}[ht]
\centering
\includegraphics[width=0.78\linewidth]{taylor_parabola.pdf}
\caption{Метод Ньютона как минимизация локальной квадратичной модели
$q_k$ (пунктир): следующая точка $x_{k+1}$ — это вершина параболы,
касающейся графика $g$ в точке $x_k$.}
\label{fig:taylor}
\end{figure}

\subsection{Пример 1. Извлечение квадратного корня}
\label{ssec:heron}

Пусть $a>0$. Вычислим $\sqrt{a}$, применяя метод Ньютона
к уравнению
\[
f(x)=x^{2}-a=0.
\]
Тогда $f'(x)=2x$, и формула~\eqref{eq:newton-root} принимает вид
\begin{equation}\label{eq:heron}
\boxed{\;x_{k+1}=\dfrac12\!\left(x_k+\dfrac{a}{x_k}\right)\;}
\end{equation}

Это — знаменитый \textbf{метод Герона}, или вавилонский метод
квадратного корня.

\begin{historybox}
Формулу~\eqref{eq:heron} в явном виде описал \emph{Герон Александрийский}
(I век н.~э.) в трактате «Метрика». Однако эту итерацию умели
применять задолго до него: на вавилонской глиняной табличке
\textsf{YBC~7289} (примерно XVIII–XVI вв. до н.~э.) сохранилось
вычисление $\sqrt{2}$ с точностью до шестого десятичного знака —
именно с помощью среднего арифметического и среднего гармонического,
что эквивалентно~\eqref{eq:heron}. Получается, что одна из
итерационных схем «пережила» три с половиной тысячи лет и
по сей день используется в библиотеках вычисления корня.
\end{historybox}

\subsubsection*{Сходимость через сжимающее отображение}

Запишем~\eqref{eq:heron} как $x_{k+1}=T(x_k)$, где
$T(x)=\tfrac12(x+a/x)$. Точка $\sqrt{a}$ — \emph{неподвижная} для $T$:
$T(\sqrt{a})=\sqrt{a}$.

\begin{theorem}{Сходимость метода Герона}{heron}
Для любого $x_0>0$ последовательность~\eqref{eq:heron} монотонно
убывает (начиная с $k=1$) и сходится к $\sqrt{a}$. Более того, для
ошибки $e_k=x_k-\sqrt{a}$ выполнено
\begin{equation}\label{eq:heron-err}
e_{k+1}\;=\;\dfrac{e_k^2}{2x_k}\,, \qquad
\text{и при }k\geq 1\colon\quad
0\le e_{k+1}\le \dfrac{e_k^{2}}{2\sqrt{a}}.
\end{equation}
\end{theorem}

\begin{proof}
\textbf{Шаг 1 (нижняя оценка).} По неравенству о средних
(AM--GM) для $x_0>0$:
$x_1=\tfrac12(x_0+a/x_0)\geq \sqrt{x_0\cdot a/x_0}=\sqrt{a}$.
Значит, начиная с $x_1$, все члены лежат в~$[\sqrt{a},\infty)$.

\textbf{Шаг 2 (рекуррентность для ошибки).}
\[
e_{k+1}=x_{k+1}-\sqrt{a}=\dfrac{x_k+a/x_k}{2}-\sqrt{a}
=\dfrac{x_k^{2}-2\sqrt{a}\,x_k+a}{2x_k}=\dfrac{(x_k-\sqrt{a})^2}{2x_k}=\dfrac{e_k^{2}}{2x_k}.
\]
При $x_k\geq\sqrt{a}$ имеем $\tfrac{1}{2x_k}\leq\tfrac{1}{2\sqrt{a}}$,
откуда вторая часть~\eqref{eq:heron-err}.

\textbf{Шаг 3 (сжатие).} Производная $T'(x)=\tfrac12(1-a/x^{2})$.
На луче $x\geq\sqrt{a}$ выполнено $T'(x)\in[0,\tfrac12)$. Значит,
$T$~— \emph{сжимающее отображение} (с константой~$\tfrac12$) на
$[\sqrt{a},\infty)$. По теореме Банаха о неподвижной точке
последовательность $x_k=T^{k}(x_1)$ сходится к единственной
неподвижной точке~$\sqrt{a}$.
\end{proof}

Из~\eqref{eq:heron-err} видно, что после первой итерации ошибка
\emph{квадрируется}: каждые два десятичных знака удваиваются за
один шаг. Поэтому при разумном $x_0$ хватает 4--5 итераций, чтобы
получить машинную точность.

\begin{example}{Расстояние до горизонта}{horizon}
Стоящий человек ростом $h=1{,}70$~м смотрит на море. До какой точки
далеко он видит? Из теоремы Пифагора (рис.~\ref{fig:horizon}) для прямой,
касательной к поверхности Земли:
\[
d^{2}+R^{2}=(R+h)^{2},\qquad
d=\sqrt{2Rh+h^{2}}\;\approx\;\sqrt{2Rh}.
\]
При $R=6371$~км, $h=1{,}7\cdot 10^{-3}$~км:
\[
2Rh\;\approx\;2\cdot 6371\cdot 0{,}0017\;=\;21{,}66\;\text{(км$^2$)}.
\]
Применим к этому числу метод Герона со стартом $x_0=5$:
\[
\begin{array}{l|l}
k & x_k \\\hline
0 & 5{,}000\,000 \\
1 & \tfrac{1}{2}(5+21{,}66/5)=4{,}666\,000 \\
2 & 4{,}654\,463 \\
3 & 4{,}654\,460
\end{array}
\]
Получаем $d\approx 4{,}65$~км — расстояние, которое запоминается
любому, кто хоть раз стоял у моря.
\end{example}

\begin{figure}[ht]
\centering
\begin{tikzpicture}[scale=0.9, line cap=round]
  % Земля
  \draw[thick, accent] (0,0) circle (2.6);
  \node[accent!80!black] at (0,-0.4) {Земля};
  % центр
  \fill (0,0) circle (1.2pt) node[below right=1pt] {$O$};
  % точка наблюдателя
  \coordinate (A) at (0,2.6);
  \coordinate (P) at (0,3.05);          % человек на высоте h
  \draw[very thick, accent2] (A) -- (P);
  \fill (P) circle (1.2pt) node[left] {$P$};
  \fill (A) circle (1pt) node[left=1pt] {$A$};
  % касательная: из P касается окружности в точке T
  % используем геометрию: |OP|=R+h, расстояние от O до T равно R, угол OTC=90.
  \pgfmathsetmacro{\Rh}{3.05}
  \pgfmathsetmacro{\R}{2.6}
  \pgfmathsetmacro{\dlen}{sqrt(\Rh*\Rh-\R*\R)}    % длина касательной
  \pgfmathsetmacro{\sina}{\dlen/\Rh}
  \pgfmathsetmacro{\cosa}{\R/\Rh}
  \coordinate (T) at ({\R*\sina}, {\R*\cosa});
  \draw[thick, accent2] (P) -- (T);
  \draw[dashed] (0,0) -- (T);
  \fill (T) circle (1pt) node[right=2pt] {$T$};
  % прямой угол в T
  \draw ($(T)+({-0.18*sin(40)},{-0.18*cos(40)})$) -- 
        ++({0.18*cos(40)},{-0.18*sin(40)}) --
        ++({0.18*sin(40)},{0.18*cos(40)});
  % метки
  \node[accent2] at (0.45,2.8) {$h$};
  \node at (0.85,1.45) {$d$};
  \node[rotate=15] at (-0.7,1.4) {$R$};
\end{tikzpicture}
\caption{Расстояние до горизонта $d=\sqrt{(R+h)^2-R^2}\approx\sqrt{2Rh}$.}
\label{fig:horizon}
\end{figure}

\subsection{Пример 2. Вычисление обратного числа}

Что, если на нашем процессоре нет аппаратного деления, и нужно
вычислить $1/a$ при $a>0$ (только сложениями и умножениями)?
Применим метод Ньютона к уравнению
\[
f(x)=\dfrac{1}{x}-a=0.
\]
Тогда $f'(x)=-1/x^{2}$, и подстановка в~\eqref{eq:newton-root} даёт
\[
x_{k+1}=x_k-\dfrac{1/x_k-a}{-1/x_k^{2}}
       =x_k+x_k^{2}\bigl(1/x_k-a\bigr)
       =x_k\bigl(2-a x_k\bigr).
\]
То есть
\begin{equation}\label{eq:ns}
\boxed{\;x_{k+1}=x_k\,(2-a x_k)\;}
\end{equation}

Эта схема замечательна тем, что в ней нет ни одного деления — только
сложение и умножение. Поэтому именно она исторически использовалась
в первых ЭВМ для реализации операции деления на аппаратном уровне.

\subsubsection*{Анализ сходимости}

Введём «относительную невязку» $u_k=1-a x_k$. Тогда
\[
u_{k+1}=1-a x_{k+1}=1-a x_k(2-ax_k)=1-2 a x_k+(a x_k)^{2}=(1-a x_k)^{2}=u_k^{2}.
\]
Таким образом
\begin{equation}\label{eq:ns-error}
u_k=u_0^{\,2^{k}}\quad\Longrightarrow\quad u_k\to 0\;\Leftrightarrow\;|u_0|<1.
\end{equation}

Условие $|1-a x_0|<1$ эквивалентно $0<x_0<2/a$. \emph{Областью
сходимости} метода является интервал $\bigl(0,\,2/a\bigr)$, и
он \textbf{конечен}: вне этого интервала ($x_0>2/a$ или $x_0<0$)
итерации расходятся. Это принципиальное отличие от метода Герона,
который сходится для любого $x_0>0$.

\begin{notebox}[title={Метод Ньютона--Шульца и оптимизатор Muon}]
Если в~\eqref{eq:ns} заменить число $x_k$ на квадратную матрицу $X_k$,
а число $a$ — на матрицу $A$, получим знаменитую матричную итерацию
\textbf{Ньютона--Шульца}: $X_{k+1}=X_k(2I-A X_k)$. Она вычисляет
$A^{-1}$ без обращения матриц на каждом шаге — лишь умножениями.

Аналогичные полиномиальные итерации применяются для приближённой
\emph{ортогонализации} матриц. Именно так устроен оптимизатор
\textbf{Muon}, появившийся в 2024–2025~гг.\ и используемый при
обучении современных больших языковых моделей: в нём матрица
градиента на каждом шаге быстро ортогонализуется несколькими
итерациями Ньютона--Шульца — это намного дешевле, чем точное
сингулярное разложение.
\end{notebox}

\subsection{Реализация на Python и численный эксперимент}

Запишем универсальную функцию, реализующую метод Ньютона
по формуле~\eqref{eq:newton-root}, и применим её к двум разобранным
примерам. Комментарии в коде даны по-английски (как принято в
большинстве библиотек): \texttt{f}, \texttt{df} — сама функция и
её производная; \texttt{x0} — начальное приближение; \texttt{n\_iter} —
максимальное число итераций; \texttt{tol} — точность остановки;
функция возвращает массив \texttt{xs} с историей итераций.

\begin{lstlisting}[label={lst:newton}]
import numpy as np
import matplotlib.pyplot as plt


def newton(f, df, x0, n_iter=20, tol=1e-15):
    """Newton's tangent method.

    Parameters
    ----------
    f, df  : callables f(x) and f'(x).
    x0     : initial guess.
    n_iter : maximum number of iterations.
    tol    : stopping tolerance on |f(x)|.

    Returns
    -------
    Array of iterates, including x0.
    """
    xs = [x0]
    x = x0
    for _ in range(n_iter):
        fx = f(x)
        if abs(fx) < tol:
            break
        x = x - fx / df(x)
        xs.append(x)
    return np.array(xs)


# ---- Example 1. Square root of a ---------------------------------
a = 2.0
xs_sqrt = newton(lambda x: x**2 - a,
                 lambda x: 2 * x,
                 x0=50.0)          # very poor initial guess
err_sqrt = np.abs(xs_sqrt - np.sqrt(a))

# ---- Example 2. Reciprocal 1/a -----------------------------------
a = 3.0
xs_inv  = newton(lambda x: 1.0 / x - a,
                 lambda x: -1.0 / x**2,
                 x0=0.20)          # inside the convergence radius (0, 2/a)
err_inv = np.abs(xs_inv - 1.0 / a)

# ---- Convergence plot in log scale -------------------------------
plt.semilogy(err_sqrt, 'o-', label=r'sqrt: $|x_k-\sqrt{2}|$')
plt.semilogy(err_inv,  's-', label=r'1/a:  $|x_k-1/3|$')
plt.xlabel('iteration $k$'); plt.ylabel('error')
plt.legend(); plt.grid(True); plt.show()
\end{lstlisting}

\noindent\small\textit{Листинг~2.1. Универсальный метод касательных
Ньютона и его применение.}\normalsize

\paragraph{Что показывает эксперимент.}

На рис.~\ref{fig:sqrt-conv} приведены кривые сходимости метода Герона
для $\sqrt{2}$ при разных стартах $x_0\in\{0{,}1;\,1;\,5;\,50\}$.
Все четыре кривые сходятся к нулю, причём поведение типичное для
квадратичной сходимости: после нескольких «согласующих» шагов
число верных знаков удваивается за итерацию. Это иллюстрирует
\textbf{глобальную сходимость} метода Герона: годится \emph{любое}
начальное приближение из $(0,\infty)$.

\begin{figure}[ht]
\centering
\includegraphics[width=0.85\linewidth]{sqrt_convergence.pdf}
\caption{Глобальная сходимость метода Герона для $\sqrt{2}$ из четырёх
очень разных стартовых точек. Шкала ошибки логарифмическая.}
\label{fig:sqrt-conv}
\end{figure}

Совсем другое поведение демонстрирует метод Ньютона--Шульца для
$1/a$ при $a=3$ (рис.~\ref{fig:rec-conv}): область сходимости — это
интервал $(0,\,2/3)\approx(0,\,0{,}667)$.
Стартуя \emph{внутри} этого интервала, метод сходится квадратично
(нижние три кривые); вблизи правой границы (например, $x_0=0{,}66$) — сходится,
но «впритык»; а уже при $x_0=0{,}70$ ошибка взрывообразно растёт:
$|u_k|=|u_0|^{2^{k}}\to\infty$ в полном согласии с
\eqref{eq:ns-error}. Это \textbf{локальная сходимость}.

\begin{figure}[ht]
\centering
\includegraphics[width=0.85\linewidth]{reciprocal_convergence.pdf}
\caption{Локальная сходимость метода Ньютона--Шульца для $1/3$.
Область сходимости — интервал $(0,\,2/3)$; вне его метод расходится.}
\label{fig:rec-conv}
\end{figure}

\subsection{Достаточные условия сверхлинейной сходимости\texorpdfstring{$^\star$}{*}}
\label{ssec:thm-superlinear}

Сравнение двух примеров наводит на вопрос: \emph{когда именно}
метод Ньютона сходится так быстро, как мы наблюдаем? Ответ
даёт следующая теорема, восходящая к Канторовичу.

\begin{definition}{Сверхлинейная и квадратичная сходимости}{rates}
Говорят, что последовательность $\{x_k\}$ сходится к $x^{\star}$
со \textbf{сверхлинейной скоростью}, если
\[
\lim_{k\to\infty}\dfrac{|x_{k+1}-x^{\star}|}{|x_k-x^{\star}|}=0.
\]
Если же выполнено более сильное условие
$|x_{k+1}-x^{\star}|\leq M |x_k-x^{\star}|^{2}$ для некоторой
константы $M>0$, то сходимость называют \textbf{квадратичной}.
\end{definition}

\begin{definition}{Сильно выпуклая функция}{strconvex}
Функция $g\colon\R\to\R$ называется \textbf{сильно выпуклой} с
константой $\mu>0$, если $g\in C^{2}$ и $g''(x)\geq\mu$ для всех $x$.
Сильная выпуклость гарантирует существование и единственность
точки минимума и, что важно для нас, отделённость второй производной
от нуля.
\end{definition}

\begin{theorem}{Сверхлинейная сходимость метода Ньютона}{newton-rate}
Пусть выполнено хотя бы одно из условий:
\begin{enumerate}\itemsep0pt
  \item[(а)] (\emph{уравнение}) $f\in C^{2}$ в окрестности точки $x^{\star}$,
        $f(x^{\star})=0$ и $f'(x^{\star})\neq 0$;
  \item[(б)] (\emph{сильно выпуклая оптимизация}) $g\in C^{3}$,
        $g$~— сильно выпуклая на интервале, содержащем точку
        минимума $x^{\star}$, и применяется итерация~\eqref{eq:newton-opt}.
\end{enumerate}
Тогда существуют $\delta>0$ и $M>0$ такие, что для всякого
$x_0\in(x^{\star}-\delta,\,x^{\star}+\delta)$ итерации Ньютона
определены и
\begin{equation}\label{eq:rate}
|x_{k+1}-x^{\star}|\;\leq\;M\,|x_k-x^{\star}|^{2}.
\end{equation}
В частности, сходимость квадратичная (а значит, сверхлинейная).
\end{theorem}

\begin{proof}
Достаточно доказать случай (а): задача (б) сводится к нему,
если положить $f=g'$, тогда условие $f'(x^{\star})\neq 0$
эквивалентно $g''(x^{\star})\geq\mu>0$.

Разложим $f$ по формуле Тейлора с остаточным членом в форме
Лагранжа в окрестности~$x_k$ и подставим $x=x^{\star}$:
\[
0\;=\;f(x^{\star})\;=\;f(x_k)+f'(x_k)(x^{\star}-x_k)+\tfrac{f''(\xi_k)}{2}(x^{\star}-x_k)^{2},
\]
где $\xi_k$~— некоторая точка между $x_k$ и $x^{\star}$. Разделим
обе части на $f'(x_k)$ (по непрерывности эта величина не равна
нулю в малой окрестности $x^{\star}$):
\[
0\;=\;\dfrac{f(x_k)}{f'(x_k)}+(x^{\star}-x_k)+\dfrac{f''(\xi_k)}{2 f'(x_k)}(x^{\star}-x_k)^{2}.
\]
Из определения $x_{k+1}=x_k-f(x_k)/f'(x_k)$ получаем
$\frac{f(x_k)}{f'(x_k)}=x_k-x_{k+1}$. Подставляя:
\[
0=(x_k-x_{k+1})+(x^{\star}-x_k)+\dfrac{f''(\xi_k)}{2 f'(x_k)}(x^{\star}-x_k)^{2},
\]
откуда после очевидных переобозначений
\begin{equation}\label{eq:newton-recur}
x_{k+1}-x^{\star}\;=\;\dfrac{f''(\xi_k)}{2 f'(x_k)}\,(x_k-x^{\star})^{2}.
\end{equation}
Это и есть искомая \emph{рекуррентность для расстояния}.

Теперь подберём константы. По непрерывности $f'$ и $f''$ существует
$\delta_1>0$ и числа $L$, $m>0$ такие, что
$|f''(\xi)|\leq L$ и $|f'(x)|\geq m$ для всех $\xi,x\in
[x^{\star}-\delta_1,\,x^{\star}+\delta_1]$.
Положим $M=L/(2m)$ и $\delta=\min\!\bigl(\delta_1,\,\tfrac{1}{2M}\bigr)$.
Тогда из~\eqref{eq:newton-recur} при $|x_k-x^{\star}|<\delta$:
\[
|x_{k+1}-x^{\star}|\leq M|x_k-x^{\star}|^{2}\leq M\delta\cdot|x_k-x^{\star}|\leq\tfrac12|x_k-x^{\star}|<\delta.
\]
Значит, последовательность не покидает $\delta$-окрестность и
оценка~\eqref{eq:rate} выполнена для всех $k$.
\end{proof}

\paragraph{Где условия выполняются для наших примеров?}

\textbf{Метод Герона.} Здесь $f(x)=x^{2}-a$, $f'(x)=2x$,
$f''(x)=2$. В точке $x^{\star}=\sqrt{a}>0$ выполнено $f'(x^{\star})=2\sqrt{a}\neq 0$.
\emph{Условия теоремы выполняются на всём луче $(0,\infty)$,
не содержащем ноль}, и константа в~\eqref{eq:rate} принимает вид
$M=\frac{1}{2\sqrt{a}}$ — в полном согласии с
явным выражением~\eqref{eq:heron-err}, полученным элементарно.

\textbf{Метод Ньютона--Шульца.} Здесь $f(x)=1/x-a$, $f'(x)=-1/x^{2}$,
$f''(x)=2/x^{3}$. В точке $x^{\star}=1/a$ имеем $f'(x^{\star})=-a^{2}\neq 0$,
и теорема применима в малой окрестности этой точки. Однако, в
отличие от метода Герона, оценка снизу $|f'(x)|\geq m>0$ \emph{ломается}
при $x\to\infty$: производная стремится к нулю. Именно поэтому
радиус сходимости в этом примере конечен и равен $1/a$ слева и справа
от $x^{\star}$, что в сумме даёт интервал $(0,\,2/a)$, найденный нами
ранее аналитически.

\medskip

Таким образом, два простейших примера полностью укладываются в
доказанную теорему: метод Ньютона сходится сверхлинейно \emph{там,
где} производная регулярна и отделена от нуля. Глобальную же
сходимость, как у Герона, приходится доказывать дополнительно — например,
ссылаясь на структуру функции (выпуклость) и теорему Банаха о
сжимающем отображении, как мы и сделали в разделе~\ref{ssec:heron}.

\begin{tasksbox}
\begin{enumerate}\itemsep1pt
  \item Применяя метод Ньютона к уравнению $f(x)=x^{n}-a=0$, получите
        итерационную формулу для $\sqrt[n]{a}$. Покажите, что для $n=3$
        она имеет вид $x_{k+1}=\frac{2 x_k+a/x_k^{2}}{3}$.
  \item Докажите, что для метода Герона из любого $x_0>0$ погрешность
        $e_1$ всегда положительна, после чего последовательность
        монотонно убывает. (Указание: используйте AM--GM.)
  \item Для метода Ньютона--Шульца при $a=2$ найдите наибольшее $x_0$,
        при котором итерации ещё сходятся. Сравните численно скорость
        сходимости при $x_0=0{,}1$ и $x_0=0{,}9$.
  \item$^{\star}$ Сформулируйте и докажите аналог
        теоремы~\ref{thm:newton-rate} для метода Ньютона минимизации
        многомерной сильно выпуклой функции $g\colon\R^{n}\to\R$.
        Какую роль играет матрица Гессе $\nabla^{2}g$?
  \item$^{\star}$ Прочитайте о теореме Канторовича в учебнике
        \href{https://en.wikipedia.org/wiki/Kantorovich_theorem}{«Kantorovich theorem»}
        и сформулируйте её в одну фразу: какие именно три величины
        должны быть «в нужном соотношении» для гарантии сходимости?
\end{enumerate}
\end{tasksbox}

%      % =====================================================================
%  Глава «Вычислительная линейная алгебра».
%  Самодостаточные сюжеты: SVD, Eigenfaces/Eigencats, PageRank,
%  оценки Липшицевости нейросети, числа с плавающей точкой.
%  Источники: nla360, optim, ds-intro.
%
%  Назначение в книге: новая глава Части II «Математические основы».
%  Если будет принято решение о вставке между Ньютоном (текущая гл. 2)
%  и Частью III, главе автоматически присвоится номер 3 (а текущая
%  глава 3 станет 4-й и т. д.). Все ссылки внутри файла — через \ref,
%  поэтому перенумерация происходит автоматически.
%
%  Подключение из main.tex (предлагаемая позиция):
%      \part{Математические основы}
%      \chapter{Численные методы оптимизации}
\label{ch:opt}

\section{Метод касательных Ньютона и его приложения}
\label{sec:newton}

\subsection{От Ньютона до Канторовича}

Задачи отыскания корня уравнения $f(x)=0$ и поиска минимума функции
$g(x)$ относятся к разным «жанрам» школьной математики, однако
численно они решаются почти одним и тем же приёмом — итерационной
схемой, придуманной Исааком Ньютоном.

\begin{historybox}
В работе \emph{De analysi per aequationes numero terminorum infinitas}
(1669) и затем в \emph{Method of Fluxions} (рукопись 1671~г., издание
1736~г.) Ньютон предложил способ приближённого решения уравнений,
который позже Джозеф Рафсон в 1690~г. переписал в современной
рекуррентной форме. Поэтому в англоязычной литературе метод часто
называют \emph{Newton–Raphson}.

Существенный вклад в строгое исследование метода внёс советский
математик \textbf{Леонид Витальевич Канторович} (1912–1986),
лауреат Нобелевской премии по экономике 1975~г. (за теорию
оптимального распределения ресурсов). В 1948~г. он опубликовал
теорему, носящую теперь его имя: она впервые дала универсальные
\emph{достаточные условия сходимости} метода Ньютона, причём не
только для уравнений с одной переменной, но и для уравнений в
бесконечномерных банаховых пространствах.
\end{historybox}

\paragraph{Геометрическая идея.} Пусть требуется найти корень
уравнения~$f(x)=0$, и нам известна точка $x_k$, близкая к корню.
Заменим график $y=f(x)$ его \emph{касательной} в точке $(x_k, f(x_k))$.
Касательная — это прямая, и её пересечение с осью $x$ найти просто.
Точку этого пересечения и объявим следующим приближением~$x_{k+1}$.
Затем повторим. Так получается метод касательных~(рис.~\ref{fig:newton-geom}).

\begin{figure}[ht]
\centering
\includegraphics[width=0.85\linewidth]{newton_geometry.pdf}
\caption{Три итерации метода Ньютона для уравнения
$x^{3}-2x-5=0$. Каждая касательная (пунктир) пересекает ось абсцисс
в точке, которая берётся в качестве следующего приближения
(квадрат). После трёх шагов корень найден с точностью лучше $10^{-6}$.}
\label{fig:newton-geom}
\end{figure}

Уравнение касательной к графику $y=f(x)$ в точке $x_k$ имеет вид
$y=f(x_k)+f'(x_k)\,(x-x_k)$.
Полагая $y=0$ и выражая $x$, получаем
\begin{equation}\label{eq:newton-root}
\boxed{\;x_{k+1}=x_k-\dfrac{f(x_k)}{f'(x_k)}\;}\qquad k=0,1,2,\dots
\end{equation}

Это и есть \emph{итерационная формула Ньютона}. Точка $x_0$ выбирается
исследователем; от её выбора может зависеть, сойдётся ли метод
вообще.

\subsection{Метод Ньютона для задач оптимизации}

Пусть теперь нужно найти минимум функции $g$. В любой точке минимума
$x^{\star}$ выполнено $g'(x^{\star})=0$. Значит, минимизация
дифференцируемой функции — это решение уравнения
\[
g'(x)=0.
\]
Применим к нему формулу~\eqref{eq:newton-root} с $f=g'$ (и тогда $f'=g''$):
\begin{equation}\label{eq:newton-opt}
\boxed{\;x_{k+1}=x_k-\dfrac{g'(x_k)}{g''(x_k)}\;}
\end{equation}

Эту же формулу можно получить иначе — и это ценная для всей теории
оптимизации точка зрения. Разложим $g$ в окрестности~$x_k$ в
\textbf{ряд Тейлора}:
\begin{equation}\label{eq:taylor}
g(x)\;=\;g(x_k)\;+\;g'(x_k)(x-x_k)\;+\;\tfrac12 g''(x_k)(x-x_k)^2
   \;+\;o\!\bigl((x-x_k)^2\bigr).
\end{equation}

Откуда берётся такое разложение? Идея проста: рядом с точкой $x_k$
функция $g$ ведёт себя «почти как многочлен». Линейное приближение
$g(x_k)+g'(x_k)(x-x_k)$ — это уже знакомая нам касательная к
графику. Если мы хотим добиться большей точности, надо учесть и
\emph{кривизну} графика — а она как раз кодируется второй производной
$g''$. Так возникает квадратичный многочлен в~\eqref{eq:taylor}.

Отбросим бесконечно малое слагаемое и обозначим оставшийся квадратичный
многочлен через
\[
q_k(x)\;=\;g(x_k)+g'(x_k)(x-x_k)+\tfrac12 g''(x_k)(x-x_k)^2.
\]
Минимизировать $q_k$ умеет любой школьник: это квадратное уравнение,
решаемое в одну строку. Условие $q_k'(x)=0$ даёт $g'(x_k)+g''(x_k)(x-x_k)=0$,
откуда $x=x_k-g'(x_k)/g''(x_k)$, что в точности совпадает
с~формулой~\eqref{eq:newton-opt}.

\begin{notebox}[title={«Разделяй и властвуй»}]
Метод Ньютона для оптимизации — классический пример принципа
\emph{разделяй и властвуй}: вместо того чтобы решать сложную задачу
минимизации произвольной функции $g$ целиком, мы заменяем её
на каждом шаге простой подзадачей — минимизацией параболы $q_k$.
Парабола подбирается так, чтобы локально (в окрестности $x_k$)
имитировать поведение $g$, и её минимум вычисляется явно.
\end{notebox}

\begin{figure}[ht]
\centering
\includegraphics[width=0.78\linewidth]{taylor_parabola.pdf}
\caption{Метод Ньютона как минимизация локальной квадратичной модели
$q_k$ (пунктир): следующая точка $x_{k+1}$ — это вершина параболы,
касающейся графика $g$ в точке $x_k$.}
\label{fig:taylor}
\end{figure}

\subsection{Пример 1. Извлечение квадратного корня}
\label{ssec:heron}

Пусть $a>0$. Вычислим $\sqrt{a}$, применяя метод Ньютона
к уравнению
\[
f(x)=x^{2}-a=0.
\]
Тогда $f'(x)=2x$, и формула~\eqref{eq:newton-root} принимает вид
\begin{equation}\label{eq:heron}
\boxed{\;x_{k+1}=\dfrac12\!\left(x_k+\dfrac{a}{x_k}\right)\;}
\end{equation}

Это — знаменитый \textbf{метод Герона}, или вавилонский метод
квадратного корня.

\begin{historybox}
Формулу~\eqref{eq:heron} в явном виде описал \emph{Герон Александрийский}
(I век н.~э.) в трактате «Метрика». Однако эту итерацию умели
применять задолго до него: на вавилонской глиняной табличке
\textsf{YBC~7289} (примерно XVIII–XVI вв. до н.~э.) сохранилось
вычисление $\sqrt{2}$ с точностью до шестого десятичного знака —
именно с помощью среднего арифметического и среднего гармонического,
что эквивалентно~\eqref{eq:heron}. Получается, что одна из
итерационных схем «пережила» три с половиной тысячи лет и
по сей день используется в библиотеках вычисления корня.
\end{historybox}

\subsubsection*{Сходимость через сжимающее отображение}

Запишем~\eqref{eq:heron} как $x_{k+1}=T(x_k)$, где
$T(x)=\tfrac12(x+a/x)$. Точка $\sqrt{a}$ — \emph{неподвижная} для $T$:
$T(\sqrt{a})=\sqrt{a}$.

\begin{theorem}{Сходимость метода Герона}{heron}
Для любого $x_0>0$ последовательность~\eqref{eq:heron} монотонно
убывает (начиная с $k=1$) и сходится к $\sqrt{a}$. Более того, для
ошибки $e_k=x_k-\sqrt{a}$ выполнено
\begin{equation}\label{eq:heron-err}
e_{k+1}\;=\;\dfrac{e_k^2}{2x_k}\,, \qquad
\text{и при }k\geq 1\colon\quad
0\le e_{k+1}\le \dfrac{e_k^{2}}{2\sqrt{a}}.
\end{equation}
\end{theorem}

\begin{proof}
\textbf{Шаг 1 (нижняя оценка).} По неравенству о средних
(AM--GM) для $x_0>0$:
$x_1=\tfrac12(x_0+a/x_0)\geq \sqrt{x_0\cdot a/x_0}=\sqrt{a}$.
Значит, начиная с $x_1$, все члены лежат в~$[\sqrt{a},\infty)$.

\textbf{Шаг 2 (рекуррентность для ошибки).}
\[
e_{k+1}=x_{k+1}-\sqrt{a}=\dfrac{x_k+a/x_k}{2}-\sqrt{a}
=\dfrac{x_k^{2}-2\sqrt{a}\,x_k+a}{2x_k}=\dfrac{(x_k-\sqrt{a})^2}{2x_k}=\dfrac{e_k^{2}}{2x_k}.
\]
При $x_k\geq\sqrt{a}$ имеем $\tfrac{1}{2x_k}\leq\tfrac{1}{2\sqrt{a}}$,
откуда вторая часть~\eqref{eq:heron-err}.

\textbf{Шаг 3 (сжатие).} Производная $T'(x)=\tfrac12(1-a/x^{2})$.
На луче $x\geq\sqrt{a}$ выполнено $T'(x)\in[0,\tfrac12)$. Значит,
$T$~— \emph{сжимающее отображение} (с константой~$\tfrac12$) на
$[\sqrt{a},\infty)$. По теореме Банаха о неподвижной точке
последовательность $x_k=T^{k}(x_1)$ сходится к единственной
неподвижной точке~$\sqrt{a}$.
\end{proof}

Из~\eqref{eq:heron-err} видно, что после первой итерации ошибка
\emph{квадрируется}: каждые два десятичных знака удваиваются за
один шаг. Поэтому при разумном $x_0$ хватает 4--5 итераций, чтобы
получить машинную точность.

\begin{example}{Расстояние до горизонта}{horizon}
Стоящий человек ростом $h=1{,}70$~м смотрит на море. До какой точки
далеко он видит? Из теоремы Пифагора (рис.~\ref{fig:horizon}) для прямой,
касательной к поверхности Земли:
\[
d^{2}+R^{2}=(R+h)^{2},\qquad
d=\sqrt{2Rh+h^{2}}\;\approx\;\sqrt{2Rh}.
\]
При $R=6371$~км, $h=1{,}7\cdot 10^{-3}$~км:
\[
2Rh\;\approx\;2\cdot 6371\cdot 0{,}0017\;=\;21{,}66\;\text{(км$^2$)}.
\]
Применим к этому числу метод Герона со стартом $x_0=5$:
\[
\begin{array}{l|l}
k & x_k \\\hline
0 & 5{,}000\,000 \\
1 & \tfrac{1}{2}(5+21{,}66/5)=4{,}666\,000 \\
2 & 4{,}654\,463 \\
3 & 4{,}654\,460
\end{array}
\]
Получаем $d\approx 4{,}65$~км — расстояние, которое запоминается
любому, кто хоть раз стоял у моря.
\end{example}

\begin{figure}[ht]
\centering
\begin{tikzpicture}[scale=0.9, line cap=round]
  % Земля
  \draw[thick, accent] (0,0) circle (2.6);
  \node[accent!80!black] at (0,-0.4) {Земля};
  % центр
  \fill (0,0) circle (1.2pt) node[below right=1pt] {$O$};
  % точка наблюдателя
  \coordinate (A) at (0,2.6);
  \coordinate (P) at (0,3.05);          % человек на высоте h
  \draw[very thick, accent2] (A) -- (P);
  \fill (P) circle (1.2pt) node[left] {$P$};
  \fill (A) circle (1pt) node[left=1pt] {$A$};
  % касательная: из P касается окружности в точке T
  % используем геометрию: |OP|=R+h, расстояние от O до T равно R, угол OTC=90.
  \pgfmathsetmacro{\Rh}{3.05}
  \pgfmathsetmacro{\R}{2.6}
  \pgfmathsetmacro{\dlen}{sqrt(\Rh*\Rh-\R*\R)}    % длина касательной
  \pgfmathsetmacro{\sina}{\dlen/\Rh}
  \pgfmathsetmacro{\cosa}{\R/\Rh}
  \coordinate (T) at ({\R*\sina}, {\R*\cosa});
  \draw[thick, accent2] (P) -- (T);
  \draw[dashed] (0,0) -- (T);
  \fill (T) circle (1pt) node[right=2pt] {$T$};
  % прямой угол в T
  \draw ($(T)+({-0.18*sin(40)},{-0.18*cos(40)})$) -- 
        ++({0.18*cos(40)},{-0.18*sin(40)}) --
        ++({0.18*sin(40)},{0.18*cos(40)});
  % метки
  \node[accent2] at (0.45,2.8) {$h$};
  \node at (0.85,1.45) {$d$};
  \node[rotate=15] at (-0.7,1.4) {$R$};
\end{tikzpicture}
\caption{Расстояние до горизонта $d=\sqrt{(R+h)^2-R^2}\approx\sqrt{2Rh}$.}
\label{fig:horizon}
\end{figure}

\subsection{Пример 2. Вычисление обратного числа}

Что, если на нашем процессоре нет аппаратного деления, и нужно
вычислить $1/a$ при $a>0$ (только сложениями и умножениями)?
Применим метод Ньютона к уравнению
\[
f(x)=\dfrac{1}{x}-a=0.
\]
Тогда $f'(x)=-1/x^{2}$, и подстановка в~\eqref{eq:newton-root} даёт
\[
x_{k+1}=x_k-\dfrac{1/x_k-a}{-1/x_k^{2}}
       =x_k+x_k^{2}\bigl(1/x_k-a\bigr)
       =x_k\bigl(2-a x_k\bigr).
\]
То есть
\begin{equation}\label{eq:ns}
\boxed{\;x_{k+1}=x_k\,(2-a x_k)\;}
\end{equation}

Эта схема замечательна тем, что в ней нет ни одного деления — только
сложение и умножение. Поэтому именно она исторически использовалась
в первых ЭВМ для реализации операции деления на аппаратном уровне.

\subsubsection*{Анализ сходимости}

Введём «относительную невязку» $u_k=1-a x_k$. Тогда
\[
u_{k+1}=1-a x_{k+1}=1-a x_k(2-ax_k)=1-2 a x_k+(a x_k)^{2}=(1-a x_k)^{2}=u_k^{2}.
\]
Таким образом
\begin{equation}\label{eq:ns-error}
u_k=u_0^{\,2^{k}}\quad\Longrightarrow\quad u_k\to 0\;\Leftrightarrow\;|u_0|<1.
\end{equation}

Условие $|1-a x_0|<1$ эквивалентно $0<x_0<2/a$. \emph{Областью
сходимости} метода является интервал $\bigl(0,\,2/a\bigr)$, и
он \textbf{конечен}: вне этого интервала ($x_0>2/a$ или $x_0<0$)
итерации расходятся. Это принципиальное отличие от метода Герона,
который сходится для любого $x_0>0$.

\begin{notebox}[title={Метод Ньютона--Шульца и оптимизатор Muon}]
Если в~\eqref{eq:ns} заменить число $x_k$ на квадратную матрицу $X_k$,
а число $a$ — на матрицу $A$, получим знаменитую матричную итерацию
\textbf{Ньютона--Шульца}: $X_{k+1}=X_k(2I-A X_k)$. Она вычисляет
$A^{-1}$ без обращения матриц на каждом шаге — лишь умножениями.

Аналогичные полиномиальные итерации применяются для приближённой
\emph{ортогонализации} матриц. Именно так устроен оптимизатор
\textbf{Muon}, появившийся в 2024–2025~гг.\ и используемый при
обучении современных больших языковых моделей: в нём матрица
градиента на каждом шаге быстро ортогонализуется несколькими
итерациями Ньютона--Шульца — это намного дешевле, чем точное
сингулярное разложение.
\end{notebox}

\subsection{Реализация на Python и численный эксперимент}

Запишем универсальную функцию, реализующую метод Ньютона
по формуле~\eqref{eq:newton-root}, и применим её к двум разобранным
примерам. Комментарии в коде даны по-английски (как принято в
большинстве библиотек): \texttt{f}, \texttt{df} — сама функция и
её производная; \texttt{x0} — начальное приближение; \texttt{n\_iter} —
максимальное число итераций; \texttt{tol} — точность остановки;
функция возвращает массив \texttt{xs} с историей итераций.

\begin{lstlisting}[label={lst:newton}]
import numpy as np
import matplotlib.pyplot as plt


def newton(f, df, x0, n_iter=20, tol=1e-15):
    """Newton's tangent method.

    Parameters
    ----------
    f, df  : callables f(x) and f'(x).
    x0     : initial guess.
    n_iter : maximum number of iterations.
    tol    : stopping tolerance on |f(x)|.

    Returns
    -------
    Array of iterates, including x0.
    """
    xs = [x0]
    x = x0
    for _ in range(n_iter):
        fx = f(x)
        if abs(fx) < tol:
            break
        x = x - fx / df(x)
        xs.append(x)
    return np.array(xs)


# ---- Example 1. Square root of a ---------------------------------
a = 2.0
xs_sqrt = newton(lambda x: x**2 - a,
                 lambda x: 2 * x,
                 x0=50.0)          # very poor initial guess
err_sqrt = np.abs(xs_sqrt - np.sqrt(a))

# ---- Example 2. Reciprocal 1/a -----------------------------------
a = 3.0
xs_inv  = newton(lambda x: 1.0 / x - a,
                 lambda x: -1.0 / x**2,
                 x0=0.20)          # inside the convergence radius (0, 2/a)
err_inv = np.abs(xs_inv - 1.0 / a)

# ---- Convergence plot in log scale -------------------------------
plt.semilogy(err_sqrt, 'o-', label=r'sqrt: $|x_k-\sqrt{2}|$')
plt.semilogy(err_inv,  's-', label=r'1/a:  $|x_k-1/3|$')
plt.xlabel('iteration $k$'); plt.ylabel('error')
plt.legend(); plt.grid(True); plt.show()
\end{lstlisting}

\noindent\small\textit{Листинг~2.1. Универсальный метод касательных
Ньютона и его применение.}\normalsize

\paragraph{Что показывает эксперимент.}

На рис.~\ref{fig:sqrt-conv} приведены кривые сходимости метода Герона
для $\sqrt{2}$ при разных стартах $x_0\in\{0{,}1;\,1;\,5;\,50\}$.
Все четыре кривые сходятся к нулю, причём поведение типичное для
квадратичной сходимости: после нескольких «согласующих» шагов
число верных знаков удваивается за итерацию. Это иллюстрирует
\textbf{глобальную сходимость} метода Герона: годится \emph{любое}
начальное приближение из $(0,\infty)$.

\begin{figure}[ht]
\centering
\includegraphics[width=0.85\linewidth]{sqrt_convergence.pdf}
\caption{Глобальная сходимость метода Герона для $\sqrt{2}$ из четырёх
очень разных стартовых точек. Шкала ошибки логарифмическая.}
\label{fig:sqrt-conv}
\end{figure}

Совсем другое поведение демонстрирует метод Ньютона--Шульца для
$1/a$ при $a=3$ (рис.~\ref{fig:rec-conv}): область сходимости — это
интервал $(0,\,2/3)\approx(0,\,0{,}667)$.
Стартуя \emph{внутри} этого интервала, метод сходится квадратично
(нижние три кривые); вблизи правой границы (например, $x_0=0{,}66$) — сходится,
но «впритык»; а уже при $x_0=0{,}70$ ошибка взрывообразно растёт:
$|u_k|=|u_0|^{2^{k}}\to\infty$ в полном согласии с
\eqref{eq:ns-error}. Это \textbf{локальная сходимость}.

\begin{figure}[ht]
\centering
\includegraphics[width=0.85\linewidth]{reciprocal_convergence.pdf}
\caption{Локальная сходимость метода Ньютона--Шульца для $1/3$.
Область сходимости — интервал $(0,\,2/3)$; вне его метод расходится.}
\label{fig:rec-conv}
\end{figure}

\subsection{Достаточные условия сверхлинейной сходимости\texorpdfstring{$^\star$}{*}}
\label{ssec:thm-superlinear}

Сравнение двух примеров наводит на вопрос: \emph{когда именно}
метод Ньютона сходится так быстро, как мы наблюдаем? Ответ
даёт следующая теорема, восходящая к Канторовичу.

\begin{definition}{Сверхлинейная и квадратичная сходимости}{rates}
Говорят, что последовательность $\{x_k\}$ сходится к $x^{\star}$
со \textbf{сверхлинейной скоростью}, если
\[
\lim_{k\to\infty}\dfrac{|x_{k+1}-x^{\star}|}{|x_k-x^{\star}|}=0.
\]
Если же выполнено более сильное условие
$|x_{k+1}-x^{\star}|\leq M |x_k-x^{\star}|^{2}$ для некоторой
константы $M>0$, то сходимость называют \textbf{квадратичной}.
\end{definition}

\begin{definition}{Сильно выпуклая функция}{strconvex}
Функция $g\colon\R\to\R$ называется \textbf{сильно выпуклой} с
константой $\mu>0$, если $g\in C^{2}$ и $g''(x)\geq\mu$ для всех $x$.
Сильная выпуклость гарантирует существование и единственность
точки минимума и, что важно для нас, отделённость второй производной
от нуля.
\end{definition}

\begin{theorem}{Сверхлинейная сходимость метода Ньютона}{newton-rate}
Пусть выполнено хотя бы одно из условий:
\begin{enumerate}\itemsep0pt
  \item[(а)] (\emph{уравнение}) $f\in C^{2}$ в окрестности точки $x^{\star}$,
        $f(x^{\star})=0$ и $f'(x^{\star})\neq 0$;
  \item[(б)] (\emph{сильно выпуклая оптимизация}) $g\in C^{3}$,
        $g$~— сильно выпуклая на интервале, содержащем точку
        минимума $x^{\star}$, и применяется итерация~\eqref{eq:newton-opt}.
\end{enumerate}
Тогда существуют $\delta>0$ и $M>0$ такие, что для всякого
$x_0\in(x^{\star}-\delta,\,x^{\star}+\delta)$ итерации Ньютона
определены и
\begin{equation}\label{eq:rate}
|x_{k+1}-x^{\star}|\;\leq\;M\,|x_k-x^{\star}|^{2}.
\end{equation}
В частности, сходимость квадратичная (а значит, сверхлинейная).
\end{theorem}

\begin{proof}
Достаточно доказать случай (а): задача (б) сводится к нему,
если положить $f=g'$, тогда условие $f'(x^{\star})\neq 0$
эквивалентно $g''(x^{\star})\geq\mu>0$.

Разложим $f$ по формуле Тейлора с остаточным членом в форме
Лагранжа в окрестности~$x_k$ и подставим $x=x^{\star}$:
\[
0\;=\;f(x^{\star})\;=\;f(x_k)+f'(x_k)(x^{\star}-x_k)+\tfrac{f''(\xi_k)}{2}(x^{\star}-x_k)^{2},
\]
где $\xi_k$~— некоторая точка между $x_k$ и $x^{\star}$. Разделим
обе части на $f'(x_k)$ (по непрерывности эта величина не равна
нулю в малой окрестности $x^{\star}$):
\[
0\;=\;\dfrac{f(x_k)}{f'(x_k)}+(x^{\star}-x_k)+\dfrac{f''(\xi_k)}{2 f'(x_k)}(x^{\star}-x_k)^{2}.
\]
Из определения $x_{k+1}=x_k-f(x_k)/f'(x_k)$ получаем
$\frac{f(x_k)}{f'(x_k)}=x_k-x_{k+1}$. Подставляя:
\[
0=(x_k-x_{k+1})+(x^{\star}-x_k)+\dfrac{f''(\xi_k)}{2 f'(x_k)}(x^{\star}-x_k)^{2},
\]
откуда после очевидных переобозначений
\begin{equation}\label{eq:newton-recur}
x_{k+1}-x^{\star}\;=\;\dfrac{f''(\xi_k)}{2 f'(x_k)}\,(x_k-x^{\star})^{2}.
\end{equation}
Это и есть искомая \emph{рекуррентность для расстояния}.

Теперь подберём константы. По непрерывности $f'$ и $f''$ существует
$\delta_1>0$ и числа $L$, $m>0$ такие, что
$|f''(\xi)|\leq L$ и $|f'(x)|\geq m$ для всех $\xi,x\in
[x^{\star}-\delta_1,\,x^{\star}+\delta_1]$.
Положим $M=L/(2m)$ и $\delta=\min\!\bigl(\delta_1,\,\tfrac{1}{2M}\bigr)$.
Тогда из~\eqref{eq:newton-recur} при $|x_k-x^{\star}|<\delta$:
\[
|x_{k+1}-x^{\star}|\leq M|x_k-x^{\star}|^{2}\leq M\delta\cdot|x_k-x^{\star}|\leq\tfrac12|x_k-x^{\star}|<\delta.
\]
Значит, последовательность не покидает $\delta$-окрестность и
оценка~\eqref{eq:rate} выполнена для всех $k$.
\end{proof}

\paragraph{Где условия выполняются для наших примеров?}

\textbf{Метод Герона.} Здесь $f(x)=x^{2}-a$, $f'(x)=2x$,
$f''(x)=2$. В точке $x^{\star}=\sqrt{a}>0$ выполнено $f'(x^{\star})=2\sqrt{a}\neq 0$.
\emph{Условия теоремы выполняются на всём луче $(0,\infty)$,
не содержащем ноль}, и константа в~\eqref{eq:rate} принимает вид
$M=\frac{1}{2\sqrt{a}}$ — в полном согласии с
явным выражением~\eqref{eq:heron-err}, полученным элементарно.

\textbf{Метод Ньютона--Шульца.} Здесь $f(x)=1/x-a$, $f'(x)=-1/x^{2}$,
$f''(x)=2/x^{3}$. В точке $x^{\star}=1/a$ имеем $f'(x^{\star})=-a^{2}\neq 0$,
и теорема применима в малой окрестности этой точки. Однако, в
отличие от метода Герона, оценка снизу $|f'(x)|\geq m>0$ \emph{ломается}
при $x\to\infty$: производная стремится к нулю. Именно поэтому
радиус сходимости в этом примере конечен и равен $1/a$ слева и справа
от $x^{\star}$, что в сумме даёт интервал $(0,\,2/a)$, найденный нами
ранее аналитически.

\medskip

Таким образом, два простейших примера полностью укладываются в
доказанную теорему: метод Ньютона сходится сверхлинейно \emph{там,
где} производная регулярна и отделена от нуля. Глобальную же
сходимость, как у Герона, приходится доказывать дополнительно — например,
ссылаясь на структуру функции (выпуклость) и теорему Банаха о
сжимающем отображении, как мы и сделали в разделе~\ref{ssec:heron}.

\begin{tasksbox}
\begin{enumerate}\itemsep1pt
  \item Применяя метод Ньютона к уравнению $f(x)=x^{n}-a=0$, получите
        итерационную формулу для $\sqrt[n]{a}$. Покажите, что для $n=3$
        она имеет вид $x_{k+1}=\frac{2 x_k+a/x_k^{2}}{3}$.
  \item Докажите, что для метода Герона из любого $x_0>0$ погрешность
        $e_1$ всегда положительна, после чего последовательность
        монотонно убывает. (Указание: используйте AM--GM.)
  \item Для метода Ньютона--Шульца при $a=2$ найдите наибольшее $x_0$,
        при котором итерации ещё сходятся. Сравните численно скорость
        сходимости при $x_0=0{,}1$ и $x_0=0{,}9$.
  \item$^{\star}$ Сформулируйте и докажите аналог
        теоремы~\ref{thm:newton-rate} для метода Ньютона минимизации
        многомерной сильно выпуклой функции $g\colon\R^{n}\to\R$.
        Какую роль играет матрица Гессе $\nabla^{2}g$?
  \item$^{\star}$ Прочитайте о теореме Канторовича в учебнике
        \href{https://en.wikipedia.org/wiki/Kantorovich_theorem}{«Kantorovich theorem»}
        и сформулируйте её в одну фразу: какие именно три величины
        должны быть «в нужном соотношении» для гарантии сходимости?
\end{enumerate}
\end{tasksbox}

%      % =====================================================================
%  Глава «Вычислительная линейная алгебра».
%  Самодостаточные сюжеты: SVD, Eigenfaces/Eigencats, PageRank,
%  оценки Липшицевости нейросети, числа с плавающей точкой.
%  Источники: nla360, optim, ds-intro.
%
%  Назначение в книге: новая глава Части II «Математические основы».
%  Если будет принято решение о вставке между Ньютоном (текущая гл. 2)
%  и Частью III, главе автоматически присвоится номер 3 (а текущая
%  глава 3 станет 4-й и т. д.). Все ссылки внутри файла — через \ref,
%  поэтому перенумерация происходит автоматически.
%
%  Подключение из main.tex (предлагаемая позиция):
%      \part{Математические основы}
%      \input{chapters/ch02_newton}
%      \input{chapters/ch_linalg}      % <— сюда
%
%  Все используемые окружения определены в structure.tex и
%  chapters/extras.tex; новых пакетов не требуется.
% =====================================================================

\chapter{Вычислительная линейная алгебра}
\label{ch:linalg}

\emph{Если метод Ньютона из главы~\ref{ch:opt} был одним
маленьким уравнением, упрямо повторяющимся, то линейная алгебра —
это язык, на котором говорят сразу о тысячах уравнений и миллионах
чисел. Этой главой мы покажем, что многие фокусы современного
искусственного интеллекта — от того, как нейросеть «узнаёт» лицо
человека, до того, как поисковик ранжирует страницы Интернета, —
держатся всего на нескольких разложениях матриц. Мы посмотрим
четыре истории: лица, котиков, Интернет и устойчивость нейросетей.}

\section{Зачем линейной алгебре отдельная глава}
\label{sec:linalg-why}

Вы уже встречались с матрицами в школьном курсе — как с прямоугольными
таблицами чисел. В современном ИИ матрица — это \term{универсальный
носитель данных}.

\begin{itemize}\itemsep1pt
  \item Чёрно-белая фотография $1024\times 1024$ — это матрица
        размера $1024\times 1024$, где каждое число — яркость пикселя
        от $0$ (чёрный) до $255$ (белый). Цветная — три такие матрицы.
  \item Текст из $N$ слов словаря — это вектор частот длины $N$
        (модель «мешка слов»), и набор из $M$ документов — матрица
        $M\times N$.
  \item Социальная сеть из $N$ пользователей — матрица $N\times N$,
        где элемент $(i,j)$ равен $1$, если $i$ подписан на $j$.
  \item Веса одного слоя нейросети — матрица $W$, через которую
        входной сигнал умножается, чтобы получить выходной.
\end{itemize}

С такими матрицами надо уметь \emph{что-то делать}: сжимать,
сравнивать, упорядочивать, обновлять. И почти всегда оказывается,
что нужная операция выражается через одно и то же фундаментальное
разложение матрицы — \term{сингулярное}. Мы изучим его и применим
к четырём задачам:

\begin{enumerate}\itemsep1pt
  \item \textbf{Узнать лицо} по фотографии — \emph{eigenfaces}
        (§\ref{sec:linalg-eigenfaces}).
  \item \textbf{Сжать пространство признаков} — \emph{eigencats} и
        теорема Эккарта–Янга (§\ref{sec:linalg-eckart},
        §\ref{sec:linalg-eigencats}).
  \item \textbf{Упорядочить весь Интернет} — \emph{PageRank}
        (§\ref{sec:linalg-pagerank}).
  \item \textbf{Понять, когда нейросеть устойчива} — \emph{оценки
        Липшицевости} (§\ref{sec:linalg-lipschitz}).
\end{enumerate}

Прежде чем перейти к сюжетам, разберёмся с одной \emph{очень} земной
проблемой: с тем, как компьютер хранит числа.

% =====================================================================
\section{Когда арифметика лжёт: числа с плавающей точкой}
\label{sec:linalg-floats}

Откройте Python и наберите:
\begin{lstlisting}[label={lst:float-paradox}]
>>> 0.1 + 0.2
0.30000000000000004
>>> 0.1 + 0.2 == 0.3
False
\end{lstlisting}

Это не баг Python. То же самое выдаст любой современный язык —
C, C++, Java, JavaScript. И этому есть фундаментальная причина:
\term{двоичная система счисления} не умеет точно представить
десятичную дробь $0{,}1$.

\subsection{Двоичные дроби и почему $0{,}1$ — это «$0{,}333\ldots$»}

В десятичной системе $1/3 = 0{,}3333\ldots$ — бесконечная дробь.
То же самое в двоичной системе случается с $0{,}1$:
\[
0{,}1_{10} \;=\; 0{,}0001\,1001\,1001\,1001\ldots_{2}
\]
Дробь периодическая. Любой реальный компьютер хранит её, оборвав
после конечного числа разрядов (обычно 52 двоичных знака после
запятой — стандарт \engterm{IEEE~754, double precision}). Поэтому
число, которое программист записывает как \texttt{0.1}, на самом
деле равно
\[
0{,}1\;\approx\;0{,}100\,000\,000\,000\,000\,005\,551\,115\ldots
\]
Когда мы складываем такое «искажённое $0{,}1$» с искажённым $0{,}2$,
ошибки усиливаются — и получается тот самый «лишний хвостик»
$4\cdot 10^{-17}$, который мы видели в листинге.

\begin{notebox}[title={Машинная точность $\varepsilon_{\text{маш}}$}]
\term{Машинной точностью} называют наименьшее положительное число
$\varepsilon$, такое что в компьютерной арифметике
$1+\varepsilon \neq 1$. Для стандарта \texttt{double} оно равно
$\varepsilon_{\text{маш}} = 2^{-52}\approx 2{,}22\cdot 10^{-16}$.

Это означает: какое бы хорошее число $x$ вы ни записали, оно
известно компьютеру лишь с относительной точностью
$\sim\varepsilon_{\text{маш}}$.
\end{notebox}

\subsection{Когда это становится опасным: катастрофическое сокращение}

Маленькая ошибка в $17$-м знаке — мелочь? Не всегда. Рассмотрим
безобидное на вид вычисление:
\[
g(x) \;=\; \dfrac{1-\cos x}{x^{2}}\qquad\text{при }x=10^{-8}.
\]
По формуле Тейлора $1-\cos x \approx x^{2}/2$, так что
$g(10^{-8})\approx 0{,}5$. Что выдаст компьютер?

\begin{lstlisting}[label={lst:cancel}]
>>> from math import cos
>>> x = 1e-8
>>> (1 - cos(x)) / x**2
0.0           # вместо 0.5!
\end{lstlisting}

Что произошло? Число $\cos(10^{-8})\approx 1-5\cdot 10^{-17}$ так
близко к единице, что компьютер просто хранит его как
\emph{ровно единицу} (поправка $5\cdot 10^{-17}$ меньше
$\varepsilon_{\text{маш}}$). Разность $1-\cos x$ обнулилась, и весь
ответ обнулился вместе с ней.

Это явление называется \term{катастрофическим сокращением}
(\engterm{catastrophic cancellation}): когда из двух близких чисел
вычитают и теряют все значащие цифры разом.

\begin{example}{Спасение через переписывание формулы}{cancel-fix}
Воспользуемся тождеством $1-\cos x = 2\sin^{2}(x/2)$:
\[
g(x) \;=\; \dfrac{2\sin^{2}(x/2)}{x^{2}}\;=\;\dfrac12\!\left(\dfrac{\sin(x/2)}{x/2}\right)^{2}.
\]
Внутри теперь нет вычитания близких чисел, и Python выдаёт
\texttt{0.49999999999999983} — погрешность лишь в последнем знаке.
\end{example}

\begin{importantbox}
\textbf{Главный практический вывод.} Если ваш численный алгоритм
выдаёт странный ответ, первая гипотеза, которую следует проверить, —
это не «бага в библиотеке», а \emph{катастрофическое сокращение} в
вашей формуле. Часто его можно убрать, переписав формулу
алгебраически эквивалентным, но численно устойчивым способом.
\end{importantbox}

\begin{historybox}
Стандарт \engterm{IEEE~754} был принят в 1985~г., в значительной
мере благодаря Уильяму Кэхэну (\engterm{William Kahan}, премия
Тьюринга 1989~г.). До этого каждый производитель процессоров делал
плавающую арифметику по-своему, и одна и та же программа на разных
машинах могла давать заметно разные ответы. Кэхэн также является
автором знаменитого \term{алгоритма компенсированного суммирования},
который позволяет складывать миллионы чисел с почти двойной
точностью.
\end{historybox}

% =====================================================================
\section{Сингулярное разложение SVD: главный фокус линейной алгебры}
\label{sec:linalg-svd}

\subsection{Геометрическая идея}

Возьмём любую матрицу $A$ размера $m\times n$ и применим её к
единичному кругу в $\R^{n}$ (то есть к множеству всех векторов $x$
с $\norm{x}=1$, где $\norm{\cdot}$ — обычная евклидова длина).
Что получится в $\R^{m}$?

\begin{theorem}{Сингулярное разложение (SVD)}{svd}
Для любой вещественной матрицы $A$ размера $m\times n$ существует
разложение
\begin{equation}\label{eq:svd}
\boxed{\;A \;=\; U\,\Sigma\,V^{\!\top}\;}
\end{equation}
где
\begin{itemize}\itemsep1pt
  \item $U$ — \emph{ортогональная} матрица $m\times m$ (поворот в $\R^{m}$);
  \item $V$ — ортогональная матрица $n\times n$ (поворот в $\R^{n}$);
  \item $\Sigma$ — «диагональная» матрица $m\times n$ с неотрицательными
        числами $\sigma_1\geq \sigma_2\geq\ldots\geq \sigma_r > 0$ на
        диагонали (и нулями ниже), где $r=\mathrm{rank}\,A$.
\end{itemize}
Числа $\sigma_i$ называются \term{сингулярными числами} матрицы~$A$.
\end{theorem}

Содержательно \eqref{eq:svd} говорит: \emph{любая} линейная операция —
это три простых действия подряд:
\[
\underbrace{V^{\!\top}}_{\text{поворот в исходном пространстве}}
\;\to\;
\underbrace{\Sigma}_{\text{растяжение по координатным осям}}
\;\to\;
\underbrace{U}_{\text{поворот в новом пространстве}}.
\]

Геометрически это значит: образ единичной сферы под действием $A$ —
это \term{эллипсоид}, у которого полуоси равны $\sigma_1,\sigma_2,\ldots$
(см.~рис.~\ref{fig:svd-ellipse}).

\begin{figure}[ht]
\centering
\begin{tikzpicture}[scale=1.05, line cap=round]
  % Левая картинка: единичный круг
  \begin{scope}[xshift=-3.6cm]
    \draw[->, gray!60] (-1.6,0) -- (1.6,0);
    \draw[->, gray!60] (0,-1.6) -- (0,1.6);
    \draw[thick, accent] (0,0) circle (1);
    \draw[->, very thick, accent2] (0,0) -- (1,0) node[below right] {$v_1$};
    \draw[->, very thick, accent2] (0,0) -- (0,1) node[above left] {$v_2$};
    \node[below] at (0,-1.85) {$\R^{n}$: $\norm{x}=1$};
  \end{scope}
  % Стрелка A
  \draw[->, very thick, accent] (-1.5,0) -- (1.5,0)
        node[midway, above, font=\large] {$A$};
  % Правая картинка: эллипс
  \begin{scope}[xshift=3.6cm, rotate=20]
    \draw[->, gray!60, rotate=-20] (-2.4,0) -- (2.4,0);
    \draw[->, gray!60, rotate=-20] (0,-1.4) -- (0,1.4);
    \draw[thick, accent] (0,0) ellipse (2 and 0.7);
    \draw[->, very thick, accent2] (0,0) -- (2,0) node[above] {$\sigma_1 u_1$};
    \draw[->, very thick, accent2] (0,0) -- (0,0.7) node[above left] {$\sigma_2 u_2$};
  \end{scope}
  \node[below] at (3.6,-1.85) {$\R^{m}$: $\{Ax:\norm{x}=1\}$};
\end{tikzpicture}
\caption{SVD геометрически. Матрица $A$ переводит единичный круг
в эллипс, полуоси которого равны сингулярным числам $\sigma_1,\sigma_2$,
а направления полуосей — столбцы матрицы $U$ (левые сингулярные
векторы). Векторы $v_1,v_2$ (столбцы $V$, правые сингулярные
векторы) — это те направления в исходном пространстве, которые
$A$ переводит в полуоси эллипса.}
\label{fig:svd-ellipse}
\end{figure}

\subsection{Связь с собственными значениями}

Сингулярные числа — это, по сути, переодетые собственные значения.
А именно:
\begin{equation}\label{eq:svd-eig}
A^{\!\top} A \;=\; (U\Sigma V^{\!\top})^{\!\top}(U\Sigma V^{\!\top})
\;=\; V\,\Sigma^{\!\top}\Sigma\,V^{\!\top}.
\end{equation}
Матрица $\Sigma^{\!\top}\Sigma$ — диагональная, на её диагонали
стоят $\sigma_1^{2},\sigma_2^{2},\ldots$. Значит, $\sigma_i^{2}$ —
это \emph{собственные значения} симметричной матрицы $A^{\!\top}A$,
а столбцы $V$ — её собственные векторы. Эта формула важна нам
по двум причинам: (1) она обосновывает существование SVD (через
известную спектральную теорему для симметричных матриц),
(2) она будет рабочим инструментом в § \ref{sec:linalg-eigenfaces},
когда мы будем вычислять собственные лица.

\begin{historybox}
Сингулярное разложение независимо открыли итальянский математик
\emph{Эудженио Бельтрами} (1873~г.) и француз \emph{Камилл Жордан}
(1874~г.). В современном виде, как разложение прямоугольной матрицы,
его сформулировал \emph{Карл Эккарт} с \emph{Гейлом Янгом} в 1936~г.
Численно надёжный алгоритм вычисления SVD появился лишь в 1965~г.
(Голуб и Кахан) — и стал одним из ключевых алгоритмов XX века. Сегодня
он реализован в любой математической библиотеке: \texttt{numpy.linalg.svd},
\texttt{scipy.linalg.svd}, \texttt{torch.linalg.svd}.
\end{historybox}

% =====================================================================
\section{Лучшее малоранговое приближение: теорема Эккарта--Янга}
\label{sec:linalg-eckart}

Самое замечательное свойство SVD состоит в том, что оно позволяет
\emph{ужимать} матрицу, выбрасывая всё «несущественное».

\subsection{Разложение по слагаемым ранга один}

Перепишем~\eqref{eq:svd} по-другому. Обозначим столбцы матрицы
$U$ через $u_1,\ldots,u_m$ (это \emph{левые сингулярные векторы}),
а столбцы $V$ — через $v_1,\ldots,v_n$ (\emph{правые сингулярные
векторы}). Тогда $A$ есть сумма $r$ слагаемых, каждое из которых
имеет ранг один:
\begin{equation}\label{eq:svd-sum}
A \;=\; \sum_{i=1}^{r}\sigma_i\,u_i v_i^{\!\top}.
\end{equation}

Каждое слагаемое $u_i v_i^{\!\top}$ — это «вертикальный вектор
умножить на горизонтальный», и оно само по себе является матрицей
ранга один (рис.~\ref{fig:rank1}).

\begin{figure}[ht]
\centering
\begin{tikzpicture}[scale=0.85]
  % колонка u
  \draw (0,0) rectangle (0.5,2.5);
  \node at (0.25,2.8) {$u$};
  \node at (0.25,-0.4) {\scriptsize $m\times 1$};
  \node at (1.2,1.25) {\Large $\times$};
  % строка v^T
  \draw (1.7,1.0) rectangle (4.2,1.5);
  \node at (2.95,1.7) {$v^{\!\top}$};
  \node at (2.95,0.5) {\scriptsize $1\times n$};
  \node at (4.8,1.25) {\Large $=$};
  % результат
  \draw[fill=accent!10] (5.3,0) rectangle (7.8,2.5);
  \node at (6.55,2.8) {$u v^{\!\top}$};
  \node at (6.55,-0.4) {\scriptsize $m\times n$, ранг $1$};
\end{tikzpicture}
\caption{Матрица ранга один: каждый столбец результата
пропорционален $u$, каждая строка — пропорциональна $v^{\!\top}$.
Чтобы её хранить, достаточно $m+n$ чисел вместо $mn$.}
\label{fig:rank1}
\end{figure}

Если оборвать сумму~\eqref{eq:svd-sum} на $k$-м слагаемом ($k<r$),
получим приближение:
\begin{equation}\label{eq:lowrank}
A_k \;=\; \sum_{i=1}^{k}\sigma_i\,u_i v_i^{\!\top}.
\end{equation}
Это матрица того же размера $m\times n$, но ранга всего лишь $k$.

\subsection{Теорема Эккарта--Янга: $A_k$ — наилучшее приближение}

\begin{theorem}{Эккарт--Янг--Мирский (1936/1960)}{eckart-young}
Среди всех матриц $B$ размера $m\times n$ ранга не более $k$
матрица $A_k$ из~\eqref{eq:lowrank} даёт \emph{наименьшую} ошибку
в смысле спектральной нормы:
\begin{equation}\label{eq:ey}
\min_{\mathrm{rank}\,B\,\leq\,k}\;\norm{A-B}_{2}\;=\;\norm{A-A_k}_{2}\;=\;\sigma_{k+1}.
\end{equation}
Здесь $\norm{M}_2 = \max_{\norm{x}=1}\norm{Mx}=\sigma_1(M)$ —
\term{спектральная норма} матрицы $M$ (равна её наибольшему
сингулярному числу). Утверждение остаётся в силе и для
\emph{фробениусовой} нормы $\norm{M}_F=\sqrt{\sum_{i,j} m_{ij}^{2}}$:
\[
\norm{A-A_k}_F\;=\;\sqrt{\sigma_{k+1}^{2}+\sigma_{k+2}^{2}+\ldots+\sigma_r^{2}}.
\]
\end{theorem}

Смысл теоремы: \emph{если хочется сжать матрицу до ранга $k$, лучшее,
что вы можете сделать, — оставить первые $k$ компонент SVD}.
Ошибка контролируется \emph{отброшенным} сингулярным числом
$\sigma_{k+1}$.

\subsection{Пример: сжатие изображения}

Возьмём фотографию $512\times 512$ в градациях серого. Хранение
«в лоб» — это $512\cdot 512 = 262\,144$ числа. Применим к матрице
фото SVD и оставим только $k$ старших компонент: придётся хранить
\[
k\cdot(512+512)+k\;=\;k\cdot 1025
\]
чисел — по одному столбцу $u_i$, одной строке $v_i^{\!\top}$ и
одному числу $\sigma_i$ на каждое слагаемое. Уже при $k=50$ это
$\sim 51\,250$ чисел — в \emph{пять} раз меньше исходного,
а визуально потеря почти незаметна (рис.~\ref{fig:img-compress}).

\begin{lstlisting}[label={lst:svd-compress}]
import numpy as np
import matplotlib.pyplot as plt
from skimage import data, color

# Load and convert to grayscale
img = color.rgb2gray(data.astronaut())        # shape (512, 512)
U, S, Vt = np.linalg.svd(img, full_matrices=False)

# Reconstruct using only top-k singular components
ranks = [1, 5, 20, 50, 200]
fig, axes = plt.subplots(1, len(ranks)+1, figsize=(14, 3))
axes[0].imshow(img, cmap='gray'); axes[0].set_title('original')
for ax, k in zip(axes[1:], ranks):
    A_k = U[:, :k] @ np.diag(S[:k]) @ Vt[:k, :]
    ax.imshow(A_k, cmap='gray')
    ax.set_title(f'rank {k}')
for ax in axes: ax.axis('off')
plt.tight_layout(); plt.show()
\end{lstlisting}

\begin{figure}[ht]
\centering
\includegraphics[width=\linewidth]{svd_image_compression.pdf}
\caption{Малоранговое приближение портрета астронавта. Уже при
ранге $50$ изображение визуально неотличимо от оригинала; при ранге
$200$ восстановление идеально. По теореме Эккарта--Янга это
\emph{лучшее} приближение фиксированного ранга.}
\label{fig:img-compress}
\end{figure}

\begin{notebox}[title={LoRA: малоранговая адаптация больших моделей}]
В 2021~г. Эдвард Ху с соавторами заметили: когда большую языковую
модель дообучают на новую задачу, \emph{обновление} матриц весов
$\Delta W$ оказывается почти малоранговой матрицей. Поэтому вместо
обучения «полной» $\Delta W$ можно сразу искать её в виде
$\Delta W = U V^{\!\top}$ с маленьким $k$, где $U,V$ —
тонкие прямоугольники.

Это метод \term{LoRA} (\engterm{Low-Rank Adaptation}). Для модели с
175 миллиардами параметров такой приём сокращает объём
дообучаемых весов в тысячи раз — и именно благодаря ему сегодня
можно «дофайнтюнить» большую модель на бытовой видеокарте.

Теоретический фундамент LoRA — та самая теорема
Эккарта--Янга~\ref{thm:eckart-young}, которой почти 90 лет.
\end{notebox}

% =====================================================================
\section{Eigenfaces: как алгоритм узнаёт лицо}
\label{sec:linalg-eigenfaces}

\subsection{Постановка задачи}

В $1991$~г. Мэттью Тёрк и Алекс Пентланд из MIT поставили простой
вопрос: \emph{можно ли обучить компьютер узнавать лицо, не прибегая
ни к каким специальным «лицевым» признакам — носу, глазам, контурам, —
а просто работая с матрицей пикселей?} Их ответ — метод
\term{собственных лиц} (\engterm{eigenfaces}) — стал классикой и
открыл целую область computer vision.

Идея: каждая фотография лица — это точка в очень многомерном
пространстве. Например, лицо $64\times 64$ — это вектор длины
$64\cdot 64 = 4\,096$ в пространстве $\R^{4096}$. Двух разных людей
снимали много раз — получаем облако точек в этом пространстве.

Главное наблюдение: \emph{это облако очень узкое}. Хотя пространство
$4\,096$-мерное, реальные лица занимают в нём подпространство
размерности всего лишь $\sim 50$. Найти это подпространство — и
есть задача.

\begin{figure}[ht]
\centering
\includegraphics[width=0.85\linewidth]{eigenfaces_dataset.pdf}
\caption{Фрагмент датасета \engterm{Olivetti Faces}: $400$ фотографий
$64\times 64$ пикселя, $40$ человек по $10$ снимков каждого. Каждая
фотография --- точка в $\R^{4096}$.}
\label{fig:olivetti}
\end{figure}

\subsection{Шаг 1. Среднее лицо и центрирование}

Пусть у нас $N$ фотографий $x_1,\ldots,x_N\in\R^{d}$ (где $d=4096$ для
$64\times 64$). Вычислим \term{среднее лицо}:
\[
\bar x \;=\; \dfrac{1}{N}\sum_{i=1}^{N} x_i.
\]
И вычтем его из каждого изображения: $\tilde x_i = x_i - \bar x$.
\emph{Среднее лицо} (рис.~\ref{fig:mean-face}) — это размытая
маска, в которой видно «лицо вообще»: усреднённые глаза, нос,
рот, овал. Дальше мы будем работать только с отклонениями
от этого среднего.

\begin{figure}[ht]
\centering
\includegraphics[width=0.55\linewidth]{eigenfaces_mean.pdf}
\caption{Среднее лицо по $400$ снимкам датасета Olivetti.}
\label{fig:mean-face}
\end{figure}

\subsection{Шаг 2. Главные направления изменчивости}

Соберём центрированные изображения в строки матрицы
$X\in\R^{N\times d}$ (одна строка = одно лицо). Применим к ней SVD:
\[
X\;=\;U\,\Sigma\,V^{\!\top}.
\]
Каждая \emph{строка} $v_i^{\!\top}$ матрицы $V^{\!\top}$ — это вектор
длины $d=4096$, который можно нарисовать как картинку $64\times 64$.
Эти картинки и называются \term{собственными лицами} — eigenfaces
(рис.~\ref{fig:eigenfaces}). Содержательно:
\begin{itemize}\itemsep1pt
  \item \emph{$v_1$} — направление наибольшей изменчивости.
        Чаще всего это «освещение слева/справа».
  \item \emph{$v_2$} — направление второй по величине изменчивости
        (часто «улыбка / не улыбка»).
  \item \emph{$v_3,\ldots$} — постепенно более тонкие признаки: очки,
        форма щёк, наклон головы, индивидуальные особенности.
\end{itemize}

\begin{figure}[ht]
\centering
\includegraphics[width=\linewidth]{eigenfaces_components.pdf}
\caption{Первые 12 eigenfaces датасета Olivetti, отрисованные как
картинки $64\times 64$. Эти странные «полу-лица» — координатные
оси в пространстве, где живут все настоящие лица.}
\label{fig:eigenfaces}
\end{figure}

\subsection{Шаг 3. Лицо как вектор из 50 чисел}

Зафиксируем $k$ (обычно $k=30\div 100$) и спроектируем каждое лицо
$x$ на эти $k$ собственных лиц:
\begin{equation}\label{eq:eigenface-proj}
\alpha_i \;=\; v_i^{\!\top}(x-\bar x),\qquad i=1,\ldots,k.
\end{equation}
Получили вектор коэффициентов $\alpha=(\alpha_1,\ldots,\alpha_k)\in\R^{k}$.
Это и есть «лицо человека» с точки зрения алгоритма — \emph{всего
$50$ чисел} вместо $4096$.

\subsection{Шаг 4. Распознавание и реконструкция}

\paragraph{Распознавание.} Чтобы узнать, кто на новой фотографии $y$,
вычислим её коэффициенты $\beta$ по~\eqref{eq:eigenface-proj} и
найдём в базе ближайшего соседа: $\hat j = \argmin_j \norm{\beta-\alpha^{(j)}}$.

\paragraph{Реконструкция.} Зная $\alpha$, можно восстановить лицо:
\[
x \;\approx\; \bar x + \sum_{i=1}^{k}\alpha_i\,v_i.
\]
Рис.~\ref{fig:eigenfaces-recon} показывает, как качество
реконструкции растёт с $k$.

\begin{figure}[ht]
\centering
\includegraphics[width=\linewidth]{eigenfaces_reconstruction.pdf}
\caption{Реконструкция одной из фотографий датасета через первые $k$
eigenfaces. При $k=10$ узнаётся «общий тип» лица; при $k=50$ — уже
конкретный человек; при $k=150$ — почти оригинал.}
\label{fig:eigenfaces-recon}
\end{figure}

\subsection{Реализация на Python}

\begin{lstlisting}[label={lst:eigenfaces}]
import numpy as np
import matplotlib.pyplot as plt
from sklearn.datasets import fetch_olivetti_faces

# 1. Load: 400 faces, 64x64 each
faces = fetch_olivetti_faces().images        # (400, 64, 64)
X = faces.reshape(400, -1)                   # (400, 4096)

# 2. Center: subtract the mean face
x_bar = X.mean(axis=0)
X_centered = X - x_bar

# 3. SVD => eigenfaces are rows of Vt
U, S, Vt = np.linalg.svd(X_centered, full_matrices=False)

# 4. Project all faces onto top-k eigenfaces
k = 50
alpha = X_centered @ Vt[:k].T                # (400, 50)

# 5. Recognise: nearest neighbour in alpha-space
def recognise(y, k=50):
    y_centered = y.flatten() - x_bar
    beta = Vt[:k] @ y_centered
    distances = np.linalg.norm(alpha - beta, axis=1)
    return distances.argmin()

# 6. Reconstruct a face from k coefficients
def reconstruct(idx, k=50):
    return x_bar + Vt[:k].T @ alpha[idx, :k]
\end{lstlisting}

\begin{interestbox}
Метод eigenfaces — родоначальник всех современных систем распознавания
лиц, включая ту, что разблокирует ваш смартфон. Современные системы
заменяют SVD на свёрточную нейросеть (см.~§\,3.3.7), но базовая
идея — \emph{сжать лицо в короткий вектор, в котором близкие лица
оказываются рядом} — осталась той же. Эти короткие векторы сегодня
называют \term{эмбеддингами}.
\end{interestbox}

% =====================================================================
\section{Eigencats: то же самое для котиков\texorpdfstring{$^\star$}{*}}
\label{sec:linalg-eigencats}

Кошки — отличный материал для проверки гипотезы, что метод eigenfaces
работает на любых однородных объектах, не только на лицах. Возьмём
датасет в $\sim 2000$ котиков\footnote{Например,
\engterm{Cats Faces Dataset} на Kaggle: фотографии $64\times 64$.
Готовые ноутбуки выложены в репозитории \href{https://github.com/MerkulovDaniil/nla360}{nla360}.}
и проведём тот же эксперимент.

\begin{figure}[ht]
\centering
\includegraphics[width=0.85\linewidth]{eigencats_components.pdf}
\caption{Первые 8 eigencats. Видны общие «коты»: положение ушей,
форма морды, направление взгляда. Каждый конкретный кот собирается
как линейная комбинация этих универсальных «прото-котов».}
\label{fig:eigencats}
\end{figure}

Качественная картина повторяется: первое eigencat ловит освещение,
второе — наклон головы, дальше идут вариации формы морды и ушей.
В сжатом виде каждый кот превращается в $\sim 50$ чисел.

\begin{notebox}[title={Что это нам говорит о ИИ}]
\textbf{Пространство «реальных» объектов всегда узкое.} Будь это
лица людей, морды котов, рукописные цифры (датасет MNIST) или
фотографии природы — реальные данные занимают в исходном
многомерном пространстве \emph{маленькое} подпространство.
Распознавание = найти это подпространство и работать в нём.

\textbf{Эта идея — фундамент современных автокодировщиков и
эмбеддингов} (см.~§\,3.4.3--3.4.6). SVD — её самая ранняя и самая
понятная реализация.
\end{notebox}

% =====================================================================
\section{PageRank: как Google нашёл порядок в Интернете}
\label{sec:linalg-pagerank}

\subsection{Задача и история}

К 1996~г. поисковики типа \engterm{AltaVista} и \engterm{Yahoo!}
ранжировали страницы по тому, сколько раз искомое слово встречается
в тексте. Это легко обманывалось: достаточно было набить страницу
популярными словами белым цветом по белому фону, и она попадала
в топ выдачи.

Лоуренс Пейдж и Сергей Брин, тогда аспиранты Стэнфорда, предложили
другой принцип ранжирования: \emph{страница важна, если на неё
ссылаются другие важные страницы}.

\subsection{Граф ссылок как матрица}

Обозначим $N$ — число всех страниц Интернета (на 1998~г. — около
$2{,}4\cdot 10^{8}$, сегодня — $\sim 10^{11}$). Построим матрицу
$P\in\R^{N\times N}$:
\[
P_{ij}\;=\;\begin{cases}
\dfrac{1}{c_j}, & \text{если со страницы $j$ есть ссылка на $i$,}\\
0, & \text{иначе,}
\end{cases}
\]
где $c_j$ — \emph{полное число исходящих ссылок} со страницы $j$.
Содержательно: $P_{ij}$ — это вероятность того, что пользователь,
сидящий на странице $j$ и случайно ткнувший на одну из её ссылок,
окажется на странице $i$. Поэтому $P$ называется матрицей переходов
\term{случайного блуждания} по графу Интернета.

\subsection{Стационарное распределение = собственный вектор}

Пусть $r\in\R^{N}$ — вектор «важностей»: $r_i$ — степень важности
страницы $i$. Принцип Пейджа–Брина формализуется системой:
\begin{equation}\label{eq:pr-fixed}
r\;=\;P r.
\end{equation}
То есть важность каждой страницы равна сумме важностей тех страниц,
которые на неё ссылаются, делённой на число их исходящих ссылок.

Уравнение~\eqref{eq:pr-fixed} говорит: $r$ — это \emph{собственный
вектор} матрицы $P$ с собственным значением $1$. И в этом — вся суть
PageRank.

\begin{theorem}{Перрон--Фробениус для PageRank}{perron}
Если все элементы матрицы $P$ положительны (или хотя бы граф ссылок
сильно связен и непериодичен), то собственное значение $1$ у $P$
существует, оно \emph{простое} и наибольшее по модулю, а
соответствующий собственный вектор $r$ единственен и имеет
положительные компоненты. Этот $r$ называется \term{стационарным
распределением} случайного блуждания.
\end{theorem}

(Эта теорема — Оскар Перрон, 1907~г., и Георг Фробениус, 1912~г. —
исторически появилась задолго до Интернета, в контексте моделей
популяций.)

\subsection{Демпфирующий множитель и формула Брина--Пейджа}

В реальной матрице Интернета у некоторых страниц нет исходящих ссылок
(\engterm{dangling pages}), а в других группах — циклы. Для таких
страниц соответствующий столбец матрицы переходов заранее заменяют
равномерным распределением по всем страницам: иначе формула $1/c_j$
при $c_j=0$ не определена, а масса случайного блуждания «исчезала» бы
в dangling pages. Чтобы
гарантировать применимость теоремы Перрона--Фробениуса, Пейдж и Брин
ввели \term{демпфирующий множитель} $\beta\in(0,1)$ (обычно
$\beta=0{,}85$):
\begin{equation}\label{eq:pagerank}
\boxed{\;r\;=\;\beta P r + \dfrac{1-\beta}{N}\,\mathbf{1}.\;}
\end{equation}
Здесь $\mathbf{1}=(1,1,\ldots,1)^{\!\top}$. Содержательно:
с вероятностью $\beta\approx 0{,}85$ случайный сёрфер кликает по
ссылке, а с вероятностью $1-\beta\approx 0{,}15$ — телепортируется
на случайную страницу Интернета.

\subsection{Как находить $r$: степенной метод}

Прямо решать систему $(I-\beta P)r=\frac{1-\beta}{N}\mathbf{1}$
с $N\sim 10^{11}$ невозможно. Но и не надо: вектор $r$ можно
получить простой итерацией.

\begin{algorithmbox}
\textbf{Степенной метод для PageRank.}
\begin{enumerate}\itemsep1pt
  \item Инициализация: $r^{(0)} \leftarrow \tfrac{1}{N}\mathbf{1}$.
  \item Итерация: $r^{(t+1)} \leftarrow \beta\,P\,r^{(t)} +
        \dfrac{1-\beta}{N}\,\mathbf{1}$.
  \item Остановка: $\norm{r^{(t+1)}-r^{(t)}} < \varepsilon$.
\end{enumerate}
\end{algorithmbox}

Сходимость гарантируется теоремой Перрона--Фробениуса:
$\norm{r^{(t)}-r}\leq C\cdot\beta^{t}$, то есть погрешность убывает
геометрически со знаменателем $\beta=0{,}85$. На практике хватает
$50\div 100$ итераций.

\subsection{Игрушечный пример: 6 страниц}

Возьмём маленький «Интернет» из шести страниц (рис.~\ref{fig:pagerank-graph})
и посчитаем для него PageRank.

\begin{figure}[ht]
\centering
\begin{tikzpicture}[scale=1.0, every node/.style={circle, draw, accent,
              fill=accent!10, minimum size=8mm, font=\small}]
  \node (A) at (0,2)    {A};
  \node (B) at (2.5,3)  {B};
  \node (C) at (5,2)    {C};
  \node (D) at (5,-0.5) {D};
  \node (E) at (2.5,-1.5){E};
  \node (F) at (0,-0.5) {F};
  \begin{scope}[->, thick, accent2]
    \draw (A) -- (B);
    \draw (A) -- (C);
    \draw (B) -- (C);
    \draw (C) -- (A);
    \draw (D) -- (C);
    \draw (D) -- (A);
    \draw (E) -- (D);
    \draw (E) -- (B);
    \draw (E) -- (F);
    \draw (F) -- (E);
    \draw (F) -- (B);
  \end{scope}
\end{tikzpicture}
\caption{Игрушечный «Интернет» из 6 страниц со ссылками. Видно,
например, что на страницу C ссылаются три другие, а на F — только
одна; будет ли C лидером по PageRank?}
\label{fig:pagerank-graph}
\end{figure}

\begin{lstlisting}[label={lst:pagerank}]
import numpy as np

# Adjacency list -> column-stochastic transition matrix P
links = {'A': ['B','C'], 'B': ['C'], 'C': ['A'],
         'D': ['C','A'], 'E': ['D','B','F'], 'F': ['E','B']}
pages = list(links); n = len(pages); idx = {p: i for i, p in enumerate(pages)}

P = np.zeros((n, n))
for j, p in enumerate(pages):
    out = links[p]
    for q in out:
        P[idx[q], j] = 1.0 / len(out)

# PageRank via power iteration
beta, eps = 0.85, 1e-10
r = np.ones(n) / n
for t in range(200):
    r_new = beta * P @ r + (1 - beta) / n
    if np.linalg.norm(r_new - r, 1) < eps: break
    r = r_new

for p, ri in sorted(zip(pages, r), key=lambda t: -t[1]):
    print(f'{p}: {ri:.4f}')
\end{lstlisting}

\noindent Выдача программы:
\[
\text{C: }0{,}295,\quad
\text{A: }0{,}252,\quad
\text{B: }0{,}175,\quad
\text{E: }0{,}098,\quad
\text{F: }0{,}092,\quad
\text{D: }0{,}087.
\]
Лидер — \texttt{C}, хотя на неё формально ссылаются столько же
страниц, сколько на \texttt{A}: PageRank учёл, что одна из ссылок
на \texttt{C} приходит с уже-важной страницы \texttt{B}.

\begin{historybox}
Алгоритм PageRank был опубликован в 1998~г. в студенческой работе
\emph{Sergey Brin, Larry Page}, «\textit{The PageRank Citation Ranking:
Bringing Order to the Web}». В том же году Брин и Пейдж основали
Google. К 2004~г. компания вышла на IPO, а к 2010-м стала одной из
крупнейших в мире. Сам Стэнфорд получил $\sim 336$ миллионов долларов
от лицензирования патента --- который изначально просто описывал
итерационный метод для собственного вектора.

Идея использовать собственный вектор графа для измерения «важности»
позднее была обобщена и применена в наукометрии (\engterm{Eigenfactor}
для ранжирования научных журналов), в рекомендательных системах,
а также в современных нейросетях для графов (\engterm{Graph Neural
Networks}).
\end{historybox}

% =====================================================================
\section{Когда нейросеть устойчива: оценки Липшицевости\texorpdfstring{$^\star$}{*}}
\label{sec:linalg-lipschitz}

\subsection{Сюжет: когда панда становится гиббоном}

В 2014~г. Иэн Гудфеллоу и соавторы опубликовали ставшую знаменитой
картинку (рис.~\ref{fig:adv-attack}): обычная фотография панды,
которую современная нейросеть уверенно опознаёт как «panda» с
вероятностью $57{,}7\%$. К каждому пикселю прибавлена крошечная,
визуально незаметная человеку поправка. На обновлённой картинке
сеть теперь уверенно (с вероятностью $99{,}3\%$) считает, что
видит \emph{гиббона}.

\begin{figure}[ht]
\centering
\includegraphics[width=0.95\linewidth]{adversarial_panda.pdf}
\caption{Adversarial-атака на нейросеть (Goodfellow et al., 2014).
Слева — оригинал, классифицирован как «panda» с уверенностью $57{,}7\%$.
В центре — почти невидимое возмущение, масштабированное для наглядности.
Справа — атакованный пример, классифицируется как «gibbon» с
уверенностью $99{,}3\%$.}
\label{fig:adv-attack}
\end{figure}

Этот эффект называется \term{adversarial-атакой} и долгое время
считался экзотикой. Сейчас понятно: причина в простом числе,
которое сопровождает любую функцию — её \emph{константе Липшица}.

\subsection{Константа Липшица: формальное определение}

\begin{definition}{Константа Липшица функции}{lipschitz}
Функция $f\colon\R^{n}\to\R^{m}$ называется \term{Липшицевой} с
константой $L>0$, если для любых $x,y\in\R^{n}$ выполнено
\begin{equation}\label{eq:lipschitz}
\norm{f(x)-f(y)}\;\leq\;L\,\norm{x-y}.
\end{equation}
Наименьшее $L$, для которого~\eqref{eq:lipschitz} выполнено,
называется \emph{константой Липшица} функции $f$ и обозначается
$\mathrm{Lip}(f)$.
\end{definition}

Содержательно $\mathrm{Lip}(f)$ показывает, во сколько раз может
\emph{усилиться} входное возмущение на выходе функции. Если
$\mathrm{Lip}(f)=L$, то малое возмущение $\delta$ на входе
может вызвать возмущение до $L\cdot\norm{\delta}$ на выходе.

\paragraph{Пример.} Для линейного отображения $f(x)=Wx$ имеем
$f(x)-f(y)=W(x-y)$, откуда $\norm{f(x)-f(y)}\leq \norm{W}_2\cdot\norm{x-y}$,
и константа Липшица — это в точности спектральная норма матрицы $W$:
\begin{equation}\label{eq:lip-linear}
\mathrm{Lip}(Wx)\;=\;\norm{W}_2\;=\;\sigma_1(W).
\end{equation}
Здесь $\sigma_1(W)$ — наибольшее сингулярное число матрицы $W$.
Вот \emph{здесь} в нашу историю входит SVD из §\,\ref{sec:linalg-svd}.

\subsection{Композиция: Липшицева константа цепочки}

\begin{theorem}{Композиция Липшицевых функций}{lip-compose}
Если $f$ и $g$ — Липшицевы с константами $L_f$ и $L_g$, то
их композиция $f\circ g$ — тоже Липшицева, причём
\begin{equation}\label{eq:lip-prod}
\mathrm{Lip}(f\circ g)\;\leq\;\mathrm{Lip}(f)\cdot\mathrm{Lip}(g).
\end{equation}
\end{theorem}

\begin{proof}
$\norm{f(g(x))-f(g(y))}\leq L_f\norm{g(x)-g(y)}\leq L_f L_g\norm{x-y}$. $\square$
\end{proof}

\subsection{Применение к нейросети}

Возьмём типичную нейросеть глубины $K$:
\[
f(x)\;=\;\sigma\bigl(W_K\,\sigma(W_{K-1}\,\sigma(\cdots\sigma(W_1 x))\bigr),
\]
где $W_i$ — матрицы весов, $\sigma$ — поэлементная активация (ReLU,
сигмоида, и т.п.). Для большинства активаций $\mathrm{Lip}(\sigma)\leq 1$
(для ReLU и $\tanh$ она равна $1$, а для логистической сигмоиды — $1/4$). По
теореме~\ref{thm:lip-compose}:
\begin{equation}\label{eq:nn-lip}
\boxed{\;\mathrm{Lip}(f)\;\leq\;\prod_{i=1}^{K}\norm{W_i}_2\;}
\end{equation}

Это и есть тот «лишний рычаг», который объясняет уязвимость
нейросетей. Если $\norm{W_i}_2 \approx 2$ для каждого слоя из $K=20$
слоёв, то возмущение на входе \emph{потенциально} усиливается
в $2^{20}\approx 10^{6}$ раз. Невидимая поправка $\delta$ с
$\norm{\delta}=10^{-3}$ может вызвать изменение выхода до $10^{3}$ —
этого хватит, чтобы перебросить классификатор из одного класса в другой.

\subsection{Как защищаться: спектральная нормализация}

Естественная стратегия: после каждого шага обучения нормализовать
матрицу $W_i$ так, чтобы $\norm{W_i}_2 = 1$. Это и есть метод
\term{Spectral Normalization} (Миято и соавт., 2018), сделавший
обучение генеративных нейросетей (GAN) намного стабильнее.

Технически: на каждом шаге считаем $\sigma_1(W)$ (с помощью одной
итерации степенного метода — почти бесплатно!) и заменяем
$W \mapsto W/\sigma_1(W)$.

\begin{lstlisting}[label={lst:spec-norm}]
def spectral_normalize(W, n_iter=1):
    """One step of power iteration to estimate sigma_1(W) and rescale."""
    u = np.random.randn(W.shape[0])
    for _ in range(n_iter):
        v = W.T @ u; v /= np.linalg.norm(v)
        u = W @ v;   u /= np.linalg.norm(u)
    sigma_1 = u @ W @ v
    return W / sigma_1
\end{lstlisting}

\begin{interestbox}
\textbf{Связь с устойчивостью в физике.} Похожая ситуация хорошо
известна в теории динамических систем: если константа Липшица
отображения $> 1$, малое возмущение начального условия экспоненциально
нарастает (классический эффект бабочки в погодных моделях).
Spectral Normalization можно рассматривать как искусственное
гашение этого эффекта.
\end{interestbox}

% =====================================================================
\section*{Подведение итогов главы}
\addcontentsline{toc}{section}{Подведение итогов главы}

\begin{itemize}\itemsep1pt
  \item Матрица — это \emph{универсальный носитель данных} для ИИ:
        картинки, тексты, графы, веса нейросети.
  \item \textbf{Числа с плавающей точкой} (§\ref{sec:linalg-floats}) —
        вычисления на компьютере всегда приближённые с относительной
        ошибкой $\sim\varepsilon_{\text{маш}}\approx 10^{-16}$.
        Самая опасная ошибка — \emph{катастрофическое сокращение} при
        вычитании близких чисел; её часто можно устранить, переписав
        формулу алгебраически эквивалентным способом.
  \item \textbf{Сингулярное разложение} $A=U\Sigma V^{\!\top}$
        (§\ref{sec:linalg-svd}) представляет любую матрицу как
        «поворот → растяжение по осям → поворот». Сингулярные числа
        $\sigma_i$ — это коэффициенты растяжения.
  \item \textbf{Теорема Эккарта--Янга} (§\ref{sec:linalg-eckart}):
        $A_k=\sum_{i\leq k}\sigma_i u_i v_i^{\!\top}$ — лучшее
        приближение ранга $k$ матрицы $A$. На этой теореме
        стоит сжатие изображений, метод eigenfaces и современный
        LoRA для дообучения больших языковых моделей.
  \item \textbf{Eigenfaces / eigencats}
        (§\ref{sec:linalg-eigenfaces}--\ref{sec:linalg-eigencats}):
        реальные данные (лица, котики, цифры) живут в очень узком
        подпространстве своего номинального многомерного пространства.
        Найти это подпространство = применить SVD к выборке.
  \item \textbf{PageRank} (§\ref{sec:linalg-pagerank}): «важность»
        страницы Интернета = компонента собственного вектора матрицы
        переходов случайного блуждания. Алгоритм находит его простой
        итерацией; теорема Перрона--Фробениуса гарантирует сходимость.
  \item \textbf{Константа Липшица} нейросети
        (§\ref{sec:linalg-lipschitz}) растёт мультипликативно от
        слоя к слою; именно поэтому глубокие сети уязвимы к
        adversarial-атакам. Спектральная нормализация ограничивает
        $\norm{W_i}_2=1$ и стабилизирует обучение.
\end{itemize}

% =====================================================================
\section*{Что почитать дальше}
\addcontentsline{toc}{section}{Что почитать дальше}

\begin{itemize}\itemsep1pt
  \item \emph{Е.\,Е.~Тыртышников.} Методы численного анализа. М.: Академия, 2007. ---
        классический российский учебник по вычислительной линейной
        алгебре университетского уровня.
  \item \emph{G.~Strang.} Linear Algebra and Learning from Data. Wellesley-Cambridge Press, 2019. ---
        свежая книга от автора знаменитого MIT-курса 18.06, написанная
        с прицелом на data science.
  \item \emph{L.\,N.~Trefethen, D.~Bau.} Numerical Linear Algebra. SIAM, 1997. ---
        классика, с очень понятным изложением SVD и итерационных методов.
  \item \emph{Видеокурс «Numerical Linear Algebra»}, И.\,В.~Оселедец, Сколтех
        (\href{https://github.com/MerkulovDaniil/nla360}{github.com/MerkulovDaniil/nla360}).
        Лекции, семинары, ноутбуки --- из которых, в частности,
        и взяты eigencats для этой главы.
  \item \emph{S.~Brin, L.~Page.} The Anatomy of a Large-Scale
        Hypertextual Web Search Engine. WWW Conference, 1998. ---
        оригинальная статья про PageRank, читается легко.
  \item \emph{I.~Goodfellow et al.} Explaining and Harnessing
        Adversarial Examples. ICLR, 2015. ---
        статья про панду-гиббона, из которой взят рис.~\ref{fig:adv-attack}.
\end{itemize}

% =====================================================================
\section*{Большой итоговый проект}
\addcontentsline{toc}{section}{Большой итоговый проект}

\textbf{Постройте свою рекомендательную систему для фильмов.}

Возьмите датасет MovieLens-100K (открытый, $100\,000$ оценок,
$\sim 1000$ пользователей $\times$ $\sim 1700$ фильмов). Постройте
матрицу «пользователь--фильм» $R$, где $R_{ij}\in\{1,\ldots,5\}$ —
оценка, поставленная пользователем $i$ фильму $j$, а $0$ — «не смотрел».

\textbf{Важно:} нули в матрице $R$ означают \emph{пропущенные} оценки,
а не настоящие нулевые рейтинги. Прямое SVD по матрице с нулями — лишь
грубый учебный baseline: непросмотренные фильмы тогда воспринимаются как
фильмы с очень низкой оценкой. Более корректный вариант — обучать
малоранговую модель только на известных оценках либо предварительно
заполнить пропуски средними значениями пользователя/фильма и затем
сравнить результат.

\begin{enumerate}\itemsep1pt
  \item Постройте малоранговое приближение $R\approx U\Sigma V^{\!\top}$
        с $k=20$ компонентами через SVD.
  \item Объясните содержательно, что такое строки $U$ (профили
        пользователей) и столбцы $V$ (профили фильмов) в этом
        приближении. Сколько чисел теперь хранит модель вместо
        $1000\cdot 1700$?
  \item Для нескольких реальных пользователей выпишите топ-5 фильмов,
        которые рекомендация модели предлагает посмотреть. Совпадают
        ли они с любимыми жанрами пользователя?
  \item Найдите два «противоположных» фильма с наибольшим расстоянием
        между их столбцами $V$ и подумайте, можно ли увидеть содержательный
        смысл в этой разделяющей оси (боевик vs. мелодрама? старое vs. новое?).
  \item ($\star$) Замените SVD на алгоритм PageRank, построив граф
        «фильмы, которые часто оценивают одинаково». Получится ли тот
        же топ?
\end{enumerate}

% =====================================================================
\begin{tasksbox}
\begin{enumerate}\itemsep1pt
  \item Сложите в Python $10\,000\,000$ копий числа $0{,}1$. Сколько
        должно получиться? Сколько получается на самом деле? Почему?
  \item Вычислите $f(x)=\sqrt{x+1}-\sqrt{x}$ при $x=10^{16}$
        напрямую и через формулу $f(x)=\frac{1}{\sqrt{x+1}+\sqrt{x}}$.
        Сравните ответы. Какой устойчив, почему?
  \item Найдите вручную SVD матрицы
        $A=\bigl(\begin{smallmatrix}3 & 0\\0 & 4\end{smallmatrix}\bigr)$
        и $A=\bigl(\begin{smallmatrix}0 & 1\\-1 & 0\end{smallmatrix}\bigr)$.
        Что в обоих случаях геометрически делают $U, \Sigma, V$?
  \item Возьмите свою фотографию $300\times 300$ в градациях серого.
        Сделайте сжатие через SVD с рангами $k=5, 20, 50, 100$. При
        каком $k$ вы перестаёте узнавать самого себя?
  \item Для игрушечного графа из рис.~\ref{fig:pagerank-graph}
        прокрутите 5 итераций степенного метода \emph{вручную},
        начав с $r^{(0)}=(\tfrac16,\ldots,\tfrac16)^{\!\top}$. Совпали ли
        ваши значения с программными после $5$ шагов? а после $20$?
  \item $^{\star}$ Покажите, что для функции $\sigma(x)=\max(0,x)$
        (ReLU) константа Липшица равна $1$. Тот же вопрос для
        сигмоиды $\sigma(x)=1/(1+e^{-x})$: чему равна $\mathrm{Lip}(\sigma)$
        и в какой точке производная максимальна?
  \item $^{\star}$ Найдите спектральную норму матрицы
        $W=\bigl(\begin{smallmatrix}1 & 2\\2 & 1\end{smallmatrix}\bigr)$
        двумя способами: (а) через определение
        $\norm{W}_2=\max_{\norm{x}=1}\norm{Wx}$ (с использованием
        собственных значений $W^{\!\top}W$); (б) одной итерацией
        степенного метода (с любым стартом). Совпали ли ответы?
\end{enumerate}
\end{tasksbox}
      % <— сюда
%
%  Все используемые окружения определены в structure.tex и
%  chapters/extras.tex; новых пакетов не требуется.
% =====================================================================

\chapter{Вычислительная линейная алгебра}
\label{ch:linalg}

\emph{Если метод Ньютона из главы~\ref{ch:opt} был одним
маленьким уравнением, упрямо повторяющимся, то линейная алгебра —
это язык, на котором говорят сразу о тысячах уравнений и миллионах
чисел. Этой главой мы покажем, что многие фокусы современного
искусственного интеллекта — от того, как нейросеть «узнаёт» лицо
человека, до того, как поисковик ранжирует страницы Интернета, —
держатся всего на нескольких разложениях матриц. Мы посмотрим
четыре истории: лица, котиков, Интернет и устойчивость нейросетей.}

\section{Зачем линейной алгебре отдельная глава}
\label{sec:linalg-why}

Вы уже встречались с матрицами в школьном курсе — как с прямоугольными
таблицами чисел. В современном ИИ матрица — это \term{универсальный
носитель данных}.

\begin{itemize}\itemsep1pt
  \item Чёрно-белая фотография $1024\times 1024$ — это матрица
        размера $1024\times 1024$, где каждое число — яркость пикселя
        от $0$ (чёрный) до $255$ (белый). Цветная — три такие матрицы.
  \item Текст из $N$ слов словаря — это вектор частот длины $N$
        (модель «мешка слов»), и набор из $M$ документов — матрица
        $M\times N$.
  \item Социальная сеть из $N$ пользователей — матрица $N\times N$,
        где элемент $(i,j)$ равен $1$, если $i$ подписан на $j$.
  \item Веса одного слоя нейросети — матрица $W$, через которую
        входной сигнал умножается, чтобы получить выходной.
\end{itemize}

С такими матрицами надо уметь \emph{что-то делать}: сжимать,
сравнивать, упорядочивать, обновлять. И почти всегда оказывается,
что нужная операция выражается через одно и то же фундаментальное
разложение матрицы — \term{сингулярное}. Мы изучим его и применим
к четырём задачам:

\begin{enumerate}\itemsep1pt
  \item \textbf{Узнать лицо} по фотографии — \emph{eigenfaces}
        (§\ref{sec:linalg-eigenfaces}).
  \item \textbf{Сжать пространство признаков} — \emph{eigencats} и
        теорема Эккарта–Янга (§\ref{sec:linalg-eckart},
        §\ref{sec:linalg-eigencats}).
  \item \textbf{Упорядочить весь Интернет} — \emph{PageRank}
        (§\ref{sec:linalg-pagerank}).
  \item \textbf{Понять, когда нейросеть устойчива} — \emph{оценки
        Липшицевости} (§\ref{sec:linalg-lipschitz}).
\end{enumerate}

Прежде чем перейти к сюжетам, разберёмся с одной \emph{очень} земной
проблемой: с тем, как компьютер хранит числа.

% =====================================================================
\section{Когда арифметика лжёт: числа с плавающей точкой}
\label{sec:linalg-floats}

Откройте Python и наберите:
\begin{lstlisting}[label={lst:float-paradox}]
>>> 0.1 + 0.2
0.30000000000000004
>>> 0.1 + 0.2 == 0.3
False
\end{lstlisting}

Это не баг Python. То же самое выдаст любой современный язык —
C, C++, Java, JavaScript. И этому есть фундаментальная причина:
\term{двоичная система счисления} не умеет точно представить
десятичную дробь $0{,}1$.

\subsection{Двоичные дроби и почему $0{,}1$ — это «$0{,}333\ldots$»}

В десятичной системе $1/3 = 0{,}3333\ldots$ — бесконечная дробь.
То же самое в двоичной системе случается с $0{,}1$:
\[
0{,}1_{10} \;=\; 0{,}0001\,1001\,1001\,1001\ldots_{2}
\]
Дробь периодическая. Любой реальный компьютер хранит её, оборвав
после конечного числа разрядов (обычно 52 двоичных знака после
запятой — стандарт \engterm{IEEE~754, double precision}). Поэтому
число, которое программист записывает как \texttt{0.1}, на самом
деле равно
\[
0{,}1\;\approx\;0{,}100\,000\,000\,000\,000\,005\,551\,115\ldots
\]
Когда мы складываем такое «искажённое $0{,}1$» с искажённым $0{,}2$,
ошибки усиливаются — и получается тот самый «лишний хвостик»
$4\cdot 10^{-17}$, который мы видели в листинге.

\begin{notebox}[title={Машинная точность $\varepsilon_{\text{маш}}$}]
\term{Машинной точностью} называют наименьшее положительное число
$\varepsilon$, такое что в компьютерной арифметике
$1+\varepsilon \neq 1$. Для стандарта \texttt{double} оно равно
$\varepsilon_{\text{маш}} = 2^{-52}\approx 2{,}22\cdot 10^{-16}$.

Это означает: какое бы хорошее число $x$ вы ни записали, оно
известно компьютеру лишь с относительной точностью
$\sim\varepsilon_{\text{маш}}$.
\end{notebox}

\subsection{Когда это становится опасным: катастрофическое сокращение}

Маленькая ошибка в $17$-м знаке — мелочь? Не всегда. Рассмотрим
безобидное на вид вычисление:
\[
g(x) \;=\; \dfrac{1-\cos x}{x^{2}}\qquad\text{при }x=10^{-8}.
\]
По формуле Тейлора $1-\cos x \approx x^{2}/2$, так что
$g(10^{-8})\approx 0{,}5$. Что выдаст компьютер?

\begin{lstlisting}[label={lst:cancel}]
>>> from math import cos
>>> x = 1e-8
>>> (1 - cos(x)) / x**2
0.0           # вместо 0.5!
\end{lstlisting}

Что произошло? Число $\cos(10^{-8})\approx 1-5\cdot 10^{-17}$ так
близко к единице, что компьютер просто хранит его как
\emph{ровно единицу} (поправка $5\cdot 10^{-17}$ меньше
$\varepsilon_{\text{маш}}$). Разность $1-\cos x$ обнулилась, и весь
ответ обнулился вместе с ней.

Это явление называется \term{катастрофическим сокращением}
(\engterm{catastrophic cancellation}): когда из двух близких чисел
вычитают и теряют все значащие цифры разом.

\begin{example}{Спасение через переписывание формулы}{cancel-fix}
Воспользуемся тождеством $1-\cos x = 2\sin^{2}(x/2)$:
\[
g(x) \;=\; \dfrac{2\sin^{2}(x/2)}{x^{2}}\;=\;\dfrac12\!\left(\dfrac{\sin(x/2)}{x/2}\right)^{2}.
\]
Внутри теперь нет вычитания близких чисел, и Python выдаёт
\texttt{0.49999999999999983} — погрешность лишь в последнем знаке.
\end{example}

\begin{importantbox}
\textbf{Главный практический вывод.} Если ваш численный алгоритм
выдаёт странный ответ, первая гипотеза, которую следует проверить, —
это не «бага в библиотеке», а \emph{катастрофическое сокращение} в
вашей формуле. Часто его можно убрать, переписав формулу
алгебраически эквивалентным, но численно устойчивым способом.
\end{importantbox}

\begin{historybox}
Стандарт \engterm{IEEE~754} был принят в 1985~г., в значительной
мере благодаря Уильяму Кэхэну (\engterm{William Kahan}, премия
Тьюринга 1989~г.). До этого каждый производитель процессоров делал
плавающую арифметику по-своему, и одна и та же программа на разных
машинах могла давать заметно разные ответы. Кэхэн также является
автором знаменитого \term{алгоритма компенсированного суммирования},
который позволяет складывать миллионы чисел с почти двойной
точностью.
\end{historybox}

% =====================================================================
\section{Сингулярное разложение SVD: главный фокус линейной алгебры}
\label{sec:linalg-svd}

\subsection{Геометрическая идея}

Возьмём любую матрицу $A$ размера $m\times n$ и применим её к
единичному кругу в $\R^{n}$ (то есть к множеству всех векторов $x$
с $\norm{x}=1$, где $\norm{\cdot}$ — обычная евклидова длина).
Что получится в $\R^{m}$?

\begin{theorem}{Сингулярное разложение (SVD)}{svd}
Для любой вещественной матрицы $A$ размера $m\times n$ существует
разложение
\begin{equation}\label{eq:svd}
\boxed{\;A \;=\; U\,\Sigma\,V^{\!\top}\;}
\end{equation}
где
\begin{itemize}\itemsep1pt
  \item $U$ — \emph{ортогональная} матрица $m\times m$ (поворот в $\R^{m}$);
  \item $V$ — ортогональная матрица $n\times n$ (поворот в $\R^{n}$);
  \item $\Sigma$ — «диагональная» матрица $m\times n$ с неотрицательными
        числами $\sigma_1\geq \sigma_2\geq\ldots\geq \sigma_r > 0$ на
        диагонали (и нулями ниже), где $r=\mathrm{rank}\,A$.
\end{itemize}
Числа $\sigma_i$ называются \term{сингулярными числами} матрицы~$A$.
\end{theorem}

Содержательно \eqref{eq:svd} говорит: \emph{любая} линейная операция —
это три простых действия подряд:
\[
\underbrace{V^{\!\top}}_{\text{поворот в исходном пространстве}}
\;\to\;
\underbrace{\Sigma}_{\text{растяжение по координатным осям}}
\;\to\;
\underbrace{U}_{\text{поворот в новом пространстве}}.
\]

Геометрически это значит: образ единичной сферы под действием $A$ —
это \term{эллипсоид}, у которого полуоси равны $\sigma_1,\sigma_2,\ldots$
(см.~рис.~\ref{fig:svd-ellipse}).

\begin{figure}[ht]
\centering
\begin{tikzpicture}[scale=1.05, line cap=round]
  % Левая картинка: единичный круг
  \begin{scope}[xshift=-3.6cm]
    \draw[->, gray!60] (-1.6,0) -- (1.6,0);
    \draw[->, gray!60] (0,-1.6) -- (0,1.6);
    \draw[thick, accent] (0,0) circle (1);
    \draw[->, very thick, accent2] (0,0) -- (1,0) node[below right] {$v_1$};
    \draw[->, very thick, accent2] (0,0) -- (0,1) node[above left] {$v_2$};
    \node[below] at (0,-1.85) {$\R^{n}$: $\norm{x}=1$};
  \end{scope}
  % Стрелка A
  \draw[->, very thick, accent] (-1.5,0) -- (1.5,0)
        node[midway, above, font=\large] {$A$};
  % Правая картинка: эллипс
  \begin{scope}[xshift=3.6cm, rotate=20]
    \draw[->, gray!60, rotate=-20] (-2.4,0) -- (2.4,0);
    \draw[->, gray!60, rotate=-20] (0,-1.4) -- (0,1.4);
    \draw[thick, accent] (0,0) ellipse (2 and 0.7);
    \draw[->, very thick, accent2] (0,0) -- (2,0) node[above] {$\sigma_1 u_1$};
    \draw[->, very thick, accent2] (0,0) -- (0,0.7) node[above left] {$\sigma_2 u_2$};
  \end{scope}
  \node[below] at (3.6,-1.85) {$\R^{m}$: $\{Ax:\norm{x}=1\}$};
\end{tikzpicture}
\caption{SVD геометрически. Матрица $A$ переводит единичный круг
в эллипс, полуоси которого равны сингулярным числам $\sigma_1,\sigma_2$,
а направления полуосей — столбцы матрицы $U$ (левые сингулярные
векторы). Векторы $v_1,v_2$ (столбцы $V$, правые сингулярные
векторы) — это те направления в исходном пространстве, которые
$A$ переводит в полуоси эллипса.}
\label{fig:svd-ellipse}
\end{figure}

\subsection{Связь с собственными значениями}

Сингулярные числа — это, по сути, переодетые собственные значения.
А именно:
\begin{equation}\label{eq:svd-eig}
A^{\!\top} A \;=\; (U\Sigma V^{\!\top})^{\!\top}(U\Sigma V^{\!\top})
\;=\; V\,\Sigma^{\!\top}\Sigma\,V^{\!\top}.
\end{equation}
Матрица $\Sigma^{\!\top}\Sigma$ — диагональная, на её диагонали
стоят $\sigma_1^{2},\sigma_2^{2},\ldots$. Значит, $\sigma_i^{2}$ —
это \emph{собственные значения} симметричной матрицы $A^{\!\top}A$,
а столбцы $V$ — её собственные векторы. Эта формула важна нам
по двум причинам: (1) она обосновывает существование SVD (через
известную спектральную теорему для симметричных матриц),
(2) она будет рабочим инструментом в § \ref{sec:linalg-eigenfaces},
когда мы будем вычислять собственные лица.

\begin{historybox}
Сингулярное разложение независимо открыли итальянский математик
\emph{Эудженио Бельтрами} (1873~г.) и француз \emph{Камилл Жордан}
(1874~г.). В современном виде, как разложение прямоугольной матрицы,
его сформулировал \emph{Карл Эккарт} с \emph{Гейлом Янгом} в 1936~г.
Численно надёжный алгоритм вычисления SVD появился лишь в 1965~г.
(Голуб и Кахан) — и стал одним из ключевых алгоритмов XX века. Сегодня
он реализован в любой математической библиотеке: \texttt{numpy.linalg.svd},
\texttt{scipy.linalg.svd}, \texttt{torch.linalg.svd}.
\end{historybox}

% =====================================================================
\section{Лучшее малоранговое приближение: теорема Эккарта--Янга}
\label{sec:linalg-eckart}

Самое замечательное свойство SVD состоит в том, что оно позволяет
\emph{ужимать} матрицу, выбрасывая всё «несущественное».

\subsection{Разложение по слагаемым ранга один}

Перепишем~\eqref{eq:svd} по-другому. Обозначим столбцы матрицы
$U$ через $u_1,\ldots,u_m$ (это \emph{левые сингулярные векторы}),
а столбцы $V$ — через $v_1,\ldots,v_n$ (\emph{правые сингулярные
векторы}). Тогда $A$ есть сумма $r$ слагаемых, каждое из которых
имеет ранг один:
\begin{equation}\label{eq:svd-sum}
A \;=\; \sum_{i=1}^{r}\sigma_i\,u_i v_i^{\!\top}.
\end{equation}

Каждое слагаемое $u_i v_i^{\!\top}$ — это «вертикальный вектор
умножить на горизонтальный», и оно само по себе является матрицей
ранга один (рис.~\ref{fig:rank1}).

\begin{figure}[ht]
\centering
\begin{tikzpicture}[scale=0.85]
  % колонка u
  \draw (0,0) rectangle (0.5,2.5);
  \node at (0.25,2.8) {$u$};
  \node at (0.25,-0.4) {\scriptsize $m\times 1$};
  \node at (1.2,1.25) {\Large $\times$};
  % строка v^T
  \draw (1.7,1.0) rectangle (4.2,1.5);
  \node at (2.95,1.7) {$v^{\!\top}$};
  \node at (2.95,0.5) {\scriptsize $1\times n$};
  \node at (4.8,1.25) {\Large $=$};
  % результат
  \draw[fill=accent!10] (5.3,0) rectangle (7.8,2.5);
  \node at (6.55,2.8) {$u v^{\!\top}$};
  \node at (6.55,-0.4) {\scriptsize $m\times n$, ранг $1$};
\end{tikzpicture}
\caption{Матрица ранга один: каждый столбец результата
пропорционален $u$, каждая строка — пропорциональна $v^{\!\top}$.
Чтобы её хранить, достаточно $m+n$ чисел вместо $mn$.}
\label{fig:rank1}
\end{figure}

Если оборвать сумму~\eqref{eq:svd-sum} на $k$-м слагаемом ($k<r$),
получим приближение:
\begin{equation}\label{eq:lowrank}
A_k \;=\; \sum_{i=1}^{k}\sigma_i\,u_i v_i^{\!\top}.
\end{equation}
Это матрица того же размера $m\times n$, но ранга всего лишь $k$.

\subsection{Теорема Эккарта--Янга: $A_k$ — наилучшее приближение}

\begin{theorem}{Эккарт--Янг--Мирский (1936/1960)}{eckart-young}
Среди всех матриц $B$ размера $m\times n$ ранга не более $k$
матрица $A_k$ из~\eqref{eq:lowrank} даёт \emph{наименьшую} ошибку
в смысле спектральной нормы:
\begin{equation}\label{eq:ey}
\min_{\mathrm{rank}\,B\,\leq\,k}\;\norm{A-B}_{2}\;=\;\norm{A-A_k}_{2}\;=\;\sigma_{k+1}.
\end{equation}
Здесь $\norm{M}_2 = \max_{\norm{x}=1}\norm{Mx}=\sigma_1(M)$ —
\term{спектральная норма} матрицы $M$ (равна её наибольшему
сингулярному числу). Утверждение остаётся в силе и для
\emph{фробениусовой} нормы $\norm{M}_F=\sqrt{\sum_{i,j} m_{ij}^{2}}$:
\[
\norm{A-A_k}_F\;=\;\sqrt{\sigma_{k+1}^{2}+\sigma_{k+2}^{2}+\ldots+\sigma_r^{2}}.
\]
\end{theorem}

Смысл теоремы: \emph{если хочется сжать матрицу до ранга $k$, лучшее,
что вы можете сделать, — оставить первые $k$ компонент SVD}.
Ошибка контролируется \emph{отброшенным} сингулярным числом
$\sigma_{k+1}$.

\subsection{Пример: сжатие изображения}

Возьмём фотографию $512\times 512$ в градациях серого. Хранение
«в лоб» — это $512\cdot 512 = 262\,144$ числа. Применим к матрице
фото SVD и оставим только $k$ старших компонент: придётся хранить
\[
k\cdot(512+512)+k\;=\;k\cdot 1025
\]
чисел — по одному столбцу $u_i$, одной строке $v_i^{\!\top}$ и
одному числу $\sigma_i$ на каждое слагаемое. Уже при $k=50$ это
$\sim 51\,250$ чисел — в \emph{пять} раз меньше исходного,
а визуально потеря почти незаметна (рис.~\ref{fig:img-compress}).

\begin{lstlisting}[label={lst:svd-compress}]
import numpy as np
import matplotlib.pyplot as plt
from skimage import data, color

# Load and convert to grayscale
img = color.rgb2gray(data.astronaut())        # shape (512, 512)
U, S, Vt = np.linalg.svd(img, full_matrices=False)

# Reconstruct using only top-k singular components
ranks = [1, 5, 20, 50, 200]
fig, axes = plt.subplots(1, len(ranks)+1, figsize=(14, 3))
axes[0].imshow(img, cmap='gray'); axes[0].set_title('original')
for ax, k in zip(axes[1:], ranks):
    A_k = U[:, :k] @ np.diag(S[:k]) @ Vt[:k, :]
    ax.imshow(A_k, cmap='gray')
    ax.set_title(f'rank {k}')
for ax in axes: ax.axis('off')
plt.tight_layout(); plt.show()
\end{lstlisting}

\begin{figure}[ht]
\centering
\includegraphics[width=\linewidth]{svd_image_compression.pdf}
\caption{Малоранговое приближение портрета астронавта. Уже при
ранге $50$ изображение визуально неотличимо от оригинала; при ранге
$200$ восстановление идеально. По теореме Эккарта--Янга это
\emph{лучшее} приближение фиксированного ранга.}
\label{fig:img-compress}
\end{figure}

\begin{notebox}[title={LoRA: малоранговая адаптация больших моделей}]
В 2021~г. Эдвард Ху с соавторами заметили: когда большую языковую
модель дообучают на новую задачу, \emph{обновление} матриц весов
$\Delta W$ оказывается почти малоранговой матрицей. Поэтому вместо
обучения «полной» $\Delta W$ можно сразу искать её в виде
$\Delta W = U V^{\!\top}$ с маленьким $k$, где $U,V$ —
тонкие прямоугольники.

Это метод \term{LoRA} (\engterm{Low-Rank Adaptation}). Для модели с
175 миллиардами параметров такой приём сокращает объём
дообучаемых весов в тысячи раз — и именно благодаря ему сегодня
можно «дофайнтюнить» большую модель на бытовой видеокарте.

Теоретический фундамент LoRA — та самая теорема
Эккарта--Янга~\ref{thm:eckart-young}, которой почти 90 лет.
\end{notebox}

% =====================================================================
\section{Eigenfaces: как алгоритм узнаёт лицо}
\label{sec:linalg-eigenfaces}

\subsection{Постановка задачи}

В $1991$~г. Мэттью Тёрк и Алекс Пентланд из MIT поставили простой
вопрос: \emph{можно ли обучить компьютер узнавать лицо, не прибегая
ни к каким специальным «лицевым» признакам — носу, глазам, контурам, —
а просто работая с матрицей пикселей?} Их ответ — метод
\term{собственных лиц} (\engterm{eigenfaces}) — стал классикой и
открыл целую область computer vision.

Идея: каждая фотография лица — это точка в очень многомерном
пространстве. Например, лицо $64\times 64$ — это вектор длины
$64\cdot 64 = 4\,096$ в пространстве $\R^{4096}$. Двух разных людей
снимали много раз — получаем облако точек в этом пространстве.

Главное наблюдение: \emph{это облако очень узкое}. Хотя пространство
$4\,096$-мерное, реальные лица занимают в нём подпространство
размерности всего лишь $\sim 50$. Найти это подпространство — и
есть задача.

\begin{figure}[ht]
\centering
\includegraphics[width=0.85\linewidth]{eigenfaces_dataset.pdf}
\caption{Фрагмент датасета \engterm{Olivetti Faces}: $400$ фотографий
$64\times 64$ пикселя, $40$ человек по $10$ снимков каждого. Каждая
фотография --- точка в $\R^{4096}$.}
\label{fig:olivetti}
\end{figure}

\subsection{Шаг 1. Среднее лицо и центрирование}

Пусть у нас $N$ фотографий $x_1,\ldots,x_N\in\R^{d}$ (где $d=4096$ для
$64\times 64$). Вычислим \term{среднее лицо}:
\[
\bar x \;=\; \dfrac{1}{N}\sum_{i=1}^{N} x_i.
\]
И вычтем его из каждого изображения: $\tilde x_i = x_i - \bar x$.
\emph{Среднее лицо} (рис.~\ref{fig:mean-face}) — это размытая
маска, в которой видно «лицо вообще»: усреднённые глаза, нос,
рот, овал. Дальше мы будем работать только с отклонениями
от этого среднего.

\begin{figure}[ht]
\centering
\includegraphics[width=0.55\linewidth]{eigenfaces_mean.pdf}
\caption{Среднее лицо по $400$ снимкам датасета Olivetti.}
\label{fig:mean-face}
\end{figure}

\subsection{Шаг 2. Главные направления изменчивости}

Соберём центрированные изображения в строки матрицы
$X\in\R^{N\times d}$ (одна строка = одно лицо). Применим к ней SVD:
\[
X\;=\;U\,\Sigma\,V^{\!\top}.
\]
Каждая \emph{строка} $v_i^{\!\top}$ матрицы $V^{\!\top}$ — это вектор
длины $d=4096$, который можно нарисовать как картинку $64\times 64$.
Эти картинки и называются \term{собственными лицами} — eigenfaces
(рис.~\ref{fig:eigenfaces}). Содержательно:
\begin{itemize}\itemsep1pt
  \item \emph{$v_1$} — направление наибольшей изменчивости.
        Чаще всего это «освещение слева/справа».
  \item \emph{$v_2$} — направление второй по величине изменчивости
        (часто «улыбка / не улыбка»).
  \item \emph{$v_3,\ldots$} — постепенно более тонкие признаки: очки,
        форма щёк, наклон головы, индивидуальные особенности.
\end{itemize}

\begin{figure}[ht]
\centering
\includegraphics[width=\linewidth]{eigenfaces_components.pdf}
\caption{Первые 12 eigenfaces датасета Olivetti, отрисованные как
картинки $64\times 64$. Эти странные «полу-лица» — координатные
оси в пространстве, где живут все настоящие лица.}
\label{fig:eigenfaces}
\end{figure}

\subsection{Шаг 3. Лицо как вектор из 50 чисел}

Зафиксируем $k$ (обычно $k=30\div 100$) и спроектируем каждое лицо
$x$ на эти $k$ собственных лиц:
\begin{equation}\label{eq:eigenface-proj}
\alpha_i \;=\; v_i^{\!\top}(x-\bar x),\qquad i=1,\ldots,k.
\end{equation}
Получили вектор коэффициентов $\alpha=(\alpha_1,\ldots,\alpha_k)\in\R^{k}$.
Это и есть «лицо человека» с точки зрения алгоритма — \emph{всего
$50$ чисел} вместо $4096$.

\subsection{Шаг 4. Распознавание и реконструкция}

\paragraph{Распознавание.} Чтобы узнать, кто на новой фотографии $y$,
вычислим её коэффициенты $\beta$ по~\eqref{eq:eigenface-proj} и
найдём в базе ближайшего соседа: $\hat j = \argmin_j \norm{\beta-\alpha^{(j)}}$.

\paragraph{Реконструкция.} Зная $\alpha$, можно восстановить лицо:
\[
x \;\approx\; \bar x + \sum_{i=1}^{k}\alpha_i\,v_i.
\]
Рис.~\ref{fig:eigenfaces-recon} показывает, как качество
реконструкции растёт с $k$.

\begin{figure}[ht]
\centering
\includegraphics[width=\linewidth]{eigenfaces_reconstruction.pdf}
\caption{Реконструкция одной из фотографий датасета через первые $k$
eigenfaces. При $k=10$ узнаётся «общий тип» лица; при $k=50$ — уже
конкретный человек; при $k=150$ — почти оригинал.}
\label{fig:eigenfaces-recon}
\end{figure}

\subsection{Реализация на Python}

\begin{lstlisting}[label={lst:eigenfaces}]
import numpy as np
import matplotlib.pyplot as plt
from sklearn.datasets import fetch_olivetti_faces

# 1. Load: 400 faces, 64x64 each
faces = fetch_olivetti_faces().images        # (400, 64, 64)
X = faces.reshape(400, -1)                   # (400, 4096)

# 2. Center: subtract the mean face
x_bar = X.mean(axis=0)
X_centered = X - x_bar

# 3. SVD => eigenfaces are rows of Vt
U, S, Vt = np.linalg.svd(X_centered, full_matrices=False)

# 4. Project all faces onto top-k eigenfaces
k = 50
alpha = X_centered @ Vt[:k].T                # (400, 50)

# 5. Recognise: nearest neighbour in alpha-space
def recognise(y, k=50):
    y_centered = y.flatten() - x_bar
    beta = Vt[:k] @ y_centered
    distances = np.linalg.norm(alpha - beta, axis=1)
    return distances.argmin()

# 6. Reconstruct a face from k coefficients
def reconstruct(idx, k=50):
    return x_bar + Vt[:k].T @ alpha[idx, :k]
\end{lstlisting}

\begin{interestbox}
Метод eigenfaces — родоначальник всех современных систем распознавания
лиц, включая ту, что разблокирует ваш смартфон. Современные системы
заменяют SVD на свёрточную нейросеть (см.~§\,3.3.7), но базовая
идея — \emph{сжать лицо в короткий вектор, в котором близкие лица
оказываются рядом} — осталась той же. Эти короткие векторы сегодня
называют \term{эмбеддингами}.
\end{interestbox}

% =====================================================================
\section{Eigencats: то же самое для котиков\texorpdfstring{$^\star$}{*}}
\label{sec:linalg-eigencats}

Кошки — отличный материал для проверки гипотезы, что метод eigenfaces
работает на любых однородных объектах, не только на лицах. Возьмём
датасет в $\sim 2000$ котиков\footnote{Например,
\engterm{Cats Faces Dataset} на Kaggle: фотографии $64\times 64$.
Готовые ноутбуки выложены в репозитории \href{https://github.com/MerkulovDaniil/nla360}{nla360}.}
и проведём тот же эксперимент.

\begin{figure}[ht]
\centering
\includegraphics[width=0.85\linewidth]{eigencats_components.pdf}
\caption{Первые 8 eigencats. Видны общие «коты»: положение ушей,
форма морды, направление взгляда. Каждый конкретный кот собирается
как линейная комбинация этих универсальных «прото-котов».}
\label{fig:eigencats}
\end{figure}

Качественная картина повторяется: первое eigencat ловит освещение,
второе — наклон головы, дальше идут вариации формы морды и ушей.
В сжатом виде каждый кот превращается в $\sim 50$ чисел.

\begin{notebox}[title={Что это нам говорит о ИИ}]
\textbf{Пространство «реальных» объектов всегда узкое.} Будь это
лица людей, морды котов, рукописные цифры (датасет MNIST) или
фотографии природы — реальные данные занимают в исходном
многомерном пространстве \emph{маленькое} подпространство.
Распознавание = найти это подпространство и работать в нём.

\textbf{Эта идея — фундамент современных автокодировщиков и
эмбеддингов} (см.~§\,3.4.3--3.4.6). SVD — её самая ранняя и самая
понятная реализация.
\end{notebox}

% =====================================================================
\section{PageRank: как Google нашёл порядок в Интернете}
\label{sec:linalg-pagerank}

\subsection{Задача и история}

К 1996~г. поисковики типа \engterm{AltaVista} и \engterm{Yahoo!}
ранжировали страницы по тому, сколько раз искомое слово встречается
в тексте. Это легко обманывалось: достаточно было набить страницу
популярными словами белым цветом по белому фону, и она попадала
в топ выдачи.

Лоуренс Пейдж и Сергей Брин, тогда аспиранты Стэнфорда, предложили
другой принцип ранжирования: \emph{страница важна, если на неё
ссылаются другие важные страницы}.

\subsection{Граф ссылок как матрица}

Обозначим $N$ — число всех страниц Интернета (на 1998~г. — около
$2{,}4\cdot 10^{8}$, сегодня — $\sim 10^{11}$). Построим матрицу
$P\in\R^{N\times N}$:
\[
P_{ij}\;=\;\begin{cases}
\dfrac{1}{c_j}, & \text{если со страницы $j$ есть ссылка на $i$,}\\
0, & \text{иначе,}
\end{cases}
\]
где $c_j$ — \emph{полное число исходящих ссылок} со страницы $j$.
Содержательно: $P_{ij}$ — это вероятность того, что пользователь,
сидящий на странице $j$ и случайно ткнувший на одну из её ссылок,
окажется на странице $i$. Поэтому $P$ называется матрицей переходов
\term{случайного блуждания} по графу Интернета.

\subsection{Стационарное распределение = собственный вектор}

Пусть $r\in\R^{N}$ — вектор «важностей»: $r_i$ — степень важности
страницы $i$. Принцип Пейджа–Брина формализуется системой:
\begin{equation}\label{eq:pr-fixed}
r\;=\;P r.
\end{equation}
То есть важность каждой страницы равна сумме важностей тех страниц,
которые на неё ссылаются, делённой на число их исходящих ссылок.

Уравнение~\eqref{eq:pr-fixed} говорит: $r$ — это \emph{собственный
вектор} матрицы $P$ с собственным значением $1$. И в этом — вся суть
PageRank.

\begin{theorem}{Перрон--Фробениус для PageRank}{perron}
Если все элементы матрицы $P$ положительны (или хотя бы граф ссылок
сильно связен и непериодичен), то собственное значение $1$ у $P$
существует, оно \emph{простое} и наибольшее по модулю, а
соответствующий собственный вектор $r$ единственен и имеет
положительные компоненты. Этот $r$ называется \term{стационарным
распределением} случайного блуждания.
\end{theorem}

(Эта теорема — Оскар Перрон, 1907~г., и Георг Фробениус, 1912~г. —
исторически появилась задолго до Интернета, в контексте моделей
популяций.)

\subsection{Демпфирующий множитель и формула Брина--Пейджа}

В реальной матрице Интернета у некоторых страниц нет исходящих ссылок
(\engterm{dangling pages}), а в других группах — циклы. Для таких
страниц соответствующий столбец матрицы переходов заранее заменяют
равномерным распределением по всем страницам: иначе формула $1/c_j$
при $c_j=0$ не определена, а масса случайного блуждания «исчезала» бы
в dangling pages. Чтобы
гарантировать применимость теоремы Перрона--Фробениуса, Пейдж и Брин
ввели \term{демпфирующий множитель} $\beta\in(0,1)$ (обычно
$\beta=0{,}85$):
\begin{equation}\label{eq:pagerank}
\boxed{\;r\;=\;\beta P r + \dfrac{1-\beta}{N}\,\mathbf{1}.\;}
\end{equation}
Здесь $\mathbf{1}=(1,1,\ldots,1)^{\!\top}$. Содержательно:
с вероятностью $\beta\approx 0{,}85$ случайный сёрфер кликает по
ссылке, а с вероятностью $1-\beta\approx 0{,}15$ — телепортируется
на случайную страницу Интернета.

\subsection{Как находить $r$: степенной метод}

Прямо решать систему $(I-\beta P)r=\frac{1-\beta}{N}\mathbf{1}$
с $N\sim 10^{11}$ невозможно. Но и не надо: вектор $r$ можно
получить простой итерацией.

\begin{algorithmbox}
\textbf{Степенной метод для PageRank.}
\begin{enumerate}\itemsep1pt
  \item Инициализация: $r^{(0)} \leftarrow \tfrac{1}{N}\mathbf{1}$.
  \item Итерация: $r^{(t+1)} \leftarrow \beta\,P\,r^{(t)} +
        \dfrac{1-\beta}{N}\,\mathbf{1}$.
  \item Остановка: $\norm{r^{(t+1)}-r^{(t)}} < \varepsilon$.
\end{enumerate}
\end{algorithmbox}

Сходимость гарантируется теоремой Перрона--Фробениуса:
$\norm{r^{(t)}-r}\leq C\cdot\beta^{t}$, то есть погрешность убывает
геометрически со знаменателем $\beta=0{,}85$. На практике хватает
$50\div 100$ итераций.

\subsection{Игрушечный пример: 6 страниц}

Возьмём маленький «Интернет» из шести страниц (рис.~\ref{fig:pagerank-graph})
и посчитаем для него PageRank.

\begin{figure}[ht]
\centering
\begin{tikzpicture}[scale=1.0, every node/.style={circle, draw, accent,
              fill=accent!10, minimum size=8mm, font=\small}]
  \node (A) at (0,2)    {A};
  \node (B) at (2.5,3)  {B};
  \node (C) at (5,2)    {C};
  \node (D) at (5,-0.5) {D};
  \node (E) at (2.5,-1.5){E};
  \node (F) at (0,-0.5) {F};
  \begin{scope}[->, thick, accent2]
    \draw (A) -- (B);
    \draw (A) -- (C);
    \draw (B) -- (C);
    \draw (C) -- (A);
    \draw (D) -- (C);
    \draw (D) -- (A);
    \draw (E) -- (D);
    \draw (E) -- (B);
    \draw (E) -- (F);
    \draw (F) -- (E);
    \draw (F) -- (B);
  \end{scope}
\end{tikzpicture}
\caption{Игрушечный «Интернет» из 6 страниц со ссылками. Видно,
например, что на страницу C ссылаются три другие, а на F — только
одна; будет ли C лидером по PageRank?}
\label{fig:pagerank-graph}
\end{figure}

\begin{lstlisting}[label={lst:pagerank}]
import numpy as np

# Adjacency list -> column-stochastic transition matrix P
links = {'A': ['B','C'], 'B': ['C'], 'C': ['A'],
         'D': ['C','A'], 'E': ['D','B','F'], 'F': ['E','B']}
pages = list(links); n = len(pages); idx = {p: i for i, p in enumerate(pages)}

P = np.zeros((n, n))
for j, p in enumerate(pages):
    out = links[p]
    for q in out:
        P[idx[q], j] = 1.0 / len(out)

# PageRank via power iteration
beta, eps = 0.85, 1e-10
r = np.ones(n) / n
for t in range(200):
    r_new = beta * P @ r + (1 - beta) / n
    if np.linalg.norm(r_new - r, 1) < eps: break
    r = r_new

for p, ri in sorted(zip(pages, r), key=lambda t: -t[1]):
    print(f'{p}: {ri:.4f}')
\end{lstlisting}

\noindent Выдача программы:
\[
\text{C: }0{,}295,\quad
\text{A: }0{,}252,\quad
\text{B: }0{,}175,\quad
\text{E: }0{,}098,\quad
\text{F: }0{,}092,\quad
\text{D: }0{,}087.
\]
Лидер — \texttt{C}, хотя на неё формально ссылаются столько же
страниц, сколько на \texttt{A}: PageRank учёл, что одна из ссылок
на \texttt{C} приходит с уже-важной страницы \texttt{B}.

\begin{historybox}
Алгоритм PageRank был опубликован в 1998~г. в студенческой работе
\emph{Sergey Brin, Larry Page}, «\textit{The PageRank Citation Ranking:
Bringing Order to the Web}». В том же году Брин и Пейдж основали
Google. К 2004~г. компания вышла на IPO, а к 2010-м стала одной из
крупнейших в мире. Сам Стэнфорд получил $\sim 336$ миллионов долларов
от лицензирования патента --- который изначально просто описывал
итерационный метод для собственного вектора.

Идея использовать собственный вектор графа для измерения «важности»
позднее была обобщена и применена в наукометрии (\engterm{Eigenfactor}
для ранжирования научных журналов), в рекомендательных системах,
а также в современных нейросетях для графов (\engterm{Graph Neural
Networks}).
\end{historybox}

% =====================================================================
\section{Когда нейросеть устойчива: оценки Липшицевости\texorpdfstring{$^\star$}{*}}
\label{sec:linalg-lipschitz}

\subsection{Сюжет: когда панда становится гиббоном}

В 2014~г. Иэн Гудфеллоу и соавторы опубликовали ставшую знаменитой
картинку (рис.~\ref{fig:adv-attack}): обычная фотография панды,
которую современная нейросеть уверенно опознаёт как «panda» с
вероятностью $57{,}7\%$. К каждому пикселю прибавлена крошечная,
визуально незаметная человеку поправка. На обновлённой картинке
сеть теперь уверенно (с вероятностью $99{,}3\%$) считает, что
видит \emph{гиббона}.

\begin{figure}[ht]
\centering
\includegraphics[width=0.95\linewidth]{adversarial_panda.pdf}
\caption{Adversarial-атака на нейросеть (Goodfellow et al., 2014).
Слева — оригинал, классифицирован как «panda» с уверенностью $57{,}7\%$.
В центре — почти невидимое возмущение, масштабированное для наглядности.
Справа — атакованный пример, классифицируется как «gibbon» с
уверенностью $99{,}3\%$.}
\label{fig:adv-attack}
\end{figure}

Этот эффект называется \term{adversarial-атакой} и долгое время
считался экзотикой. Сейчас понятно: причина в простом числе,
которое сопровождает любую функцию — её \emph{константе Липшица}.

\subsection{Константа Липшица: формальное определение}

\begin{definition}{Константа Липшица функции}{lipschitz}
Функция $f\colon\R^{n}\to\R^{m}$ называется \term{Липшицевой} с
константой $L>0$, если для любых $x,y\in\R^{n}$ выполнено
\begin{equation}\label{eq:lipschitz}
\norm{f(x)-f(y)}\;\leq\;L\,\norm{x-y}.
\end{equation}
Наименьшее $L$, для которого~\eqref{eq:lipschitz} выполнено,
называется \emph{константой Липшица} функции $f$ и обозначается
$\mathrm{Lip}(f)$.
\end{definition}

Содержательно $\mathrm{Lip}(f)$ показывает, во сколько раз может
\emph{усилиться} входное возмущение на выходе функции. Если
$\mathrm{Lip}(f)=L$, то малое возмущение $\delta$ на входе
может вызвать возмущение до $L\cdot\norm{\delta}$ на выходе.

\paragraph{Пример.} Для линейного отображения $f(x)=Wx$ имеем
$f(x)-f(y)=W(x-y)$, откуда $\norm{f(x)-f(y)}\leq \norm{W}_2\cdot\norm{x-y}$,
и константа Липшица — это в точности спектральная норма матрицы $W$:
\begin{equation}\label{eq:lip-linear}
\mathrm{Lip}(Wx)\;=\;\norm{W}_2\;=\;\sigma_1(W).
\end{equation}
Здесь $\sigma_1(W)$ — наибольшее сингулярное число матрицы $W$.
Вот \emph{здесь} в нашу историю входит SVD из §\,\ref{sec:linalg-svd}.

\subsection{Композиция: Липшицева константа цепочки}

\begin{theorem}{Композиция Липшицевых функций}{lip-compose}
Если $f$ и $g$ — Липшицевы с константами $L_f$ и $L_g$, то
их композиция $f\circ g$ — тоже Липшицева, причём
\begin{equation}\label{eq:lip-prod}
\mathrm{Lip}(f\circ g)\;\leq\;\mathrm{Lip}(f)\cdot\mathrm{Lip}(g).
\end{equation}
\end{theorem}

\begin{proof}
$\norm{f(g(x))-f(g(y))}\leq L_f\norm{g(x)-g(y)}\leq L_f L_g\norm{x-y}$. $\square$
\end{proof}

\subsection{Применение к нейросети}

Возьмём типичную нейросеть глубины $K$:
\[
f(x)\;=\;\sigma\bigl(W_K\,\sigma(W_{K-1}\,\sigma(\cdots\sigma(W_1 x))\bigr),
\]
где $W_i$ — матрицы весов, $\sigma$ — поэлементная активация (ReLU,
сигмоида, и т.п.). Для большинства активаций $\mathrm{Lip}(\sigma)\leq 1$
(для ReLU и $\tanh$ она равна $1$, а для логистической сигмоиды — $1/4$). По
теореме~\ref{thm:lip-compose}:
\begin{equation}\label{eq:nn-lip}
\boxed{\;\mathrm{Lip}(f)\;\leq\;\prod_{i=1}^{K}\norm{W_i}_2\;}
\end{equation}

Это и есть тот «лишний рычаг», который объясняет уязвимость
нейросетей. Если $\norm{W_i}_2 \approx 2$ для каждого слоя из $K=20$
слоёв, то возмущение на входе \emph{потенциально} усиливается
в $2^{20}\approx 10^{6}$ раз. Невидимая поправка $\delta$ с
$\norm{\delta}=10^{-3}$ может вызвать изменение выхода до $10^{3}$ —
этого хватит, чтобы перебросить классификатор из одного класса в другой.

\subsection{Как защищаться: спектральная нормализация}

Естественная стратегия: после каждого шага обучения нормализовать
матрицу $W_i$ так, чтобы $\norm{W_i}_2 = 1$. Это и есть метод
\term{Spectral Normalization} (Миято и соавт., 2018), сделавший
обучение генеративных нейросетей (GAN) намного стабильнее.

Технически: на каждом шаге считаем $\sigma_1(W)$ (с помощью одной
итерации степенного метода — почти бесплатно!) и заменяем
$W \mapsto W/\sigma_1(W)$.

\begin{lstlisting}[label={lst:spec-norm}]
def spectral_normalize(W, n_iter=1):
    """One step of power iteration to estimate sigma_1(W) and rescale."""
    u = np.random.randn(W.shape[0])
    for _ in range(n_iter):
        v = W.T @ u; v /= np.linalg.norm(v)
        u = W @ v;   u /= np.linalg.norm(u)
    sigma_1 = u @ W @ v
    return W / sigma_1
\end{lstlisting}

\begin{interestbox}
\textbf{Связь с устойчивостью в физике.} Похожая ситуация хорошо
известна в теории динамических систем: если константа Липшица
отображения $> 1$, малое возмущение начального условия экспоненциально
нарастает (классический эффект бабочки в погодных моделях).
Spectral Normalization можно рассматривать как искусственное
гашение этого эффекта.
\end{interestbox}

% =====================================================================
\section*{Подведение итогов главы}
\addcontentsline{toc}{section}{Подведение итогов главы}

\begin{itemize}\itemsep1pt
  \item Матрица — это \emph{универсальный носитель данных} для ИИ:
        картинки, тексты, графы, веса нейросети.
  \item \textbf{Числа с плавающей точкой} (§\ref{sec:linalg-floats}) —
        вычисления на компьютере всегда приближённые с относительной
        ошибкой $\sim\varepsilon_{\text{маш}}\approx 10^{-16}$.
        Самая опасная ошибка — \emph{катастрофическое сокращение} при
        вычитании близких чисел; её часто можно устранить, переписав
        формулу алгебраически эквивалентным способом.
  \item \textbf{Сингулярное разложение} $A=U\Sigma V^{\!\top}$
        (§\ref{sec:linalg-svd}) представляет любую матрицу как
        «поворот → растяжение по осям → поворот». Сингулярные числа
        $\sigma_i$ — это коэффициенты растяжения.
  \item \textbf{Теорема Эккарта--Янга} (§\ref{sec:linalg-eckart}):
        $A_k=\sum_{i\leq k}\sigma_i u_i v_i^{\!\top}$ — лучшее
        приближение ранга $k$ матрицы $A$. На этой теореме
        стоит сжатие изображений, метод eigenfaces и современный
        LoRA для дообучения больших языковых моделей.
  \item \textbf{Eigenfaces / eigencats}
        (§\ref{sec:linalg-eigenfaces}--\ref{sec:linalg-eigencats}):
        реальные данные (лица, котики, цифры) живут в очень узком
        подпространстве своего номинального многомерного пространства.
        Найти это подпространство = применить SVD к выборке.
  \item \textbf{PageRank} (§\ref{sec:linalg-pagerank}): «важность»
        страницы Интернета = компонента собственного вектора матрицы
        переходов случайного блуждания. Алгоритм находит его простой
        итерацией; теорема Перрона--Фробениуса гарантирует сходимость.
  \item \textbf{Константа Липшица} нейросети
        (§\ref{sec:linalg-lipschitz}) растёт мультипликативно от
        слоя к слою; именно поэтому глубокие сети уязвимы к
        adversarial-атакам. Спектральная нормализация ограничивает
        $\norm{W_i}_2=1$ и стабилизирует обучение.
\end{itemize}

% =====================================================================
\section*{Что почитать дальше}
\addcontentsline{toc}{section}{Что почитать дальше}

\begin{itemize}\itemsep1pt
  \item \emph{Е.\,Е.~Тыртышников.} Методы численного анализа. М.: Академия, 2007. ---
        классический российский учебник по вычислительной линейной
        алгебре университетского уровня.
  \item \emph{G.~Strang.} Linear Algebra and Learning from Data. Wellesley-Cambridge Press, 2019. ---
        свежая книга от автора знаменитого MIT-курса 18.06, написанная
        с прицелом на data science.
  \item \emph{L.\,N.~Trefethen, D.~Bau.} Numerical Linear Algebra. SIAM, 1997. ---
        классика, с очень понятным изложением SVD и итерационных методов.
  \item \emph{Видеокурс «Numerical Linear Algebra»}, И.\,В.~Оселедец, Сколтех
        (\href{https://github.com/MerkulovDaniil/nla360}{github.com/MerkulovDaniil/nla360}).
        Лекции, семинары, ноутбуки --- из которых, в частности,
        и взяты eigencats для этой главы.
  \item \emph{S.~Brin, L.~Page.} The Anatomy of a Large-Scale
        Hypertextual Web Search Engine. WWW Conference, 1998. ---
        оригинальная статья про PageRank, читается легко.
  \item \emph{I.~Goodfellow et al.} Explaining and Harnessing
        Adversarial Examples. ICLR, 2015. ---
        статья про панду-гиббона, из которой взят рис.~\ref{fig:adv-attack}.
\end{itemize}

% =====================================================================
\section*{Большой итоговый проект}
\addcontentsline{toc}{section}{Большой итоговый проект}

\textbf{Постройте свою рекомендательную систему для фильмов.}

Возьмите датасет MovieLens-100K (открытый, $100\,000$ оценок,
$\sim 1000$ пользователей $\times$ $\sim 1700$ фильмов). Постройте
матрицу «пользователь--фильм» $R$, где $R_{ij}\in\{1,\ldots,5\}$ —
оценка, поставленная пользователем $i$ фильму $j$, а $0$ — «не смотрел».

\textbf{Важно:} нули в матрице $R$ означают \emph{пропущенные} оценки,
а не настоящие нулевые рейтинги. Прямое SVD по матрице с нулями — лишь
грубый учебный baseline: непросмотренные фильмы тогда воспринимаются как
фильмы с очень низкой оценкой. Более корректный вариант — обучать
малоранговую модель только на известных оценках либо предварительно
заполнить пропуски средними значениями пользователя/фильма и затем
сравнить результат.

\begin{enumerate}\itemsep1pt
  \item Постройте малоранговое приближение $R\approx U\Sigma V^{\!\top}$
        с $k=20$ компонентами через SVD.
  \item Объясните содержательно, что такое строки $U$ (профили
        пользователей) и столбцы $V$ (профили фильмов) в этом
        приближении. Сколько чисел теперь хранит модель вместо
        $1000\cdot 1700$?
  \item Для нескольких реальных пользователей выпишите топ-5 фильмов,
        которые рекомендация модели предлагает посмотреть. Совпадают
        ли они с любимыми жанрами пользователя?
  \item Найдите два «противоположных» фильма с наибольшим расстоянием
        между их столбцами $V$ и подумайте, можно ли увидеть содержательный
        смысл в этой разделяющей оси (боевик vs. мелодрама? старое vs. новое?).
  \item ($\star$) Замените SVD на алгоритм PageRank, построив граф
        «фильмы, которые часто оценивают одинаково». Получится ли тот
        же топ?
\end{enumerate}

% =====================================================================
\begin{tasksbox}
\begin{enumerate}\itemsep1pt
  \item Сложите в Python $10\,000\,000$ копий числа $0{,}1$. Сколько
        должно получиться? Сколько получается на самом деле? Почему?
  \item Вычислите $f(x)=\sqrt{x+1}-\sqrt{x}$ при $x=10^{16}$
        напрямую и через формулу $f(x)=\frac{1}{\sqrt{x+1}+\sqrt{x}}$.
        Сравните ответы. Какой устойчив, почему?
  \item Найдите вручную SVD матрицы
        $A=\bigl(\begin{smallmatrix}3 & 0\\0 & 4\end{smallmatrix}\bigr)$
        и $A=\bigl(\begin{smallmatrix}0 & 1\\-1 & 0\end{smallmatrix}\bigr)$.
        Что в обоих случаях геометрически делают $U, \Sigma, V$?
  \item Возьмите свою фотографию $300\times 300$ в градациях серого.
        Сделайте сжатие через SVD с рангами $k=5, 20, 50, 100$. При
        каком $k$ вы перестаёте узнавать самого себя?
  \item Для игрушечного графа из рис.~\ref{fig:pagerank-graph}
        прокрутите 5 итераций степенного метода \emph{вручную},
        начав с $r^{(0)}=(\tfrac16,\ldots,\tfrac16)^{\!\top}$. Совпали ли
        ваши значения с программными после $5$ шагов? а после $20$?
  \item $^{\star}$ Покажите, что для функции $\sigma(x)=\max(0,x)$
        (ReLU) константа Липшица равна $1$. Тот же вопрос для
        сигмоиды $\sigma(x)=1/(1+e^{-x})$: чему равна $\mathrm{Lip}(\sigma)$
        и в какой точке производная максимальна?
  \item $^{\star}$ Найдите спектральную норму матрицы
        $W=\bigl(\begin{smallmatrix}1 & 2\\2 & 1\end{smallmatrix}\bigr)$
        двумя способами: (а) через определение
        $\norm{W}_2=\max_{\norm{x}=1}\norm{Wx}$ (с использованием
        собственных значений $W^{\!\top}W$); (б) одной итерацией
        степенного метода (с любым стартом). Совпали ли ответы?
\end{enumerate}
\end{tasksbox}
      % <— сюда
%
%  Все используемые окружения определены в structure.tex и
%  chapters/extras.tex; новых пакетов не требуется.
% =====================================================================

\chapter{Вычислительная линейная алгебра}
\label{ch:linalg}

\emph{Если метод Ньютона из главы~\ref{ch:opt} был одним
маленьким уравнением, упрямо повторяющимся, то линейная алгебра —
это язык, на котором говорят сразу о тысячах уравнений и миллионах
чисел. Этой главой мы покажем, что многие фокусы современного
искусственного интеллекта — от того, как нейросеть «узнаёт» лицо
человека, до того, как поисковик ранжирует страницы Интернета, —
держатся всего на нескольких разложениях матриц. Мы посмотрим
четыре истории: лица, котиков, Интернет и устойчивость нейросетей.}

\section{Зачем линейной алгебре отдельная глава}
\label{sec:linalg-why}

Вы уже встречались с матрицами в школьном курсе — как с прямоугольными
таблицами чисел. В современном ИИ матрица — это \term{универсальный
носитель данных}.

\begin{itemize}\itemsep1pt
  \item Чёрно-белая фотография $1024\times 1024$ — это матрица
        размера $1024\times 1024$, где каждое число — яркость пикселя
        от $0$ (чёрный) до $255$ (белый). Цветная — три такие матрицы.
  \item Текст из $N$ слов словаря — это вектор частот длины $N$
        (модель «мешка слов»), и набор из $M$ документов — матрица
        $M\times N$.
  \item Социальная сеть из $N$ пользователей — матрица $N\times N$,
        где элемент $(i,j)$ равен $1$, если $i$ подписан на $j$.
  \item Веса одного слоя нейросети — матрица $W$, через которую
        входной сигнал умножается, чтобы получить выходной.
\end{itemize}

С такими матрицами надо уметь \emph{что-то делать}: сжимать,
сравнивать, упорядочивать, обновлять. И почти всегда оказывается,
что нужная операция выражается через одно и то же фундаментальное
разложение матрицы — \term{сингулярное}. Мы изучим его и применим
к четырём задачам:

\begin{enumerate}\itemsep1pt
  \item \textbf{Узнать лицо} по фотографии — \emph{eigenfaces}
        (§\ref{sec:linalg-eigenfaces}).
  \item \textbf{Сжать пространство признаков} — \emph{eigencats} и
        теорема Эккарта–Янга (§\ref{sec:linalg-eckart},
        §\ref{sec:linalg-eigencats}).
  \item \textbf{Упорядочить весь Интернет} — \emph{PageRank}
        (§\ref{sec:linalg-pagerank}).
  \item \textbf{Понять, когда нейросеть устойчива} — \emph{оценки
        Липшицевости} (§\ref{sec:linalg-lipschitz}).
\end{enumerate}

Прежде чем перейти к сюжетам, разберёмся с одной \emph{очень} земной
проблемой: с тем, как компьютер хранит числа.

% =====================================================================
\section{Когда арифметика лжёт: числа с плавающей точкой}
\label{sec:linalg-floats}

Откройте Python и наберите:
\begin{lstlisting}[label={lst:float-paradox}]
>>> 0.1 + 0.2
0.30000000000000004
>>> 0.1 + 0.2 == 0.3
False
\end{lstlisting}

Это не баг Python. То же самое выдаст любой современный язык —
C, C++, Java, JavaScript. И этому есть фундаментальная причина:
\term{двоичная система счисления} не умеет точно представить
десятичную дробь $0{,}1$.

\subsection{Двоичные дроби и почему $0{,}1$ — это «$0{,}333\ldots$»}

В десятичной системе $1/3 = 0{,}3333\ldots$ — бесконечная дробь.
То же самое в двоичной системе случается с $0{,}1$:
\[
0{,}1_{10} \;=\; 0{,}0001\,1001\,1001\,1001\ldots_{2}
\]
Дробь периодическая. Любой реальный компьютер хранит её, оборвав
после конечного числа разрядов (обычно 52 двоичных знака после
запятой — стандарт \engterm{IEEE~754, double precision}). Поэтому
число, которое программист записывает как \texttt{0.1}, на самом
деле равно
\[
0{,}1\;\approx\;0{,}100\,000\,000\,000\,000\,005\,551\,115\ldots
\]
Когда мы складываем такое «искажённое $0{,}1$» с искажённым $0{,}2$,
ошибки усиливаются — и получается тот самый «лишний хвостик»
$4\cdot 10^{-17}$, который мы видели в листинге.

\begin{notebox}[title={Машинная точность $\varepsilon_{\text{маш}}$}]
\term{Машинной точностью} называют наименьшее положительное число
$\varepsilon$, такое что в компьютерной арифметике
$1+\varepsilon \neq 1$. Для стандарта \texttt{double} оно равно
$\varepsilon_{\text{маш}} = 2^{-52}\approx 2{,}22\cdot 10^{-16}$.

Это означает: какое бы хорошее число $x$ вы ни записали, оно
известно компьютеру лишь с относительной точностью
$\sim\varepsilon_{\text{маш}}$.
\end{notebox}

\subsection{Когда это становится опасным: катастрофическое сокращение}

Маленькая ошибка в $17$-м знаке — мелочь? Не всегда. Рассмотрим
безобидное на вид вычисление:
\[
g(x) \;=\; \dfrac{1-\cos x}{x^{2}}\qquad\text{при }x=10^{-8}.
\]
По формуле Тейлора $1-\cos x \approx x^{2}/2$, так что
$g(10^{-8})\approx 0{,}5$. Что выдаст компьютер?

\begin{lstlisting}[label={lst:cancel}]
>>> from math import cos
>>> x = 1e-8
>>> (1 - cos(x)) / x**2
0.0           # вместо 0.5!
\end{lstlisting}

Что произошло? Число $\cos(10^{-8})\approx 1-5\cdot 10^{-17}$ так
близко к единице, что компьютер просто хранит его как
\emph{ровно единицу} (поправка $5\cdot 10^{-17}$ меньше
$\varepsilon_{\text{маш}}$). Разность $1-\cos x$ обнулилась, и весь
ответ обнулился вместе с ней.

Это явление называется \term{катастрофическим сокращением}
(\engterm{catastrophic cancellation}): когда из двух близких чисел
вычитают и теряют все значащие цифры разом.

\begin{example}{Спасение через переписывание формулы}{cancel-fix}
Воспользуемся тождеством $1-\cos x = 2\sin^{2}(x/2)$:
\[
g(x) \;=\; \dfrac{2\sin^{2}(x/2)}{x^{2}}\;=\;\dfrac12\!\left(\dfrac{\sin(x/2)}{x/2}\right)^{2}.
\]
Внутри теперь нет вычитания близких чисел, и Python выдаёт
\texttt{0.49999999999999983} — погрешность лишь в последнем знаке.
\end{example}

\begin{importantbox}
\textbf{Главный практический вывод.} Если ваш численный алгоритм
выдаёт странный ответ, первая гипотеза, которую следует проверить, —
это не «бага в библиотеке», а \emph{катастрофическое сокращение} в
вашей формуле. Часто его можно убрать, переписав формулу
алгебраически эквивалентным, но численно устойчивым способом.
\end{importantbox}

\begin{historybox}
Стандарт \engterm{IEEE~754} был принят в 1985~г., в значительной
мере благодаря Уильяму Кэхэну (\engterm{William Kahan}, премия
Тьюринга 1989~г.). До этого каждый производитель процессоров делал
плавающую арифметику по-своему, и одна и та же программа на разных
машинах могла давать заметно разные ответы. Кэхэн также является
автором знаменитого \term{алгоритма компенсированного суммирования},
который позволяет складывать миллионы чисел с почти двойной
точностью.
\end{historybox}

% =====================================================================
\section{Сингулярное разложение SVD: главный фокус линейной алгебры}
\label{sec:linalg-svd}

\subsection{Геометрическая идея}

Возьмём любую матрицу $A$ размера $m\times n$ и применим её к
единичному кругу в $\R^{n}$ (то есть к множеству всех векторов $x$
с $\norm{x}=1$, где $\norm{\cdot}$ — обычная евклидова длина).
Что получится в $\R^{m}$?

\begin{theorem}{Сингулярное разложение (SVD)}{svd}
Для любой вещественной матрицы $A$ размера $m\times n$ существует
разложение
\begin{equation}\label{eq:svd}
\boxed{\;A \;=\; U\,\Sigma\,V^{\!\top}\;}
\end{equation}
где
\begin{itemize}\itemsep1pt
  \item $U$ — \emph{ортогональная} матрица $m\times m$ (поворот в $\R^{m}$);
  \item $V$ — ортогональная матрица $n\times n$ (поворот в $\R^{n}$);
  \item $\Sigma$ — «диагональная» матрица $m\times n$ с неотрицательными
        числами $\sigma_1\geq \sigma_2\geq\ldots\geq \sigma_r > 0$ на
        диагонали (и нулями ниже), где $r=\mathrm{rank}\,A$.
\end{itemize}
Числа $\sigma_i$ называются \term{сингулярными числами} матрицы~$A$.
\end{theorem}

Содержательно \eqref{eq:svd} говорит: \emph{любая} линейная операция —
это три простых действия подряд:
\[
\underbrace{V^{\!\top}}_{\text{поворот в исходном пространстве}}
\;\to\;
\underbrace{\Sigma}_{\text{растяжение по координатным осям}}
\;\to\;
\underbrace{U}_{\text{поворот в новом пространстве}}.
\]

Геометрически это значит: образ единичной сферы под действием $A$ —
это \term{эллипсоид}, у которого полуоси равны $\sigma_1,\sigma_2,\ldots$
(см.~рис.~\ref{fig:svd-ellipse}).

\begin{figure}[ht]
\centering
\begin{tikzpicture}[scale=1.05, line cap=round]
  % Левая картинка: единичный круг
  \begin{scope}[xshift=-3.6cm]
    \draw[->, gray!60] (-1.6,0) -- (1.6,0);
    \draw[->, gray!60] (0,-1.6) -- (0,1.6);
    \draw[thick, accent] (0,0) circle (1);
    \draw[->, very thick, accent2] (0,0) -- (1,0) node[below right] {$v_1$};
    \draw[->, very thick, accent2] (0,0) -- (0,1) node[above left] {$v_2$};
    \node[below] at (0,-1.85) {$\R^{n}$: $\norm{x}=1$};
  \end{scope}
  % Стрелка A
  \draw[->, very thick, accent] (-1.5,0) -- (1.5,0)
        node[midway, above, font=\large] {$A$};
  % Правая картинка: эллипс
  \begin{scope}[xshift=3.6cm, rotate=20]
    \draw[->, gray!60, rotate=-20] (-2.4,0) -- (2.4,0);
    \draw[->, gray!60, rotate=-20] (0,-1.4) -- (0,1.4);
    \draw[thick, accent] (0,0) ellipse (2 and 0.7);
    \draw[->, very thick, accent2] (0,0) -- (2,0) node[above] {$\sigma_1 u_1$};
    \draw[->, very thick, accent2] (0,0) -- (0,0.7) node[above left] {$\sigma_2 u_2$};
  \end{scope}
  \node[below] at (3.6,-1.85) {$\R^{m}$: $\{Ax:\norm{x}=1\}$};
\end{tikzpicture}
\caption{SVD геометрически. Матрица $A$ переводит единичный круг
в эллипс, полуоси которого равны сингулярным числам $\sigma_1,\sigma_2$,
а направления полуосей — столбцы матрицы $U$ (левые сингулярные
векторы). Векторы $v_1,v_2$ (столбцы $V$, правые сингулярные
векторы) — это те направления в исходном пространстве, которые
$A$ переводит в полуоси эллипса.}
\label{fig:svd-ellipse}
\end{figure}

\subsection{Связь с собственными значениями}

Сингулярные числа — это, по сути, переодетые собственные значения.
А именно:
\begin{equation}\label{eq:svd-eig}
A^{\!\top} A \;=\; (U\Sigma V^{\!\top})^{\!\top}(U\Sigma V^{\!\top})
\;=\; V\,\Sigma^{\!\top}\Sigma\,V^{\!\top}.
\end{equation}
Матрица $\Sigma^{\!\top}\Sigma$ — диагональная, на её диагонали
стоят $\sigma_1^{2},\sigma_2^{2},\ldots$. Значит, $\sigma_i^{2}$ —
это \emph{собственные значения} симметричной матрицы $A^{\!\top}A$,
а столбцы $V$ — её собственные векторы. Эта формула важна нам
по двум причинам: (1) она обосновывает существование SVD (через
известную спектральную теорему для симметричных матриц),
(2) она будет рабочим инструментом в § \ref{sec:linalg-eigenfaces},
когда мы будем вычислять собственные лица.

\begin{historybox}
Сингулярное разложение независимо открыли итальянский математик
\emph{Эудженио Бельтрами} (1873~г.) и француз \emph{Камилл Жордан}
(1874~г.). В современном виде, как разложение прямоугольной матрицы,
его сформулировал \emph{Карл Эккарт} с \emph{Гейлом Янгом} в 1936~г.
Численно надёжный алгоритм вычисления SVD появился лишь в 1965~г.
(Голуб и Кахан) — и стал одним из ключевых алгоритмов XX века. Сегодня
он реализован в любой математической библиотеке: \texttt{numpy.linalg.svd},
\texttt{scipy.linalg.svd}, \texttt{torch.linalg.svd}.
\end{historybox}

% =====================================================================
\section{Лучшее малоранговое приближение: теорема Эккарта--Янга}
\label{sec:linalg-eckart}

Самое замечательное свойство SVD состоит в том, что оно позволяет
\emph{ужимать} матрицу, выбрасывая всё «несущественное».

\subsection{Разложение по слагаемым ранга один}

Перепишем~\eqref{eq:svd} по-другому. Обозначим столбцы матрицы
$U$ через $u_1,\ldots,u_m$ (это \emph{левые сингулярные векторы}),
а столбцы $V$ — через $v_1,\ldots,v_n$ (\emph{правые сингулярные
векторы}). Тогда $A$ есть сумма $r$ слагаемых, каждое из которых
имеет ранг один:
\begin{equation}\label{eq:svd-sum}
A \;=\; \sum_{i=1}^{r}\sigma_i\,u_i v_i^{\!\top}.
\end{equation}

Каждое слагаемое $u_i v_i^{\!\top}$ — это «вертикальный вектор
умножить на горизонтальный», и оно само по себе является матрицей
ранга один (рис.~\ref{fig:rank1}).

\begin{figure}[ht]
\centering
\begin{tikzpicture}[scale=0.85]
  % колонка u
  \draw (0,0) rectangle (0.5,2.5);
  \node at (0.25,2.8) {$u$};
  \node at (0.25,-0.4) {\scriptsize $m\times 1$};
  \node at (1.2,1.25) {\Large $\times$};
  % строка v^T
  \draw (1.7,1.0) rectangle (4.2,1.5);
  \node at (2.95,1.7) {$v^{\!\top}$};
  \node at (2.95,0.5) {\scriptsize $1\times n$};
  \node at (4.8,1.25) {\Large $=$};
  % результат
  \draw[fill=accent!10] (5.3,0) rectangle (7.8,2.5);
  \node at (6.55,2.8) {$u v^{\!\top}$};
  \node at (6.55,-0.4) {\scriptsize $m\times n$, ранг $1$};
\end{tikzpicture}
\caption{Матрица ранга один: каждый столбец результата
пропорционален $u$, каждая строка — пропорциональна $v^{\!\top}$.
Чтобы её хранить, достаточно $m+n$ чисел вместо $mn$.}
\label{fig:rank1}
\end{figure}

Если оборвать сумму~\eqref{eq:svd-sum} на $k$-м слагаемом ($k<r$),
получим приближение:
\begin{equation}\label{eq:lowrank}
A_k \;=\; \sum_{i=1}^{k}\sigma_i\,u_i v_i^{\!\top}.
\end{equation}
Это матрица того же размера $m\times n$, но ранга всего лишь $k$.

\subsection{Теорема Эккарта--Янга: $A_k$ — наилучшее приближение}

\begin{theorem}{Эккарт--Янг--Мирский (1936/1960)}{eckart-young}
Среди всех матриц $B$ размера $m\times n$ ранга не более $k$
матрица $A_k$ из~\eqref{eq:lowrank} даёт \emph{наименьшую} ошибку
в смысле спектральной нормы:
\begin{equation}\label{eq:ey}
\min_{\mathrm{rank}\,B\,\leq\,k}\;\norm{A-B}_{2}\;=\;\norm{A-A_k}_{2}\;=\;\sigma_{k+1}.
\end{equation}
Здесь $\norm{M}_2 = \max_{\norm{x}=1}\norm{Mx}=\sigma_1(M)$ —
\term{спектральная норма} матрицы $M$ (равна её наибольшему
сингулярному числу). Утверждение остаётся в силе и для
\emph{фробениусовой} нормы $\norm{M}_F=\sqrt{\sum_{i,j} m_{ij}^{2}}$:
\[
\norm{A-A_k}_F\;=\;\sqrt{\sigma_{k+1}^{2}+\sigma_{k+2}^{2}+\ldots+\sigma_r^{2}}.
\]
\end{theorem}

Смысл теоремы: \emph{если хочется сжать матрицу до ранга $k$, лучшее,
что вы можете сделать, — оставить первые $k$ компонент SVD}.
Ошибка контролируется \emph{отброшенным} сингулярным числом
$\sigma_{k+1}$.

\subsection{Пример: сжатие изображения}

Возьмём фотографию $512\times 512$ в градациях серого. Хранение
«в лоб» — это $512\cdot 512 = 262\,144$ числа. Применим к матрице
фото SVD и оставим только $k$ старших компонент: придётся хранить
\[
k\cdot(512+512)+k\;=\;k\cdot 1025
\]
чисел — по одному столбцу $u_i$, одной строке $v_i^{\!\top}$ и
одному числу $\sigma_i$ на каждое слагаемое. Уже при $k=50$ это
$\sim 51\,250$ чисел — в \emph{пять} раз меньше исходного,
а визуально потеря почти незаметна (рис.~\ref{fig:img-compress}).

\begin{lstlisting}[label={lst:svd-compress}]
import numpy as np
import matplotlib.pyplot as plt
from skimage import data, color

# Load and convert to grayscale
img = color.rgb2gray(data.astronaut())        # shape (512, 512)
U, S, Vt = np.linalg.svd(img, full_matrices=False)

# Reconstruct using only top-k singular components
ranks = [1, 5, 20, 50, 200]
fig, axes = plt.subplots(1, len(ranks)+1, figsize=(14, 3))
axes[0].imshow(img, cmap='gray'); axes[0].set_title('original')
for ax, k in zip(axes[1:], ranks):
    A_k = U[:, :k] @ np.diag(S[:k]) @ Vt[:k, :]
    ax.imshow(A_k, cmap='gray')
    ax.set_title(f'rank {k}')
for ax in axes: ax.axis('off')
plt.tight_layout(); plt.show()
\end{lstlisting}

\begin{figure}[ht]
\centering
\includegraphics[width=\linewidth]{svd_image_compression.pdf}
\caption{Малоранговое приближение портрета астронавта. Уже при
ранге $50$ изображение визуально неотличимо от оригинала; при ранге
$200$ восстановление идеально. По теореме Эккарта--Янга это
\emph{лучшее} приближение фиксированного ранга.}
\label{fig:img-compress}
\end{figure}

\begin{notebox}[title={LoRA: малоранговая адаптация больших моделей}]
В 2021~г. Эдвард Ху с соавторами заметили: когда большую языковую
модель дообучают на новую задачу, \emph{обновление} матриц весов
$\Delta W$ оказывается почти малоранговой матрицей. Поэтому вместо
обучения «полной» $\Delta W$ можно сразу искать её в виде
$\Delta W = U V^{\!\top}$ с маленьким $k$, где $U,V$ —
тонкие прямоугольники.

Это метод \term{LoRA} (\engterm{Low-Rank Adaptation}). Для модели с
175 миллиардами параметров такой приём сокращает объём
дообучаемых весов в тысячи раз — и именно благодаря ему сегодня
можно «дофайнтюнить» большую модель на бытовой видеокарте.

Теоретический фундамент LoRA — та самая теорема
Эккарта--Янга~\ref{thm:eckart-young}, которой почти 90 лет.
\end{notebox}

% =====================================================================
\section{Eigenfaces: как алгоритм узнаёт лицо}
\label{sec:linalg-eigenfaces}

\subsection{Постановка задачи}

В $1991$~г. Мэттью Тёрк и Алекс Пентланд из MIT поставили простой
вопрос: \emph{можно ли обучить компьютер узнавать лицо, не прибегая
ни к каким специальным «лицевым» признакам — носу, глазам, контурам, —
а просто работая с матрицей пикселей?} Их ответ — метод
\term{собственных лиц} (\engterm{eigenfaces}) — стал классикой и
открыл целую область computer vision.

Идея: каждая фотография лица — это точка в очень многомерном
пространстве. Например, лицо $64\times 64$ — это вектор длины
$64\cdot 64 = 4\,096$ в пространстве $\R^{4096}$. Двух разных людей
снимали много раз — получаем облако точек в этом пространстве.

Главное наблюдение: \emph{это облако очень узкое}. Хотя пространство
$4\,096$-мерное, реальные лица занимают в нём подпространство
размерности всего лишь $\sim 50$. Найти это подпространство — и
есть задача.

\begin{figure}[ht]
\centering
\includegraphics[width=0.85\linewidth]{eigenfaces_dataset.pdf}
\caption{Фрагмент датасета \engterm{Olivetti Faces}: $400$ фотографий
$64\times 64$ пикселя, $40$ человек по $10$ снимков каждого. Каждая
фотография --- точка в $\R^{4096}$.}
\label{fig:olivetti}
\end{figure}

\subsection{Шаг 1. Среднее лицо и центрирование}

Пусть у нас $N$ фотографий $x_1,\ldots,x_N\in\R^{d}$ (где $d=4096$ для
$64\times 64$). Вычислим \term{среднее лицо}:
\[
\bar x \;=\; \dfrac{1}{N}\sum_{i=1}^{N} x_i.
\]
И вычтем его из каждого изображения: $\tilde x_i = x_i - \bar x$.
\emph{Среднее лицо} (рис.~\ref{fig:mean-face}) — это размытая
маска, в которой видно «лицо вообще»: усреднённые глаза, нос,
рот, овал. Дальше мы будем работать только с отклонениями
от этого среднего.

\begin{figure}[ht]
\centering
\includegraphics[width=0.55\linewidth]{eigenfaces_mean.pdf}
\caption{Среднее лицо по $400$ снимкам датасета Olivetti.}
\label{fig:mean-face}
\end{figure}

\subsection{Шаг 2. Главные направления изменчивости}

Соберём центрированные изображения в строки матрицы
$X\in\R^{N\times d}$ (одна строка = одно лицо). Применим к ней SVD:
\[
X\;=\;U\,\Sigma\,V^{\!\top}.
\]
Каждая \emph{строка} $v_i^{\!\top}$ матрицы $V^{\!\top}$ — это вектор
длины $d=4096$, который можно нарисовать как картинку $64\times 64$.
Эти картинки и называются \term{собственными лицами} — eigenfaces
(рис.~\ref{fig:eigenfaces}). Содержательно:
\begin{itemize}\itemsep1pt
  \item \emph{$v_1$} — направление наибольшей изменчивости.
        Чаще всего это «освещение слева/справа».
  \item \emph{$v_2$} — направление второй по величине изменчивости
        (часто «улыбка / не улыбка»).
  \item \emph{$v_3,\ldots$} — постепенно более тонкие признаки: очки,
        форма щёк, наклон головы, индивидуальные особенности.
\end{itemize}

\begin{figure}[ht]
\centering
\includegraphics[width=\linewidth]{eigenfaces_components.pdf}
\caption{Первые 12 eigenfaces датасета Olivetti, отрисованные как
картинки $64\times 64$. Эти странные «полу-лица» — координатные
оси в пространстве, где живут все настоящие лица.}
\label{fig:eigenfaces}
\end{figure}

\subsection{Шаг 3. Лицо как вектор из 50 чисел}

Зафиксируем $k$ (обычно $k=30\div 100$) и спроектируем каждое лицо
$x$ на эти $k$ собственных лиц:
\begin{equation}\label{eq:eigenface-proj}
\alpha_i \;=\; v_i^{\!\top}(x-\bar x),\qquad i=1,\ldots,k.
\end{equation}
Получили вектор коэффициентов $\alpha=(\alpha_1,\ldots,\alpha_k)\in\R^{k}$.
Это и есть «лицо человека» с точки зрения алгоритма — \emph{всего
$50$ чисел} вместо $4096$.

\subsection{Шаг 4. Распознавание и реконструкция}

\paragraph{Распознавание.} Чтобы узнать, кто на новой фотографии $y$,
вычислим её коэффициенты $\beta$ по~\eqref{eq:eigenface-proj} и
найдём в базе ближайшего соседа: $\hat j = \argmin_j \norm{\beta-\alpha^{(j)}}$.

\paragraph{Реконструкция.} Зная $\alpha$, можно восстановить лицо:
\[
x \;\approx\; \bar x + \sum_{i=1}^{k}\alpha_i\,v_i.
\]
Рис.~\ref{fig:eigenfaces-recon} показывает, как качество
реконструкции растёт с $k$.

\begin{figure}[ht]
\centering
\includegraphics[width=\linewidth]{eigenfaces_reconstruction.pdf}
\caption{Реконструкция одной из фотографий датасета через первые $k$
eigenfaces. При $k=10$ узнаётся «общий тип» лица; при $k=50$ — уже
конкретный человек; при $k=150$ — почти оригинал.}
\label{fig:eigenfaces-recon}
\end{figure}

\subsection{Реализация на Python}

\begin{lstlisting}[label={lst:eigenfaces}]
import numpy as np
import matplotlib.pyplot as plt
from sklearn.datasets import fetch_olivetti_faces

# 1. Load: 400 faces, 64x64 each
faces = fetch_olivetti_faces().images        # (400, 64, 64)
X = faces.reshape(400, -1)                   # (400, 4096)

# 2. Center: subtract the mean face
x_bar = X.mean(axis=0)
X_centered = X - x_bar

# 3. SVD => eigenfaces are rows of Vt
U, S, Vt = np.linalg.svd(X_centered, full_matrices=False)

# 4. Project all faces onto top-k eigenfaces
k = 50
alpha = X_centered @ Vt[:k].T                # (400, 50)

# 5. Recognise: nearest neighbour in alpha-space
def recognise(y, k=50):
    y_centered = y.flatten() - x_bar
    beta = Vt[:k] @ y_centered
    distances = np.linalg.norm(alpha - beta, axis=1)
    return distances.argmin()

# 6. Reconstruct a face from k coefficients
def reconstruct(idx, k=50):
    return x_bar + Vt[:k].T @ alpha[idx, :k]
\end{lstlisting}

\begin{interestbox}
Метод eigenfaces — родоначальник всех современных систем распознавания
лиц, включая ту, что разблокирует ваш смартфон. Современные системы
заменяют SVD на свёрточную нейросеть (см.~§\,3.3.7), но базовая
идея — \emph{сжать лицо в короткий вектор, в котором близкие лица
оказываются рядом} — осталась той же. Эти короткие векторы сегодня
называют \term{эмбеддингами}.
\end{interestbox}

% =====================================================================
\section{Eigencats: то же самое для котиков\texorpdfstring{$^\star$}{*}}
\label{sec:linalg-eigencats}

Кошки — отличный материал для проверки гипотезы, что метод eigenfaces
работает на любых однородных объектах, не только на лицах. Возьмём
датасет в $\sim 2000$ котиков\footnote{Например,
\engterm{Cats Faces Dataset} на Kaggle: фотографии $64\times 64$.
Готовые ноутбуки выложены в репозитории \href{https://github.com/MerkulovDaniil/nla360}{nla360}.}
и проведём тот же эксперимент.

\begin{figure}[ht]
\centering
\includegraphics[width=0.85\linewidth]{eigencats_components.pdf}
\caption{Первые 8 eigencats. Видны общие «коты»: положение ушей,
форма морды, направление взгляда. Каждый конкретный кот собирается
как линейная комбинация этих универсальных «прото-котов».}
\label{fig:eigencats}
\end{figure}

Качественная картина повторяется: первое eigencat ловит освещение,
второе — наклон головы, дальше идут вариации формы морды и ушей.
В сжатом виде каждый кот превращается в $\sim 50$ чисел.

\begin{notebox}[title={Что это нам говорит о ИИ}]
\textbf{Пространство «реальных» объектов всегда узкое.} Будь это
лица людей, морды котов, рукописные цифры (датасет MNIST) или
фотографии природы — реальные данные занимают в исходном
многомерном пространстве \emph{маленькое} подпространство.
Распознавание = найти это подпространство и работать в нём.

\textbf{Эта идея — фундамент современных автокодировщиков и
эмбеддингов} (см.~§\,3.4.3--3.4.6). SVD — её самая ранняя и самая
понятная реализация.
\end{notebox}

% =====================================================================
\section{PageRank: как Google нашёл порядок в Интернете}
\label{sec:linalg-pagerank}

\subsection{Задача и история}

К 1996~г. поисковики типа \engterm{AltaVista} и \engterm{Yahoo!}
ранжировали страницы по тому, сколько раз искомое слово встречается
в тексте. Это легко обманывалось: достаточно было набить страницу
популярными словами белым цветом по белому фону, и она попадала
в топ выдачи.

Лоуренс Пейдж и Сергей Брин, тогда аспиранты Стэнфорда, предложили
другой принцип ранжирования: \emph{страница важна, если на неё
ссылаются другие важные страницы}.

\subsection{Граф ссылок как матрица}

Обозначим $N$ — число всех страниц Интернета (на 1998~г. — около
$2{,}4\cdot 10^{8}$, сегодня — $\sim 10^{11}$). Построим матрицу
$P\in\R^{N\times N}$:
\[
P_{ij}\;=\;\begin{cases}
\dfrac{1}{c_j}, & \text{если со страницы $j$ есть ссылка на $i$,}\\
0, & \text{иначе,}
\end{cases}
\]
где $c_j$ — \emph{полное число исходящих ссылок} со страницы $j$.
Содержательно: $P_{ij}$ — это вероятность того, что пользователь,
сидящий на странице $j$ и случайно ткнувший на одну из её ссылок,
окажется на странице $i$. Поэтому $P$ называется матрицей переходов
\term{случайного блуждания} по графу Интернета.

\subsection{Стационарное распределение = собственный вектор}

Пусть $r\in\R^{N}$ — вектор «важностей»: $r_i$ — степень важности
страницы $i$. Принцип Пейджа–Брина формализуется системой:
\begin{equation}\label{eq:pr-fixed}
r\;=\;P r.
\end{equation}
То есть важность каждой страницы равна сумме важностей тех страниц,
которые на неё ссылаются, делённой на число их исходящих ссылок.

Уравнение~\eqref{eq:pr-fixed} говорит: $r$ — это \emph{собственный
вектор} матрицы $P$ с собственным значением $1$. И в этом — вся суть
PageRank.

\begin{theorem}{Перрон--Фробениус для PageRank}{perron}
Если все элементы матрицы $P$ положительны (или хотя бы граф ссылок
сильно связен и непериодичен), то собственное значение $1$ у $P$
существует, оно \emph{простое} и наибольшее по модулю, а
соответствующий собственный вектор $r$ единственен и имеет
положительные компоненты. Этот $r$ называется \term{стационарным
распределением} случайного блуждания.
\end{theorem}

(Эта теорема — Оскар Перрон, 1907~г., и Георг Фробениус, 1912~г. —
исторически появилась задолго до Интернета, в контексте моделей
популяций.)

\subsection{Демпфирующий множитель и формула Брина--Пейджа}

В реальной матрице Интернета у некоторых страниц нет исходящих ссылок
(\engterm{dangling pages}), а в других группах — циклы. Для таких
страниц соответствующий столбец матрицы переходов заранее заменяют
равномерным распределением по всем страницам: иначе формула $1/c_j$
при $c_j=0$ не определена, а масса случайного блуждания «исчезала» бы
в dangling pages. Чтобы
гарантировать применимость теоремы Перрона--Фробениуса, Пейдж и Брин
ввели \term{демпфирующий множитель} $\beta\in(0,1)$ (обычно
$\beta=0{,}85$):
\begin{equation}\label{eq:pagerank}
\boxed{\;r\;=\;\beta P r + \dfrac{1-\beta}{N}\,\mathbf{1}.\;}
\end{equation}
Здесь $\mathbf{1}=(1,1,\ldots,1)^{\!\top}$. Содержательно:
с вероятностью $\beta\approx 0{,}85$ случайный сёрфер кликает по
ссылке, а с вероятностью $1-\beta\approx 0{,}15$ — телепортируется
на случайную страницу Интернета.

\subsection{Как находить $r$: степенной метод}

Прямо решать систему $(I-\beta P)r=\frac{1-\beta}{N}\mathbf{1}$
с $N\sim 10^{11}$ невозможно. Но и не надо: вектор $r$ можно
получить простой итерацией.

\begin{algorithmbox}
\textbf{Степенной метод для PageRank.}
\begin{enumerate}\itemsep1pt
  \item Инициализация: $r^{(0)} \leftarrow \tfrac{1}{N}\mathbf{1}$.
  \item Итерация: $r^{(t+1)} \leftarrow \beta\,P\,r^{(t)} +
        \dfrac{1-\beta}{N}\,\mathbf{1}$.
  \item Остановка: $\norm{r^{(t+1)}-r^{(t)}} < \varepsilon$.
\end{enumerate}
\end{algorithmbox}

Сходимость гарантируется теоремой Перрона--Фробениуса:
$\norm{r^{(t)}-r}\leq C\cdot\beta^{t}$, то есть погрешность убывает
геометрически со знаменателем $\beta=0{,}85$. На практике хватает
$50\div 100$ итераций.

\subsection{Игрушечный пример: 6 страниц}

Возьмём маленький «Интернет» из шести страниц (рис.~\ref{fig:pagerank-graph})
и посчитаем для него PageRank.

\begin{figure}[ht]
\centering
\begin{tikzpicture}[scale=1.0, every node/.style={circle, draw, accent,
              fill=accent!10, minimum size=8mm, font=\small}]
  \node (A) at (0,2)    {A};
  \node (B) at (2.5,3)  {B};
  \node (C) at (5,2)    {C};
  \node (D) at (5,-0.5) {D};
  \node (E) at (2.5,-1.5){E};
  \node (F) at (0,-0.5) {F};
  \begin{scope}[->, thick, accent2]
    \draw (A) -- (B);
    \draw (A) -- (C);
    \draw (B) -- (C);
    \draw (C) -- (A);
    \draw (D) -- (C);
    \draw (D) -- (A);
    \draw (E) -- (D);
    \draw (E) -- (B);
    \draw (E) -- (F);
    \draw (F) -- (E);
    \draw (F) -- (B);
  \end{scope}
\end{tikzpicture}
\caption{Игрушечный «Интернет» из 6 страниц со ссылками. Видно,
например, что на страницу C ссылаются три другие, а на F — только
одна; будет ли C лидером по PageRank?}
\label{fig:pagerank-graph}
\end{figure}

\begin{lstlisting}[label={lst:pagerank}]
import numpy as np

# Adjacency list -> column-stochastic transition matrix P
links = {'A': ['B','C'], 'B': ['C'], 'C': ['A'],
         'D': ['C','A'], 'E': ['D','B','F'], 'F': ['E','B']}
pages = list(links); n = len(pages); idx = {p: i for i, p in enumerate(pages)}

P = np.zeros((n, n))
for j, p in enumerate(pages):
    out = links[p]
    for q in out:
        P[idx[q], j] = 1.0 / len(out)

# PageRank via power iteration
beta, eps = 0.85, 1e-10
r = np.ones(n) / n
for t in range(200):
    r_new = beta * P @ r + (1 - beta) / n
    if np.linalg.norm(r_new - r, 1) < eps: break
    r = r_new

for p, ri in sorted(zip(pages, r), key=lambda t: -t[1]):
    print(f'{p}: {ri:.4f}')
\end{lstlisting}

\noindent Выдача программы:
\[
\text{C: }0{,}295,\quad
\text{A: }0{,}252,\quad
\text{B: }0{,}175,\quad
\text{E: }0{,}098,\quad
\text{F: }0{,}092,\quad
\text{D: }0{,}087.
\]
Лидер — \texttt{C}, хотя на неё формально ссылаются столько же
страниц, сколько на \texttt{A}: PageRank учёл, что одна из ссылок
на \texttt{C} приходит с уже-важной страницы \texttt{B}.

\begin{historybox}
Алгоритм PageRank был опубликован в 1998~г. в студенческой работе
\emph{Sergey Brin, Larry Page}, «\textit{The PageRank Citation Ranking:
Bringing Order to the Web}». В том же году Брин и Пейдж основали
Google. К 2004~г. компания вышла на IPO, а к 2010-м стала одной из
крупнейших в мире. Сам Стэнфорд получил $\sim 336$ миллионов долларов
от лицензирования патента --- который изначально просто описывал
итерационный метод для собственного вектора.

Идея использовать собственный вектор графа для измерения «важности»
позднее была обобщена и применена в наукометрии (\engterm{Eigenfactor}
для ранжирования научных журналов), в рекомендательных системах,
а также в современных нейросетях для графов (\engterm{Graph Neural
Networks}).
\end{historybox}

% =====================================================================
\section{Когда нейросеть устойчива: оценки Липшицевости\texorpdfstring{$^\star$}{*}}
\label{sec:linalg-lipschitz}

\subsection{Сюжет: когда панда становится гиббоном}

В 2014~г. Иэн Гудфеллоу и соавторы опубликовали ставшую знаменитой
картинку (рис.~\ref{fig:adv-attack}): обычная фотография панды,
которую современная нейросеть уверенно опознаёт как «panda» с
вероятностью $57{,}7\%$. К каждому пикселю прибавлена крошечная,
визуально незаметная человеку поправка. На обновлённой картинке
сеть теперь уверенно (с вероятностью $99{,}3\%$) считает, что
видит \emph{гиббона}.

\begin{figure}[ht]
\centering
\includegraphics[width=0.95\linewidth]{adversarial_panda.pdf}
\caption{Adversarial-атака на нейросеть (Goodfellow et al., 2014).
Слева — оригинал, классифицирован как «panda» с уверенностью $57{,}7\%$.
В центре — почти невидимое возмущение, масштабированное для наглядности.
Справа — атакованный пример, классифицируется как «gibbon» с
уверенностью $99{,}3\%$.}
\label{fig:adv-attack}
\end{figure}

Этот эффект называется \term{adversarial-атакой} и долгое время
считался экзотикой. Сейчас понятно: причина в простом числе,
которое сопровождает любую функцию — её \emph{константе Липшица}.

\subsection{Константа Липшица: формальное определение}

\begin{definition}{Константа Липшица функции}{lipschitz}
Функция $f\colon\R^{n}\to\R^{m}$ называется \term{Липшицевой} с
константой $L>0$, если для любых $x,y\in\R^{n}$ выполнено
\begin{equation}\label{eq:lipschitz}
\norm{f(x)-f(y)}\;\leq\;L\,\norm{x-y}.
\end{equation}
Наименьшее $L$, для которого~\eqref{eq:lipschitz} выполнено,
называется \emph{константой Липшица} функции $f$ и обозначается
$\mathrm{Lip}(f)$.
\end{definition}

Содержательно $\mathrm{Lip}(f)$ показывает, во сколько раз может
\emph{усилиться} входное возмущение на выходе функции. Если
$\mathrm{Lip}(f)=L$, то малое возмущение $\delta$ на входе
может вызвать возмущение до $L\cdot\norm{\delta}$ на выходе.

\paragraph{Пример.} Для линейного отображения $f(x)=Wx$ имеем
$f(x)-f(y)=W(x-y)$, откуда $\norm{f(x)-f(y)}\leq \norm{W}_2\cdot\norm{x-y}$,
и константа Липшица — это в точности спектральная норма матрицы $W$:
\begin{equation}\label{eq:lip-linear}
\mathrm{Lip}(Wx)\;=\;\norm{W}_2\;=\;\sigma_1(W).
\end{equation}
Здесь $\sigma_1(W)$ — наибольшее сингулярное число матрицы $W$.
Вот \emph{здесь} в нашу историю входит SVD из §\,\ref{sec:linalg-svd}.

\subsection{Композиция: Липшицева константа цепочки}

\begin{theorem}{Композиция Липшицевых функций}{lip-compose}
Если $f$ и $g$ — Липшицевы с константами $L_f$ и $L_g$, то
их композиция $f\circ g$ — тоже Липшицева, причём
\begin{equation}\label{eq:lip-prod}
\mathrm{Lip}(f\circ g)\;\leq\;\mathrm{Lip}(f)\cdot\mathrm{Lip}(g).
\end{equation}
\end{theorem}

\begin{proof}
$\norm{f(g(x))-f(g(y))}\leq L_f\norm{g(x)-g(y)}\leq L_f L_g\norm{x-y}$. $\square$
\end{proof}

\subsection{Применение к нейросети}

Возьмём типичную нейросеть глубины $K$:
\[
f(x)\;=\;\sigma\bigl(W_K\,\sigma(W_{K-1}\,\sigma(\cdots\sigma(W_1 x))\bigr),
\]
где $W_i$ — матрицы весов, $\sigma$ — поэлементная активация (ReLU,
сигмоида, и т.п.). Для большинства активаций $\mathrm{Lip}(\sigma)\leq 1$
(для ReLU и $\tanh$ она равна $1$, а для логистической сигмоиды — $1/4$). По
теореме~\ref{thm:lip-compose}:
\begin{equation}\label{eq:nn-lip}
\boxed{\;\mathrm{Lip}(f)\;\leq\;\prod_{i=1}^{K}\norm{W_i}_2\;}
\end{equation}

Это и есть тот «лишний рычаг», который объясняет уязвимость
нейросетей. Если $\norm{W_i}_2 \approx 2$ для каждого слоя из $K=20$
слоёв, то возмущение на входе \emph{потенциально} усиливается
в $2^{20}\approx 10^{6}$ раз. Невидимая поправка $\delta$ с
$\norm{\delta}=10^{-3}$ может вызвать изменение выхода до $10^{3}$ —
этого хватит, чтобы перебросить классификатор из одного класса в другой.

\subsection{Как защищаться: спектральная нормализация}

Естественная стратегия: после каждого шага обучения нормализовать
матрицу $W_i$ так, чтобы $\norm{W_i}_2 = 1$. Это и есть метод
\term{Spectral Normalization} (Миято и соавт., 2018), сделавший
обучение генеративных нейросетей (GAN) намного стабильнее.

Технически: на каждом шаге считаем $\sigma_1(W)$ (с помощью одной
итерации степенного метода — почти бесплатно!) и заменяем
$W \mapsto W/\sigma_1(W)$.

\begin{lstlisting}[label={lst:spec-norm}]
def spectral_normalize(W, n_iter=1):
    """One step of power iteration to estimate sigma_1(W) and rescale."""
    u = np.random.randn(W.shape[0])
    for _ in range(n_iter):
        v = W.T @ u; v /= np.linalg.norm(v)
        u = W @ v;   u /= np.linalg.norm(u)
    sigma_1 = u @ W @ v
    return W / sigma_1
\end{lstlisting}

\begin{interestbox}
\textbf{Связь с устойчивостью в физике.} Похожая ситуация хорошо
известна в теории динамических систем: если константа Липшица
отображения $> 1$, малое возмущение начального условия экспоненциально
нарастает (классический эффект бабочки в погодных моделях).
Spectral Normalization можно рассматривать как искусственное
гашение этого эффекта.
\end{interestbox}

% =====================================================================
\section*{Подведение итогов главы}
\addcontentsline{toc}{section}{Подведение итогов главы}

\begin{itemize}\itemsep1pt
  \item Матрица — это \emph{универсальный носитель данных} для ИИ:
        картинки, тексты, графы, веса нейросети.
  \item \textbf{Числа с плавающей точкой} (§\ref{sec:linalg-floats}) —
        вычисления на компьютере всегда приближённые с относительной
        ошибкой $\sim\varepsilon_{\text{маш}}\approx 10^{-16}$.
        Самая опасная ошибка — \emph{катастрофическое сокращение} при
        вычитании близких чисел; её часто можно устранить, переписав
        формулу алгебраически эквивалентным способом.
  \item \textbf{Сингулярное разложение} $A=U\Sigma V^{\!\top}$
        (§\ref{sec:linalg-svd}) представляет любую матрицу как
        «поворот → растяжение по осям → поворот». Сингулярные числа
        $\sigma_i$ — это коэффициенты растяжения.
  \item \textbf{Теорема Эккарта--Янга} (§\ref{sec:linalg-eckart}):
        $A_k=\sum_{i\leq k}\sigma_i u_i v_i^{\!\top}$ — лучшее
        приближение ранга $k$ матрицы $A$. На этой теореме
        стоит сжатие изображений, метод eigenfaces и современный
        LoRA для дообучения больших языковых моделей.
  \item \textbf{Eigenfaces / eigencats}
        (§\ref{sec:linalg-eigenfaces}--\ref{sec:linalg-eigencats}):
        реальные данные (лица, котики, цифры) живут в очень узком
        подпространстве своего номинального многомерного пространства.
        Найти это подпространство = применить SVD к выборке.
  \item \textbf{PageRank} (§\ref{sec:linalg-pagerank}): «важность»
        страницы Интернета = компонента собственного вектора матрицы
        переходов случайного блуждания. Алгоритм находит его простой
        итерацией; теорема Перрона--Фробениуса гарантирует сходимость.
  \item \textbf{Константа Липшица} нейросети
        (§\ref{sec:linalg-lipschitz}) растёт мультипликативно от
        слоя к слою; именно поэтому глубокие сети уязвимы к
        adversarial-атакам. Спектральная нормализация ограничивает
        $\norm{W_i}_2=1$ и стабилизирует обучение.
\end{itemize}

% =====================================================================
\section*{Что почитать дальше}
\addcontentsline{toc}{section}{Что почитать дальше}

\begin{itemize}\itemsep1pt
  \item \emph{Е.\,Е.~Тыртышников.} Методы численного анализа. М.: Академия, 2007. ---
        классический российский учебник по вычислительной линейной
        алгебре университетского уровня.
  \item \emph{G.~Strang.} Linear Algebra and Learning from Data. Wellesley-Cambridge Press, 2019. ---
        свежая книга от автора знаменитого MIT-курса 18.06, написанная
        с прицелом на data science.
  \item \emph{L.\,N.~Trefethen, D.~Bau.} Numerical Linear Algebra. SIAM, 1997. ---
        классика, с очень понятным изложением SVD и итерационных методов.
  \item \emph{Видеокурс «Numerical Linear Algebra»}, И.\,В.~Оселедец, Сколтех
        (\href{https://github.com/MerkulovDaniil/nla360}{github.com/MerkulovDaniil/nla360}).
        Лекции, семинары, ноутбуки --- из которых, в частности,
        и взяты eigencats для этой главы.
  \item \emph{S.~Brin, L.~Page.} The Anatomy of a Large-Scale
        Hypertextual Web Search Engine. WWW Conference, 1998. ---
        оригинальная статья про PageRank, читается легко.
  \item \emph{I.~Goodfellow et al.} Explaining and Harnessing
        Adversarial Examples. ICLR, 2015. ---
        статья про панду-гиббона, из которой взят рис.~\ref{fig:adv-attack}.
\end{itemize}

% =====================================================================
\section*{Большой итоговый проект}
\addcontentsline{toc}{section}{Большой итоговый проект}

\textbf{Постройте свою рекомендательную систему для фильмов.}

Возьмите датасет MovieLens-100K (открытый, $100\,000$ оценок,
$\sim 1000$ пользователей $\times$ $\sim 1700$ фильмов). Постройте
матрицу «пользователь--фильм» $R$, где $R_{ij}\in\{1,\ldots,5\}$ —
оценка, поставленная пользователем $i$ фильму $j$, а $0$ — «не смотрел».

\textbf{Важно:} нули в матрице $R$ означают \emph{пропущенные} оценки,
а не настоящие нулевые рейтинги. Прямое SVD по матрице с нулями — лишь
грубый учебный baseline: непросмотренные фильмы тогда воспринимаются как
фильмы с очень низкой оценкой. Более корректный вариант — обучать
малоранговую модель только на известных оценках либо предварительно
заполнить пропуски средними значениями пользователя/фильма и затем
сравнить результат.

\begin{enumerate}\itemsep1pt
  \item Постройте малоранговое приближение $R\approx U\Sigma V^{\!\top}$
        с $k=20$ компонентами через SVD.
  \item Объясните содержательно, что такое строки $U$ (профили
        пользователей) и столбцы $V$ (профили фильмов) в этом
        приближении. Сколько чисел теперь хранит модель вместо
        $1000\cdot 1700$?
  \item Для нескольких реальных пользователей выпишите топ-5 фильмов,
        которые рекомендация модели предлагает посмотреть. Совпадают
        ли они с любимыми жанрами пользователя?
  \item Найдите два «противоположных» фильма с наибольшим расстоянием
        между их столбцами $V$ и подумайте, можно ли увидеть содержательный
        смысл в этой разделяющей оси (боевик vs. мелодрама? старое vs. новое?).
  \item ($\star$) Замените SVD на алгоритм PageRank, построив граф
        «фильмы, которые часто оценивают одинаково». Получится ли тот
        же топ?
\end{enumerate}

% =====================================================================
\begin{tasksbox}
\begin{enumerate}\itemsep1pt
  \item Сложите в Python $10\,000\,000$ копий числа $0{,}1$. Сколько
        должно получиться? Сколько получается на самом деле? Почему?
  \item Вычислите $f(x)=\sqrt{x+1}-\sqrt{x}$ при $x=10^{16}$
        напрямую и через формулу $f(x)=\frac{1}{\sqrt{x+1}+\sqrt{x}}$.
        Сравните ответы. Какой устойчив, почему?
  \item Найдите вручную SVD матрицы
        $A=\bigl(\begin{smallmatrix}3 & 0\\0 & 4\end{smallmatrix}\bigr)$
        и $A=\bigl(\begin{smallmatrix}0 & 1\\-1 & 0\end{smallmatrix}\bigr)$.
        Что в обоих случаях геометрически делают $U, \Sigma, V$?
  \item Возьмите свою фотографию $300\times 300$ в градациях серого.
        Сделайте сжатие через SVD с рангами $k=5, 20, 50, 100$. При
        каком $k$ вы перестаёте узнавать самого себя?
  \item Для игрушечного графа из рис.~\ref{fig:pagerank-graph}
        прокрутите 5 итераций степенного метода \emph{вручную},
        начав с $r^{(0)}=(\tfrac16,\ldots,\tfrac16)^{\!\top}$. Совпали ли
        ваши значения с программными после $5$ шагов? а после $20$?
  \item $^{\star}$ Покажите, что для функции $\sigma(x)=\max(0,x)$
        (ReLU) константа Липшица равна $1$. Тот же вопрос для
        сигмоиды $\sigma(x)=1/(1+e^{-x})$: чему равна $\mathrm{Lip}(\sigma)$
        и в какой точке производная максимальна?
  \item $^{\star}$ Найдите спектральную норму матрицы
        $W=\bigl(\begin{smallmatrix}1 & 2\\2 & 1\end{smallmatrix}\bigr)$
        двумя способами: (а) через определение
        $\norm{W}_2=\max_{\norm{x}=1}\norm{Wx}$ (с использованием
        собственных значений $W^{\!\top}W$); (б) одной итерацией
        степенного метода (с любым стартом). Совпали ли ответы?
\end{enumerate}
\end{tasksbox}
      % <— сюда
%
%  Все используемые окружения определены в structure.tex и
%  chapters/extras.tex; новых пакетов не требуется.
% =====================================================================

\chapter{Вычислительная линейная алгебра}
\label{ch:linalg}

\emph{Если метод Ньютона из главы~\ref{ch:opt} был одним
маленьким уравнением, упрямо повторяющимся, то линейная алгебра —
это язык, на котором говорят сразу о тысячах уравнений и миллионах
чисел. Этой главой мы покажем, что многие фокусы современного
искусственного интеллекта — от того, как нейросеть «узнаёт» лицо
человека, до того, как поисковик ранжирует страницы Интернета, —
держатся всего на нескольких разложениях матриц. Мы посмотрим
четыре истории: лица, котиков, Интернет и устойчивость нейросетей.}

\section{Зачем линейной алгебре отдельная глава}
\label{sec:linalg-why}

Вы уже встречались с матрицами в школьном курсе — как с прямоугольными
таблицами чисел. В современном ИИ матрица — это \term{универсальный
носитель данных}.

\begin{itemize}\itemsep1pt
  \item Чёрно-белая фотография $1024\times 1024$ — это матрица
        размера $1024\times 1024$, где каждое число — яркость пикселя
        от $0$ (чёрный) до $255$ (белый). Цветная — три такие матрицы.
  \item Текст из $N$ слов словаря — это вектор частот длины $N$
        (модель «мешка слов»), и набор из $M$ документов — матрица
        $M\times N$.
  \item Социальная сеть из $N$ пользователей — матрица $N\times N$,
        где элемент $(i,j)$ равен $1$, если $i$ подписан на $j$.
  \item Веса одного слоя нейросети — матрица $W$, через которую
        входной сигнал умножается, чтобы получить выходной.
\end{itemize}

С такими матрицами надо уметь \emph{что-то делать}: сжимать,
сравнивать, упорядочивать, обновлять. И почти всегда оказывается,
что нужная операция выражается через одно и то же фундаментальное
разложение матрицы — \term{сингулярное}. Мы изучим его и применим
к четырём задачам:

\begin{enumerate}\itemsep1pt
  \item \textbf{Узнать лицо} по фотографии — \emph{eigenfaces}
        (§\ref{sec:linalg-eigenfaces}).
  \item \textbf{Сжать пространство признаков} — \emph{eigencats} и
        теорема Эккарта–Янга (§\ref{sec:linalg-eckart},
        §\ref{sec:linalg-eigencats}).
  \item \textbf{Упорядочить весь Интернет} — \emph{PageRank}
        (§\ref{sec:linalg-pagerank}).
  \item \textbf{Понять, когда нейросеть устойчива} — \emph{оценки
        Липшицевости} (§\ref{sec:linalg-lipschitz}).
\end{enumerate}

Прежде чем перейти к сюжетам, разберёмся с одной \emph{очень} земной
проблемой: с тем, как компьютер хранит числа.

% =====================================================================
\section{Когда арифметика лжёт: числа с плавающей точкой}
\label{sec:linalg-floats}

Откройте Python и наберите:
\begin{lstlisting}[label={lst:float-paradox}]
>>> 0.1 + 0.2
0.30000000000000004
>>> 0.1 + 0.2 == 0.3
False
\end{lstlisting}

Это не баг Python. То же самое выдаст любой современный язык —
C, C++, Java, JavaScript. И этому есть фундаментальная причина:
\term{двоичная система счисления} не умеет точно представить
десятичную дробь $0{,}1$.

\subsection{Двоичные дроби и почему $0{,}1$ — это «$0{,}333\ldots$»}

В десятичной системе $1/3 = 0{,}3333\ldots$ — бесконечная дробь.
То же самое в двоичной системе случается с $0{,}1$:
\[
0{,}1_{10} \;=\; 0{,}0001\,1001\,1001\,1001\ldots_{2}
\]
Дробь периодическая. Любой реальный компьютер хранит её, оборвав
после конечного числа разрядов (обычно 52 двоичных знака после
запятой — стандарт \engterm{IEEE~754, double precision}). Поэтому
число, которое программист записывает как \texttt{0.1}, на самом
деле равно
\[
0{,}1\;\approx\;0{,}100\,000\,000\,000\,000\,005\,551\,115\ldots
\]
Когда мы складываем такое «искажённое $0{,}1$» с искажённым $0{,}2$,
ошибки усиливаются — и получается тот самый «лишний хвостик»
$4\cdot 10^{-17}$, который мы видели в листинге.

\begin{notebox}[title={Машинная точность $\varepsilon_{\text{маш}}$}]
\term{Машинной точностью} называют наименьшее положительное число
$\varepsilon$, такое что в компьютерной арифметике
$1+\varepsilon \neq 1$. Для стандарта \texttt{double} оно равно
$\varepsilon_{\text{маш}} = 2^{-52}\approx 2{,}22\cdot 10^{-16}$.

Это означает: какое бы хорошее число $x$ вы ни записали, оно
известно компьютеру лишь с относительной точностью
$\sim\varepsilon_{\text{маш}}$.
\end{notebox}

\subsection{Когда это становится опасным: катастрофическое сокращение}

Маленькая ошибка в $17$-м знаке — мелочь? Не всегда. Рассмотрим
безобидное на вид вычисление:
\[
g(x) \;=\; \dfrac{1-\cos x}{x^{2}}\qquad\text{при }x=10^{-8}.
\]
По формуле Тейлора $1-\cos x \approx x^{2}/2$, так что
$g(10^{-8})\approx 0{,}5$. Что выдаст компьютер?

\begin{lstlisting}[label={lst:cancel}]
>>> from math import cos
>>> x = 1e-8
>>> (1 - cos(x)) / x**2
0.0           # вместо 0.5!
\end{lstlisting}

Что произошло? Число $\cos(10^{-8})\approx 1-5\cdot 10^{-17}$ так
близко к единице, что компьютер просто хранит его как
\emph{ровно единицу} (поправка $5\cdot 10^{-17}$ меньше
$\varepsilon_{\text{маш}}$). Разность $1-\cos x$ обнулилась, и весь
ответ обнулился вместе с ней.

Это явление называется \term{катастрофическим сокращением}
(\engterm{catastrophic cancellation}): когда из двух близких чисел
вычитают и теряют все значащие цифры разом.

\begin{example}{Спасение через переписывание формулы}{cancel-fix}
Воспользуемся тождеством $1-\cos x = 2\sin^{2}(x/2)$:
\[
g(x) \;=\; \dfrac{2\sin^{2}(x/2)}{x^{2}}\;=\;\dfrac12\!\left(\dfrac{\sin(x/2)}{x/2}\right)^{2}.
\]
Внутри теперь нет вычитания близких чисел, и Python выдаёт
\texttt{0.49999999999999983} — погрешность лишь в последнем знаке.
\end{example}

\begin{importantbox}
\textbf{Главный практический вывод.} Если ваш численный алгоритм
выдаёт странный ответ, первая гипотеза, которую следует проверить, —
это не «бага в библиотеке», а \emph{катастрофическое сокращение} в
вашей формуле. Часто его можно убрать, переписав формулу
алгебраически эквивалентным, но численно устойчивым способом.
\end{importantbox}

\begin{historybox}
Стандарт \engterm{IEEE~754} был принят в 1985~г., в значительной
мере благодаря Уильяму Кэхэну (\engterm{William Kahan}, премия
Тьюринга 1989~г.). До этого каждый производитель процессоров делал
плавающую арифметику по-своему, и одна и та же программа на разных
машинах могла давать заметно разные ответы. Кэхэн также является
автором знаменитого \term{алгоритма компенсированного суммирования},
который позволяет складывать миллионы чисел с почти двойной
точностью.
\end{historybox}

% =====================================================================
\section{Сингулярное разложение SVD: главный фокус линейной алгебры}
\label{sec:linalg-svd}

\subsection{Геометрическая идея}

Возьмём любую матрицу $A$ размера $m\times n$ и применим её к
единичному кругу в $\R^{n}$ (то есть к множеству всех векторов $x$
с $\norm{x}=1$, где $\norm{\cdot}$ — обычная евклидова длина).
Что получится в $\R^{m}$?

\begin{theorem}{Сингулярное разложение (SVD)}{svd}
Для любой вещественной матрицы $A$ размера $m\times n$ существует
разложение
\begin{equation}\label{eq:svd}
\boxed{\;A \;=\; U\,\Sigma\,V^{\!\top}\;}
\end{equation}
где
\begin{itemize}\itemsep1pt
  \item $U$ — \emph{ортогональная} матрица $m\times m$ (поворот в $\R^{m}$);
  \item $V$ — ортогональная матрица $n\times n$ (поворот в $\R^{n}$);
  \item $\Sigma$ — «диагональная» матрица $m\times n$ с неотрицательными
        числами $\sigma_1\geq \sigma_2\geq\ldots\geq \sigma_r > 0$ на
        диагонали (и нулями ниже), где $r=\mathrm{rank}\,A$.
\end{itemize}
Числа $\sigma_i$ называются \term{сингулярными числами} матрицы~$A$.
\end{theorem}

Содержательно \eqref{eq:svd} говорит: \emph{любая} линейная операция —
это три простых действия подряд:
\[
\underbrace{V^{\!\top}}_{\text{поворот в исходном пространстве}}
\;\to\;
\underbrace{\Sigma}_{\text{растяжение по координатным осям}}
\;\to\;
\underbrace{U}_{\text{поворот в новом пространстве}}.
\]

Геометрически это значит: образ единичной сферы под действием $A$ —
это \term{эллипсоид}, у которого полуоси равны $\sigma_1,\sigma_2,\ldots$
(см.~рис.~\ref{fig:svd-ellipse}).

\begin{figure}[ht]
\centering
\begin{tikzpicture}[scale=1.05, line cap=round]
  % Левая картинка: единичный круг
  \begin{scope}[xshift=-3.6cm]
    \draw[->, gray!60] (-1.6,0) -- (1.6,0);
    \draw[->, gray!60] (0,-1.6) -- (0,1.6);
    \draw[thick, accent] (0,0) circle (1);
    \draw[->, very thick, accent2] (0,0) -- (1,0) node[below right] {$v_1$};
    \draw[->, very thick, accent2] (0,0) -- (0,1) node[above left] {$v_2$};
    \node[below] at (0,-1.85) {$\R^{n}$: $\norm{x}=1$};
  \end{scope}
  % Стрелка A
  \draw[->, very thick, accent] (-1.5,0) -- (1.5,0)
        node[midway, above, font=\large] {$A$};
  % Правая картинка: эллипс
  \begin{scope}[xshift=3.6cm, rotate=20]
    \draw[->, gray!60, rotate=-20] (-2.4,0) -- (2.4,0);
    \draw[->, gray!60, rotate=-20] (0,-1.4) -- (0,1.4);
    \draw[thick, accent] (0,0) ellipse (2 and 0.7);
    \draw[->, very thick, accent2] (0,0) -- (2,0) node[above] {$\sigma_1 u_1$};
    \draw[->, very thick, accent2] (0,0) -- (0,0.7) node[above left] {$\sigma_2 u_2$};
  \end{scope}
  \node[below] at (3.6,-1.85) {$\R^{m}$: $\{Ax:\norm{x}=1\}$};
\end{tikzpicture}
\caption{SVD геометрически. Матрица $A$ переводит единичный круг
в эллипс, полуоси которого равны сингулярным числам $\sigma_1,\sigma_2$,
а направления полуосей — столбцы матрицы $U$ (левые сингулярные
векторы). Векторы $v_1,v_2$ (столбцы $V$, правые сингулярные
векторы) — это те направления в исходном пространстве, которые
$A$ переводит в полуоси эллипса.}
\label{fig:svd-ellipse}
\end{figure}

\subsection{Связь с собственными значениями}

Сингулярные числа — это, по сути, переодетые собственные значения.
А именно:
\begin{equation}\label{eq:svd-eig}
A^{\!\top} A \;=\; (U\Sigma V^{\!\top})^{\!\top}(U\Sigma V^{\!\top})
\;=\; V\,\Sigma^{\!\top}\Sigma\,V^{\!\top}.
\end{equation}
Матрица $\Sigma^{\!\top}\Sigma$ — диагональная, на её диагонали
стоят $\sigma_1^{2},\sigma_2^{2},\ldots$. Значит, $\sigma_i^{2}$ —
это \emph{собственные значения} симметричной матрицы $A^{\!\top}A$,
а столбцы $V$ — её собственные векторы. Эта формула важна нам
по двум причинам: (1) она обосновывает существование SVD (через
известную спектральную теорему для симметричных матриц),
(2) она будет рабочим инструментом в § \ref{sec:linalg-eigenfaces},
когда мы будем вычислять собственные лица.

\begin{historybox}
Сингулярное разложение независимо открыли итальянский математик
\emph{Эудженио Бельтрами} (1873~г.) и француз \emph{Камилл Жордан}
(1874~г.). В современном виде, как разложение прямоугольной матрицы,
его сформулировал \emph{Карл Эккарт} с \emph{Гейлом Янгом} в 1936~г.
Численно надёжный алгоритм вычисления SVD появился лишь в 1965~г.
(Голуб и Кахан) — и стал одним из ключевых алгоритмов XX века. Сегодня
он реализован в любой математической библиотеке: \texttt{numpy.linalg.svd},
\texttt{scipy.linalg.svd}, \texttt{torch.linalg.svd}.
\end{historybox}

% =====================================================================
\section{Лучшее малоранговое приближение: теорема Эккарта--Янга}
\label{sec:linalg-eckart}

Самое замечательное свойство SVD состоит в том, что оно позволяет
\emph{ужимать} матрицу, выбрасывая всё «несущественное».

\subsection{Разложение по слагаемым ранга один}

Перепишем~\eqref{eq:svd} по-другому. Обозначим столбцы матрицы
$U$ через $u_1,\ldots,u_m$ (это \emph{левые сингулярные векторы}),
а столбцы $V$ — через $v_1,\ldots,v_n$ (\emph{правые сингулярные
векторы}). Тогда $A$ есть сумма $r$ слагаемых, каждое из которых
имеет ранг один:
\begin{equation}\label{eq:svd-sum}
A \;=\; \sum_{i=1}^{r}\sigma_i\,u_i v_i^{\!\top}.
\end{equation}

Каждое слагаемое $u_i v_i^{\!\top}$ — это «вертикальный вектор
умножить на горизонтальный», и оно само по себе является матрицей
ранга один (рис.~\ref{fig:rank1}).

\begin{figure}[ht]
\centering
\begin{tikzpicture}[scale=0.85]
  % колонка u
  \draw (0,0) rectangle (0.5,2.5);
  \node at (0.25,2.8) {$u$};
  \node at (0.25,-0.4) {\scriptsize $m\times 1$};
  \node at (1.2,1.25) {\Large $\times$};
  % строка v^T
  \draw (1.7,1.0) rectangle (4.2,1.5);
  \node at (2.95,1.7) {$v^{\!\top}$};
  \node at (2.95,0.5) {\scriptsize $1\times n$};
  \node at (4.8,1.25) {\Large $=$};
  % результат
  \draw[fill=accent!10] (5.3,0) rectangle (7.8,2.5);
  \node at (6.55,2.8) {$u v^{\!\top}$};
  \node at (6.55,-0.4) {\scriptsize $m\times n$, ранг $1$};
\end{tikzpicture}
\caption{Матрица ранга один: каждый столбец результата
пропорционален $u$, каждая строка — пропорциональна $v^{\!\top}$.
Чтобы её хранить, достаточно $m+n$ чисел вместо $mn$.}
\label{fig:rank1}
\end{figure}

Если оборвать сумму~\eqref{eq:svd-sum} на $k$-м слагаемом ($k<r$),
получим приближение:
\begin{equation}\label{eq:lowrank}
A_k \;=\; \sum_{i=1}^{k}\sigma_i\,u_i v_i^{\!\top}.
\end{equation}
Это матрица того же размера $m\times n$, но ранга всего лишь $k$.

\subsection{Теорема Эккарта--Янга: $A_k$ — наилучшее приближение}

\begin{theorem}{Эккарт--Янг--Мирский (1936/1960)}{eckart-young}
Среди всех матриц $B$ размера $m\times n$ ранга не более $k$
матрица $A_k$ из~\eqref{eq:lowrank} даёт \emph{наименьшую} ошибку
в смысле спектральной нормы:
\begin{equation}\label{eq:ey}
\min_{\mathrm{rank}\,B\,\leq\,k}\;\norm{A-B}_{2}\;=\;\norm{A-A_k}_{2}\;=\;\sigma_{k+1}.
\end{equation}
Здесь $\norm{M}_2 = \max_{\norm{x}=1}\norm{Mx}=\sigma_1(M)$ —
\term{спектральная норма} матрицы $M$ (равна её наибольшему
сингулярному числу). Утверждение остаётся в силе и для
\emph{фробениусовой} нормы $\norm{M}_F=\sqrt{\sum_{i,j} m_{ij}^{2}}$:
\[
\norm{A-A_k}_F\;=\;\sqrt{\sigma_{k+1}^{2}+\sigma_{k+2}^{2}+\ldots+\sigma_r^{2}}.
\]
\end{theorem}

Смысл теоремы: \emph{если хочется сжать матрицу до ранга $k$, лучшее,
что вы можете сделать, — оставить первые $k$ компонент SVD}.
Ошибка контролируется \emph{отброшенным} сингулярным числом
$\sigma_{k+1}$.

\subsection{Пример: сжатие изображения}

Возьмём фотографию $512\times 512$ в градациях серого. Хранение
«в лоб» — это $512\cdot 512 = 262\,144$ числа. Применим к матрице
фото SVD и оставим только $k$ старших компонент: придётся хранить
\[
k\cdot(512+512)+k\;=\;k\cdot 1025
\]
чисел — по одному столбцу $u_i$, одной строке $v_i^{\!\top}$ и
одному числу $\sigma_i$ на каждое слагаемое. Уже при $k=50$ это
$\sim 51\,250$ чисел — в \emph{пять} раз меньше исходного,
а визуально потеря почти незаметна (рис.~\ref{fig:img-compress}).

\begin{lstlisting}[label={lst:svd-compress}]
import numpy as np
import matplotlib.pyplot as plt
from skimage import data, color

# Load and convert to grayscale
img = color.rgb2gray(data.astronaut())        # shape (512, 512)
U, S, Vt = np.linalg.svd(img, full_matrices=False)

# Reconstruct using only top-k singular components
ranks = [1, 5, 20, 50, 200]
fig, axes = plt.subplots(1, len(ranks)+1, figsize=(14, 3))
axes[0].imshow(img, cmap='gray'); axes[0].set_title('original')
for ax, k in zip(axes[1:], ranks):
    A_k = U[:, :k] @ np.diag(S[:k]) @ Vt[:k, :]
    ax.imshow(A_k, cmap='gray')
    ax.set_title(f'rank {k}')
for ax in axes: ax.axis('off')
plt.tight_layout(); plt.show()
\end{lstlisting}

\begin{figure}[ht]
\centering
\includegraphics[width=\linewidth]{svd_image_compression.pdf}
\caption{Малоранговое приближение портрета астронавта. Уже при
ранге $50$ изображение визуально неотличимо от оригинала; при ранге
$200$ восстановление идеально. По теореме Эккарта--Янга это
\emph{лучшее} приближение фиксированного ранга.}
\label{fig:img-compress}
\end{figure}

\begin{notebox}[title={LoRA: малоранговая адаптация больших моделей}]
В 2021~г. Эдвард Ху с соавторами заметили: когда большую языковую
модель дообучают на новую задачу, \emph{обновление} матриц весов
$\Delta W$ оказывается почти малоранговой матрицей. Поэтому вместо
обучения «полной» $\Delta W$ можно сразу искать её в виде
$\Delta W = U V^{\!\top}$ с маленьким $k$, где $U,V$ —
тонкие прямоугольники.

Это метод \term{LoRA} (\engterm{Low-Rank Adaptation}). Для модели с
175 миллиардами параметров такой приём сокращает объём
дообучаемых весов в тысячи раз — и именно благодаря ему сегодня
можно «дофайнтюнить» большую модель на бытовой видеокарте.

Теоретический фундамент LoRA — та самая теорема
Эккарта--Янга~\ref{thm:eckart-young}, которой почти 90 лет.
\end{notebox}

% =====================================================================
\section{Eigenfaces: как алгоритм узнаёт лицо}
\label{sec:linalg-eigenfaces}

\subsection{Постановка задачи}

В $1991$~г. Мэттью Тёрк и Алекс Пентланд из MIT поставили простой
вопрос: \emph{можно ли обучить компьютер узнавать лицо, не прибегая
ни к каким специальным «лицевым» признакам — носу, глазам, контурам, —
а просто работая с матрицей пикселей?} Их ответ — метод
\term{собственных лиц} (\engterm{eigenfaces}) — стал классикой и
открыл целую область computer vision.

Идея: каждая фотография лица — это точка в очень многомерном
пространстве. Например, лицо $64\times 64$ — это вектор длины
$64\cdot 64 = 4\,096$ в пространстве $\R^{4096}$. Двух разных людей
снимали много раз — получаем облако точек в этом пространстве.

Главное наблюдение: \emph{это облако очень узкое}. Хотя пространство
$4\,096$-мерное, реальные лица занимают в нём подпространство
размерности всего лишь $\sim 50$. Найти это подпространство — и
есть задача.

\begin{figure}[ht]
\centering
\includegraphics[width=0.85\linewidth]{eigenfaces_dataset.pdf}
\caption{Фрагмент датасета \engterm{Olivetti Faces}: $400$ фотографий
$64\times 64$ пикселя, $40$ человек по $10$ снимков каждого. Каждая
фотография --- точка в $\R^{4096}$.}
\label{fig:olivetti}
\end{figure}

\subsection{Шаг 1. Среднее лицо и центрирование}

Пусть у нас $N$ фотографий $x_1,\ldots,x_N\in\R^{d}$ (где $d=4096$ для
$64\times 64$). Вычислим \term{среднее лицо}:
\[
\bar x \;=\; \dfrac{1}{N}\sum_{i=1}^{N} x_i.
\]
И вычтем его из каждого изображения: $\tilde x_i = x_i - \bar x$.
\emph{Среднее лицо} (рис.~\ref{fig:mean-face}) — это размытая
маска, в которой видно «лицо вообще»: усреднённые глаза, нос,
рот, овал. Дальше мы будем работать только с отклонениями
от этого среднего.

\begin{figure}[ht]
\centering
\includegraphics[width=0.55\linewidth]{eigenfaces_mean.pdf}
\caption{Среднее лицо по $400$ снимкам датасета Olivetti.}
\label{fig:mean-face}
\end{figure}

\subsection{Шаг 2. Главные направления изменчивости}

Соберём центрированные изображения в строки матрицы
$X\in\R^{N\times d}$ (одна строка = одно лицо). Применим к ней SVD:
\[
X\;=\;U\,\Sigma\,V^{\!\top}.
\]
Каждая \emph{строка} $v_i^{\!\top}$ матрицы $V^{\!\top}$ — это вектор
длины $d=4096$, который можно нарисовать как картинку $64\times 64$.
Эти картинки и называются \term{собственными лицами} — eigenfaces
(рис.~\ref{fig:eigenfaces}). Содержательно:
\begin{itemize}\itemsep1pt
  \item \emph{$v_1$} — направление наибольшей изменчивости.
        Чаще всего это «освещение слева/справа».
  \item \emph{$v_2$} — направление второй по величине изменчивости
        (часто «улыбка / не улыбка»).
  \item \emph{$v_3,\ldots$} — постепенно более тонкие признаки: очки,
        форма щёк, наклон головы, индивидуальные особенности.
\end{itemize}

\begin{figure}[ht]
\centering
\includegraphics[width=\linewidth]{eigenfaces_components.pdf}
\caption{Первые 12 eigenfaces датасета Olivetti, отрисованные как
картинки $64\times 64$. Эти странные «полу-лица» — координатные
оси в пространстве, где живут все настоящие лица.}
\label{fig:eigenfaces}
\end{figure}

\subsection{Шаг 3. Лицо как вектор из 50 чисел}

Зафиксируем $k$ (обычно $k=30\div 100$) и спроектируем каждое лицо
$x$ на эти $k$ собственных лиц:
\begin{equation}\label{eq:eigenface-proj}
\alpha_i \;=\; v_i^{\!\top}(x-\bar x),\qquad i=1,\ldots,k.
\end{equation}
Получили вектор коэффициентов $\alpha=(\alpha_1,\ldots,\alpha_k)\in\R^{k}$.
Это и есть «лицо человека» с точки зрения алгоритма — \emph{всего
$50$ чисел} вместо $4096$.

\subsection{Шаг 4. Распознавание и реконструкция}

\paragraph{Распознавание.} Чтобы узнать, кто на новой фотографии $y$,
вычислим её коэффициенты $\beta$ по~\eqref{eq:eigenface-proj} и
найдём в базе ближайшего соседа: $\hat j = \argmin_j \norm{\beta-\alpha^{(j)}}$.

\paragraph{Реконструкция.} Зная $\alpha$, можно восстановить лицо:
\[
x \;\approx\; \bar x + \sum_{i=1}^{k}\alpha_i\,v_i.
\]
Рис.~\ref{fig:eigenfaces-recon} показывает, как качество
реконструкции растёт с $k$.

\begin{figure}[ht]
\centering
\includegraphics[width=\linewidth]{eigenfaces_reconstruction.pdf}
\caption{Реконструкция одной из фотографий датасета через первые $k$
eigenfaces. При $k=10$ узнаётся «общий тип» лица; при $k=50$ — уже
конкретный человек; при $k=150$ — почти оригинал.}
\label{fig:eigenfaces-recon}
\end{figure}

\subsection{Реализация на Python}

\begin{lstlisting}[label={lst:eigenfaces}]
import numpy as np
import matplotlib.pyplot as plt
from sklearn.datasets import fetch_olivetti_faces

# 1. Load: 400 faces, 64x64 each
faces = fetch_olivetti_faces().images        # (400, 64, 64)
X = faces.reshape(400, -1)                   # (400, 4096)

# 2. Center: subtract the mean face
x_bar = X.mean(axis=0)
X_centered = X - x_bar

# 3. SVD => eigenfaces are rows of Vt
U, S, Vt = np.linalg.svd(X_centered, full_matrices=False)

# 4. Project all faces onto top-k eigenfaces
k = 50
alpha = X_centered @ Vt[:k].T                # (400, 50)

# 5. Recognise: nearest neighbour in alpha-space
def recognise(y, k=50):
    y_centered = y.flatten() - x_bar
    beta = Vt[:k] @ y_centered
    distances = np.linalg.norm(alpha - beta, axis=1)
    return distances.argmin()

# 6. Reconstruct a face from k coefficients
def reconstruct(idx, k=50):
    return x_bar + Vt[:k].T @ alpha[idx, :k]
\end{lstlisting}

\begin{interestbox}
Метод eigenfaces — родоначальник всех современных систем распознавания
лиц, включая ту, что разблокирует ваш смартфон. Современные системы
заменяют SVD на свёрточную нейросеть (см.~§\,3.3.7), но базовая
идея — \emph{сжать лицо в короткий вектор, в котором близкие лица
оказываются рядом} — осталась той же. Эти короткие векторы сегодня
называют \term{эмбеддингами}.
\end{interestbox}

% =====================================================================
\section{Eigencats: то же самое для котиков\texorpdfstring{$^\star$}{*}}
\label{sec:linalg-eigencats}

Кошки — отличный материал для проверки гипотезы, что метод eigenfaces
работает на любых однородных объектах, не только на лицах. Возьмём
датасет в $\sim 2000$ котиков\footnote{Например,
\engterm{Cats Faces Dataset} на Kaggle: фотографии $64\times 64$.
Готовые ноутбуки выложены в репозитории \href{https://github.com/MerkulovDaniil/nla360}{nla360}.}
и проведём тот же эксперимент.

\begin{figure}[ht]
\centering
\includegraphics[width=0.85\linewidth]{eigencats_components.pdf}
\caption{Первые 8 eigencats. Видны общие «коты»: положение ушей,
форма морды, направление взгляда. Каждый конкретный кот собирается
как линейная комбинация этих универсальных «прото-котов».}
\label{fig:eigencats}
\end{figure}

Качественная картина повторяется: первое eigencat ловит освещение,
второе — наклон головы, дальше идут вариации формы морды и ушей.
В сжатом виде каждый кот превращается в $\sim 50$ чисел.

\begin{notebox}[title={Что это нам говорит о ИИ}]
\textbf{Пространство «реальных» объектов всегда узкое.} Будь это
лица людей, морды котов, рукописные цифры (датасет MNIST) или
фотографии природы — реальные данные занимают в исходном
многомерном пространстве \emph{маленькое} подпространство.
Распознавание = найти это подпространство и работать в нём.

\textbf{Эта идея — фундамент современных автокодировщиков и
эмбеддингов} (см.~§\,3.4.3--3.4.6). SVD — её самая ранняя и самая
понятная реализация.
\end{notebox}

% =====================================================================
\section{PageRank: как Google нашёл порядок в Интернете}
\label{sec:linalg-pagerank}

\subsection{Задача и история}

К 1996~г. поисковики типа \engterm{AltaVista} и \engterm{Yahoo!}
ранжировали страницы по тому, сколько раз искомое слово встречается
в тексте. Это легко обманывалось: достаточно было набить страницу
популярными словами белым цветом по белому фону, и она попадала
в топ выдачи.

Лоуренс Пейдж и Сергей Брин, тогда аспиранты Стэнфорда, предложили
другой принцип ранжирования: \emph{страница важна, если на неё
ссылаются другие важные страницы}.

\subsection{Граф ссылок как матрица}

Обозначим $N$ — число всех страниц Интернета (на 1998~г. — около
$2{,}4\cdot 10^{8}$, сегодня — $\sim 10^{11}$). Построим матрицу
$P\in\R^{N\times N}$:
\[
P_{ij}\;=\;\begin{cases}
\dfrac{1}{c_j}, & \text{если со страницы $j$ есть ссылка на $i$,}\\
0, & \text{иначе,}
\end{cases}
\]
где $c_j$ — \emph{полное число исходящих ссылок} со страницы $j$.
Содержательно: $P_{ij}$ — это вероятность того, что пользователь,
сидящий на странице $j$ и случайно ткнувший на одну из её ссылок,
окажется на странице $i$. Поэтому $P$ называется матрицей переходов
\term{случайного блуждания} по графу Интернета.

\subsection{Стационарное распределение = собственный вектор}

Пусть $r\in\R^{N}$ — вектор «важностей»: $r_i$ — степень важности
страницы $i$. Принцип Пейджа–Брина формализуется системой:
\begin{equation}\label{eq:pr-fixed}
r\;=\;P r.
\end{equation}
То есть важность каждой страницы равна сумме важностей тех страниц,
которые на неё ссылаются, делённой на число их исходящих ссылок.

Уравнение~\eqref{eq:pr-fixed} говорит: $r$ — это \emph{собственный
вектор} матрицы $P$ с собственным значением $1$. И в этом — вся суть
PageRank.

\begin{theorem}{Перрон--Фробениус для PageRank}{perron}
Если все элементы матрицы $P$ положительны (или хотя бы граф ссылок
сильно связен и непериодичен), то собственное значение $1$ у $P$
существует, оно \emph{простое} и наибольшее по модулю, а
соответствующий собственный вектор $r$ единственен и имеет
положительные компоненты. Этот $r$ называется \term{стационарным
распределением} случайного блуждания.
\end{theorem}

(Эта теорема — Оскар Перрон, 1907~г., и Георг Фробениус, 1912~г. —
исторически появилась задолго до Интернета, в контексте моделей
популяций.)

\subsection{Демпфирующий множитель и формула Брина--Пейджа}

В реальной матрице Интернета у некоторых страниц нет исходящих ссылок
(\engterm{dangling pages}), а в других группах — циклы. Для таких
страниц соответствующий столбец матрицы переходов заранее заменяют
равномерным распределением по всем страницам: иначе формула $1/c_j$
при $c_j=0$ не определена, а масса случайного блуждания «исчезала» бы
в dangling pages. Чтобы
гарантировать применимость теоремы Перрона--Фробениуса, Пейдж и Брин
ввели \term{демпфирующий множитель} $\beta\in(0,1)$ (обычно
$\beta=0{,}85$):
\begin{equation}\label{eq:pagerank}
\boxed{\;r\;=\;\beta P r + \dfrac{1-\beta}{N}\,\mathbf{1}.\;}
\end{equation}
Здесь $\mathbf{1}=(1,1,\ldots,1)^{\!\top}$. Содержательно:
с вероятностью $\beta\approx 0{,}85$ случайный сёрфер кликает по
ссылке, а с вероятностью $1-\beta\approx 0{,}15$ — телепортируется
на случайную страницу Интернета.

\subsection{Как находить $r$: степенной метод}

Прямо решать систему $(I-\beta P)r=\frac{1-\beta}{N}\mathbf{1}$
с $N\sim 10^{11}$ невозможно. Но и не надо: вектор $r$ можно
получить простой итерацией.

\begin{algorithmbox}
\textbf{Степенной метод для PageRank.}
\begin{enumerate}\itemsep1pt
  \item Инициализация: $r^{(0)} \leftarrow \tfrac{1}{N}\mathbf{1}$.
  \item Итерация: $r^{(t+1)} \leftarrow \beta\,P\,r^{(t)} +
        \dfrac{1-\beta}{N}\,\mathbf{1}$.
  \item Остановка: $\norm{r^{(t+1)}-r^{(t)}} < \varepsilon$.
\end{enumerate}
\end{algorithmbox}

Сходимость гарантируется теоремой Перрона--Фробениуса:
$\norm{r^{(t)}-r}\leq C\cdot\beta^{t}$, то есть погрешность убывает
геометрически со знаменателем $\beta=0{,}85$. На практике хватает
$50\div 100$ итераций.

\subsection{Игрушечный пример: 6 страниц}

Возьмём маленький «Интернет» из шести страниц (рис.~\ref{fig:pagerank-graph})
и посчитаем для него PageRank.

\begin{figure}[ht]
\centering
\begin{tikzpicture}[scale=1.0, every node/.style={circle, draw, accent,
              fill=accent!10, minimum size=8mm, font=\small}]
  \node (A) at (0,2)    {A};
  \node (B) at (2.5,3)  {B};
  \node (C) at (5,2)    {C};
  \node (D) at (5,-0.5) {D};
  \node (E) at (2.5,-1.5){E};
  \node (F) at (0,-0.5) {F};
  \begin{scope}[->, thick, accent2]
    \draw (A) -- (B);
    \draw (A) -- (C);
    \draw (B) -- (C);
    \draw (C) -- (A);
    \draw (D) -- (C);
    \draw (D) -- (A);
    \draw (E) -- (D);
    \draw (E) -- (B);
    \draw (E) -- (F);
    \draw (F) -- (E);
    \draw (F) -- (B);
  \end{scope}
\end{tikzpicture}
\caption{Игрушечный «Интернет» из 6 страниц со ссылками. Видно,
например, что на страницу C ссылаются три другие, а на F — только
одна; будет ли C лидером по PageRank?}
\label{fig:pagerank-graph}
\end{figure}

\begin{lstlisting}[label={lst:pagerank}]
import numpy as np

# Adjacency list -> column-stochastic transition matrix P
links = {'A': ['B','C'], 'B': ['C'], 'C': ['A'],
         'D': ['C','A'], 'E': ['D','B','F'], 'F': ['E','B']}
pages = list(links); n = len(pages); idx = {p: i for i, p in enumerate(pages)}

P = np.zeros((n, n))
for j, p in enumerate(pages):
    out = links[p]
    for q in out:
        P[idx[q], j] = 1.0 / len(out)

# PageRank via power iteration
beta, eps = 0.85, 1e-10
r = np.ones(n) / n
for t in range(200):
    r_new = beta * P @ r + (1 - beta) / n
    if np.linalg.norm(r_new - r, 1) < eps: break
    r = r_new

for p, ri in sorted(zip(pages, r), key=lambda t: -t[1]):
    print(f'{p}: {ri:.4f}')
\end{lstlisting}

\noindent Выдача программы:
\[
\text{C: }0{,}295,\quad
\text{A: }0{,}252,\quad
\text{B: }0{,}175,\quad
\text{E: }0{,}098,\quad
\text{F: }0{,}092,\quad
\text{D: }0{,}087.
\]
Лидер — \texttt{C}, хотя на неё формально ссылаются столько же
страниц, сколько на \texttt{A}: PageRank учёл, что одна из ссылок
на \texttt{C} приходит с уже-важной страницы \texttt{B}.

\begin{historybox}
Алгоритм PageRank был опубликован в 1998~г. в студенческой работе
\emph{Sergey Brin, Larry Page}, «\textit{The PageRank Citation Ranking:
Bringing Order to the Web}». В том же году Брин и Пейдж основали
Google. К 2004~г. компания вышла на IPO, а к 2010-м стала одной из
крупнейших в мире. Сам Стэнфорд получил $\sim 336$ миллионов долларов
от лицензирования патента --- который изначально просто описывал
итерационный метод для собственного вектора.

Идея использовать собственный вектор графа для измерения «важности»
позднее была обобщена и применена в наукометрии (\engterm{Eigenfactor}
для ранжирования научных журналов), в рекомендательных системах,
а также в современных нейросетях для графов (\engterm{Graph Neural
Networks}).
\end{historybox}

% =====================================================================
\section{Когда нейросеть устойчива: оценки Липшицевости\texorpdfstring{$^\star$}{*}}
\label{sec:linalg-lipschitz}

\subsection{Сюжет: когда панда становится гиббоном}

В 2014~г. Иэн Гудфеллоу и соавторы опубликовали ставшую знаменитой
картинку (рис.~\ref{fig:adv-attack}): обычная фотография панды,
которую современная нейросеть уверенно опознаёт как «panda» с
вероятностью $57{,}7\%$. К каждому пикселю прибавлена крошечная,
визуально незаметная человеку поправка. На обновлённой картинке
сеть теперь уверенно (с вероятностью $99{,}3\%$) считает, что
видит \emph{гиббона}.

\begin{figure}[ht]
\centering
\includegraphics[width=0.95\linewidth]{adversarial_panda.pdf}
\caption{Adversarial-атака на нейросеть (Goodfellow et al., 2014).
Слева — оригинал, классифицирован как «panda» с уверенностью $57{,}7\%$.
В центре — почти невидимое возмущение, масштабированное для наглядности.
Справа — атакованный пример, классифицируется как «gibbon» с
уверенностью $99{,}3\%$.}
\label{fig:adv-attack}
\end{figure}

Этот эффект называется \term{adversarial-атакой} и долгое время
считался экзотикой. Сейчас понятно: причина в простом числе,
которое сопровождает любую функцию — её \emph{константе Липшица}.

\subsection{Константа Липшица: формальное определение}

\begin{definition}{Константа Липшица функции}{lipschitz}
Функция $f\colon\R^{n}\to\R^{m}$ называется \term{Липшицевой} с
константой $L>0$, если для любых $x,y\in\R^{n}$ выполнено
\begin{equation}\label{eq:lipschitz}
\norm{f(x)-f(y)}\;\leq\;L\,\norm{x-y}.
\end{equation}
Наименьшее $L$, для которого~\eqref{eq:lipschitz} выполнено,
называется \emph{константой Липшица} функции $f$ и обозначается
$\mathrm{Lip}(f)$.
\end{definition}

Содержательно $\mathrm{Lip}(f)$ показывает, во сколько раз может
\emph{усилиться} входное возмущение на выходе функции. Если
$\mathrm{Lip}(f)=L$, то малое возмущение $\delta$ на входе
может вызвать возмущение до $L\cdot\norm{\delta}$ на выходе.

\paragraph{Пример.} Для линейного отображения $f(x)=Wx$ имеем
$f(x)-f(y)=W(x-y)$, откуда $\norm{f(x)-f(y)}\leq \norm{W}_2\cdot\norm{x-y}$,
и константа Липшица — это в точности спектральная норма матрицы $W$:
\begin{equation}\label{eq:lip-linear}
\mathrm{Lip}(Wx)\;=\;\norm{W}_2\;=\;\sigma_1(W).
\end{equation}
Здесь $\sigma_1(W)$ — наибольшее сингулярное число матрицы $W$.
Вот \emph{здесь} в нашу историю входит SVD из §\,\ref{sec:linalg-svd}.

\subsection{Композиция: Липшицева константа цепочки}

\begin{theorem}{Композиция Липшицевых функций}{lip-compose}
Если $f$ и $g$ — Липшицевы с константами $L_f$ и $L_g$, то
их композиция $f\circ g$ — тоже Липшицева, причём
\begin{equation}\label{eq:lip-prod}
\mathrm{Lip}(f\circ g)\;\leq\;\mathrm{Lip}(f)\cdot\mathrm{Lip}(g).
\end{equation}
\end{theorem}

\begin{proof}
$\norm{f(g(x))-f(g(y))}\leq L_f\norm{g(x)-g(y)}\leq L_f L_g\norm{x-y}$. $\square$
\end{proof}

\subsection{Применение к нейросети}

Возьмём типичную нейросеть глубины $K$:
\[
f(x)\;=\;\sigma\bigl(W_K\,\sigma(W_{K-1}\,\sigma(\cdots\sigma(W_1 x))\bigr),
\]
где $W_i$ — матрицы весов, $\sigma$ — поэлементная активация (ReLU,
сигмоида, и т.п.). Для большинства активаций $\mathrm{Lip}(\sigma)\leq 1$
(для ReLU и $\tanh$ она равна $1$, а для логистической сигмоиды — $1/4$). По
теореме~\ref{thm:lip-compose}:
\begin{equation}\label{eq:nn-lip}
\boxed{\;\mathrm{Lip}(f)\;\leq\;\prod_{i=1}^{K}\norm{W_i}_2\;}
\end{equation}

Это и есть тот «лишний рычаг», который объясняет уязвимость
нейросетей. Если $\norm{W_i}_2 \approx 2$ для каждого слоя из $K=20$
слоёв, то возмущение на входе \emph{потенциально} усиливается
в $2^{20}\approx 10^{6}$ раз. Невидимая поправка $\delta$ с
$\norm{\delta}=10^{-3}$ может вызвать изменение выхода до $10^{3}$ —
этого хватит, чтобы перебросить классификатор из одного класса в другой.

\subsection{Как защищаться: спектральная нормализация}

Естественная стратегия: после каждого шага обучения нормализовать
матрицу $W_i$ так, чтобы $\norm{W_i}_2 = 1$. Это и есть метод
\term{Spectral Normalization} (Миято и соавт., 2018), сделавший
обучение генеративных нейросетей (GAN) намного стабильнее.

Технически: на каждом шаге считаем $\sigma_1(W)$ (с помощью одной
итерации степенного метода — почти бесплатно!) и заменяем
$W \mapsto W/\sigma_1(W)$.

\begin{lstlisting}[label={lst:spec-norm}]
def spectral_normalize(W, n_iter=1):
    """One step of power iteration to estimate sigma_1(W) and rescale."""
    u = np.random.randn(W.shape[0])
    for _ in range(n_iter):
        v = W.T @ u; v /= np.linalg.norm(v)
        u = W @ v;   u /= np.linalg.norm(u)
    sigma_1 = u @ W @ v
    return W / sigma_1
\end{lstlisting}

\begin{interestbox}
\textbf{Связь с устойчивостью в физике.} Похожая ситуация хорошо
известна в теории динамических систем: если константа Липшица
отображения $> 1$, малое возмущение начального условия экспоненциально
нарастает (классический эффект бабочки в погодных моделях).
Spectral Normalization можно рассматривать как искусственное
гашение этого эффекта.
\end{interestbox}

% =====================================================================
\section*{Подведение итогов главы}
\addcontentsline{toc}{section}{Подведение итогов главы}

\begin{itemize}\itemsep1pt
  \item Матрица — это \emph{универсальный носитель данных} для ИИ:
        картинки, тексты, графы, веса нейросети.
  \item \textbf{Числа с плавающей точкой} (§\ref{sec:linalg-floats}) —
        вычисления на компьютере всегда приближённые с относительной
        ошибкой $\sim\varepsilon_{\text{маш}}\approx 10^{-16}$.
        Самая опасная ошибка — \emph{катастрофическое сокращение} при
        вычитании близких чисел; её часто можно устранить, переписав
        формулу алгебраически эквивалентным способом.
  \item \textbf{Сингулярное разложение} $A=U\Sigma V^{\!\top}$
        (§\ref{sec:linalg-svd}) представляет любую матрицу как
        «поворот → растяжение по осям → поворот». Сингулярные числа
        $\sigma_i$ — это коэффициенты растяжения.
  \item \textbf{Теорема Эккарта--Янга} (§\ref{sec:linalg-eckart}):
        $A_k=\sum_{i\leq k}\sigma_i u_i v_i^{\!\top}$ — лучшее
        приближение ранга $k$ матрицы $A$. На этой теореме
        стоит сжатие изображений, метод eigenfaces и современный
        LoRA для дообучения больших языковых моделей.
  \item \textbf{Eigenfaces / eigencats}
        (§\ref{sec:linalg-eigenfaces}--\ref{sec:linalg-eigencats}):
        реальные данные (лица, котики, цифры) живут в очень узком
        подпространстве своего номинального многомерного пространства.
        Найти это подпространство = применить SVD к выборке.
  \item \textbf{PageRank} (§\ref{sec:linalg-pagerank}): «важность»
        страницы Интернета = компонента собственного вектора матрицы
        переходов случайного блуждания. Алгоритм находит его простой
        итерацией; теорема Перрона--Фробениуса гарантирует сходимость.
  \item \textbf{Константа Липшица} нейросети
        (§\ref{sec:linalg-lipschitz}) растёт мультипликативно от
        слоя к слою; именно поэтому глубокие сети уязвимы к
        adversarial-атакам. Спектральная нормализация ограничивает
        $\norm{W_i}_2=1$ и стабилизирует обучение.
\end{itemize}

% =====================================================================
\section*{Что почитать дальше}
\addcontentsline{toc}{section}{Что почитать дальше}

\begin{itemize}\itemsep1pt
  \item \emph{Е.\,Е.~Тыртышников.} Методы численного анализа. М.: Академия, 2007. ---
        классический российский учебник по вычислительной линейной
        алгебре университетского уровня.
  \item \emph{G.~Strang.} Linear Algebra and Learning from Data. Wellesley-Cambridge Press, 2019. ---
        свежая книга от автора знаменитого MIT-курса 18.06, написанная
        с прицелом на data science.
  \item \emph{L.\,N.~Trefethen, D.~Bau.} Numerical Linear Algebra. SIAM, 1997. ---
        классика, с очень понятным изложением SVD и итерационных методов.
  \item \emph{Видеокурс «Numerical Linear Algebra»}, И.\,В.~Оселедец, Сколтех
        (\href{https://github.com/MerkulovDaniil/nla360}{github.com/MerkulovDaniil/nla360}).
        Лекции, семинары, ноутбуки --- из которых, в частности,
        и взяты eigencats для этой главы.
  \item \emph{S.~Brin, L.~Page.} The Anatomy of a Large-Scale
        Hypertextual Web Search Engine. WWW Conference, 1998. ---
        оригинальная статья про PageRank, читается легко.
  \item \emph{I.~Goodfellow et al.} Explaining and Harnessing
        Adversarial Examples. ICLR, 2015. ---
        статья про панду-гиббона, из которой взят рис.~\ref{fig:adv-attack}.
\end{itemize}

% =====================================================================
\section*{Большой итоговый проект}
\addcontentsline{toc}{section}{Большой итоговый проект}

\textbf{Постройте свою рекомендательную систему для фильмов.}

Возьмите датасет MovieLens-100K (открытый, $100\,000$ оценок,
$\sim 1000$ пользователей $\times$ $\sim 1700$ фильмов). Постройте
матрицу «пользователь--фильм» $R$, где $R_{ij}\in\{1,\ldots,5\}$ —
оценка, поставленная пользователем $i$ фильму $j$, а $0$ — «не смотрел».

\textbf{Важно:} нули в матрице $R$ означают \emph{пропущенные} оценки,
а не настоящие нулевые рейтинги. Прямое SVD по матрице с нулями — лишь
грубый учебный baseline: непросмотренные фильмы тогда воспринимаются как
фильмы с очень низкой оценкой. Более корректный вариант — обучать
малоранговую модель только на известных оценках либо предварительно
заполнить пропуски средними значениями пользователя/фильма и затем
сравнить результат.

\begin{enumerate}\itemsep1pt
  \item Постройте малоранговое приближение $R\approx U\Sigma V^{\!\top}$
        с $k=20$ компонентами через SVD.
  \item Объясните содержательно, что такое строки $U$ (профили
        пользователей) и столбцы $V$ (профили фильмов) в этом
        приближении. Сколько чисел теперь хранит модель вместо
        $1000\cdot 1700$?
  \item Для нескольких реальных пользователей выпишите топ-5 фильмов,
        которые рекомендация модели предлагает посмотреть. Совпадают
        ли они с любимыми жанрами пользователя?
  \item Найдите два «противоположных» фильма с наибольшим расстоянием
        между их столбцами $V$ и подумайте, можно ли увидеть содержательный
        смысл в этой разделяющей оси (боевик vs. мелодрама? старое vs. новое?).
  \item ($\star$) Замените SVD на алгоритм PageRank, построив граф
        «фильмы, которые часто оценивают одинаково». Получится ли тот
        же топ?
\end{enumerate}

% =====================================================================
\begin{tasksbox}
\begin{enumerate}\itemsep1pt
  \item Сложите в Python $10\,000\,000$ копий числа $0{,}1$. Сколько
        должно получиться? Сколько получается на самом деле? Почему?
  \item Вычислите $f(x)=\sqrt{x+1}-\sqrt{x}$ при $x=10^{16}$
        напрямую и через формулу $f(x)=\frac{1}{\sqrt{x+1}+\sqrt{x}}$.
        Сравните ответы. Какой устойчив, почему?
  \item Найдите вручную SVD матрицы
        $A=\bigl(\begin{smallmatrix}3 & 0\\0 & 4\end{smallmatrix}\bigr)$
        и $A=\bigl(\begin{smallmatrix}0 & 1\\-1 & 0\end{smallmatrix}\bigr)$.
        Что в обоих случаях геометрически делают $U, \Sigma, V$?
  \item Возьмите свою фотографию $300\times 300$ в градациях серого.
        Сделайте сжатие через SVD с рангами $k=5, 20, 50, 100$. При
        каком $k$ вы перестаёте узнавать самого себя?
  \item Для игрушечного графа из рис.~\ref{fig:pagerank-graph}
        прокрутите 5 итераций степенного метода \emph{вручную},
        начав с $r^{(0)}=(\tfrac16,\ldots,\tfrac16)^{\!\top}$. Совпали ли
        ваши значения с программными после $5$ шагов? а после $20$?
  \item $^{\star}$ Покажите, что для функции $\sigma(x)=\max(0,x)$
        (ReLU) константа Липшица равна $1$. Тот же вопрос для
        сигмоиды $\sigma(x)=1/(1+e^{-x})$: чему равна $\mathrm{Lip}(\sigma)$
        и в какой точке производная максимальна?
  \item $^{\star}$ Найдите спектральную норму матрицы
        $W=\bigl(\begin{smallmatrix}1 & 2\\2 & 1\end{smallmatrix}\bigr)$
        двумя способами: (а) через определение
        $\norm{W}_2=\max_{\norm{x}=1}\norm{Wx}$ (с использованием
        собственных значений $W^{\!\top}W$); (б) одной итерацией
        степенного метода (с любым стартом). Совпали ли ответы?
\end{enumerate}
\end{tasksbox}
